% ============================================================================
% Fichier     : SAFI_main.tex
% Titre       : Supervision, Analyse Forensique et Gestion des Incidents
% Auteur      : MINKA MI NGUIDJOÏ Thierry Emmanuel
% Institution : Chaire UNESCO d'Afrique Centrale Accès TIC : ENSP / LIMSI
% Version     : 0.1 -- octobre 2026 (squelette)
% ============================================================================

\documentclass[12pt, twoside, openright, a4paper]{book}

% --- Feuille de style unifiée ---
\usepackage{Style_SAFI}

% --- Bibliographie ---
\addbibresource{bibliographie/SAFI_references.bib}

% ============================================================================
% MÉTADONNÉES PDF
% ============================================================================
\hypersetup{
    pdftitle={Supervision, Analyse Forensique et Gestion des Incidents},
    pdfauthor={MINKA MI NGUIDJOÏ Thierry Emmanuel},
    pdfsubject={SIS589 · Master Professionnel en Sécurité des Systèmes Informatiques et des Données · Université de Garoua},
    pdfkeywords={SOC, SIEM, IDS, forensique, ISO 27035, ISO 27037, réponse à incident, CSIRT, SAFI}
}

\begin{document}

% --------------------------------------------------------------------------
% ÉLÉMENTS PRÉLIMINAIRES (pages romaines)
% --------------------------------------------------------------------------
\frontmatter

    % ==============================================================================
% Fichier : front/00_page_titre.tex
% Titre   : Page de titre - \titreCours
% Auteur  : MINKA MI NGUIDJOÏ Thierry Emmanuel
% Version : 1.1 -- novembre 2026 (charte des cours, modèle de HE2)
% ==============================================================================
% Logos : Université de Garoua à gauche et à droite.
% ==============================================================================

\begin{titlepage}
\thispagestyle{empty}

\begin{tikzpicture}[remember picture, overlay]
    \fill[CouleurPrimaire]
        (current page.north west) rectangle
        ([yshift=-3.6cm] current page.north east);
    \fill[CouleurAccent]
        ([yshift=-3.6cm] current page.north west) rectangle
        ([yshift=-4.0cm] current page.north east);
    \fill[CouleurPrimaire]
        (current page.south west) rectangle
        ([yshift=2.5cm] current page.south east);
    \fill[CouleurAccent]
        ([yshift=2.5cm] current page.south west) rectangle
        ([yshift=2.9cm] current page.south east);
    % En-tête institutionnel, posé dans le bandeau supérieur
    \node[align=center, text=white] at ([yshift=-1.8cm] current page.north) {%
        {\fontsize{11}{13}\selectfont
        \textsc{République du Cameroun}}\\[0.3ex]
        {\color{CouleurAccent!60!white}\rule{0.4\paperwidth}{0.4pt}}\\[0.6ex]
        {\fontsize{10}{12}\selectfont\textbf{\tutelle}}\\[0.2ex]
        {\fontsize{9}{11}\selectfont \faculte{} · \departement}\\[0.2ex]
        {\fontsize{8}{10}\selectfont \unitepedago}};
    % Logos : l'Université de Garoua à gauche et à droite (un seul établissement)
    \node[fill=white, rounded corners=4pt, inner sep=2pt]
        at ([xshift=1.75cm, yshift=-1.8cm] current page.north west)
        {\includegraphics[height=2.6cm]{figures/logos/garoua_logo.png}};
    \node[fill=white, rounded corners=4pt, inner sep=2pt]
        at ([xshift=-1.75cm, yshift=-1.8cm] current page.north east)
        {\includegraphics[height=2.6cm]{figures/logos/garoua_logo.png}};
    % Année et version, posées dans le bandeau inférieur
    \node[align=center] at ([yshift=1.25cm] current page.south) {%
        {\color{white}\bfseries\large Année académique \anneescolaire}\\[0.3ex]
        {\color{CouleurAccent!70!white}\small
        Version 1.0 | Document pédagogique de l'UE \codeUE}};
\end{tikzpicture}

\vspace*{0.6cm}

% --- Titre principal (sur deux lignes) ---
\begin{center}
    {\color{CouleurPrimaire}\rule{0.85\textwidth}{1.2pt}}\\[1.2ex]
    {\fontsize{24}{30}\selectfont\bfseries\color{CouleurPrimaire}
        Supervision, Analyse Forensique\\[0.4ex]
        et Gestion des Incidents}\\[0.8ex]
    {\color{CouleurAccent}\rule{0.85\textwidth}{2pt}}\\[0.5ex]
    {\color{CouleurPrimaire}\rule{0.85\textwidth}{0.5pt}}
\end{center}

\vspace{0.8cm}

\begin{center}
    {\fontsize{13}{16}\selectfont\itshape\color{CouleurSecondaire}
        Détecter, investiguer, maîtriser~:\\
        du SOC à la preuve opposable}
\end{center}

\vspace{0.5cm}

% --- Cartouche sigle et code UE ---
\begin{center}
    \begin{tikzpicture}
        \node[draw=CouleurPrimaire, fill=CouleurPrimaire, rounded corners=6pt,
              inner xsep=14pt, inner ysep=6pt]{%
            \color{white}\bfseries\large \sigleCours{} - UE~\codeUE
        };
    \end{tikzpicture}
\end{center}

\vspace{1.0cm}

% --- Schéma de la discipline : le cycle de gestion d'un incident ---
\begin{center}
\begin{tikzpicture}[
    f/.style = {draw=CouleurBordDef, fill=CouleurDef, rounded corners=3pt,
        minimum width=2.2cm, minimum height=0.8cm, align=center,
        font=\small\bfseries\color{CouleurPrimaire}},
    arrow/.style = {-{Stealth[length=4pt]}, thick, color=CouleurSecondaire}]
\foreach \t [count=\i] in {Surveiller, Détecter, Qualifier, Investiguer, Restaurer}
    \node[f] (n\i) at (2.8*\i - 2.8, 0) {\t};
\foreach \i [evaluate=\i as \j using int(\i+1)] in {1,2,3,4}
    \draw[arrow] (n\i) -- (n\j);
\draw[arrow] (n5.south) to[bend left=18] node[below, font=\footnotesize\itshape,
    text=CouleurSecondaire] {retour d'expérience} (n1.south);
\end{tikzpicture}
\end{center}

\vspace{0.6cm}

\begin{center}
    {\large\bfseries\color{CouleurPrimaire}
        \auteurcours}\\[0.5ex]
    {\normalsize\color{CouleurSecondaire}
       Enseignant Chercheur | \tutelle}\\[0.2ex]
    {\small\color{CouleurBordInfo}
        \faculte{} · \departement}\\[0.2ex]
    {\small\color{CouleurBordInfo}
        \href{mailto:maletyon@proton.me}{maletyon@proton.me}}
\end{center}

\end{titlepage}

    \cleardoublepage

    % ==============================================================================
% Fichier : front/01_avant_propos.tex : SAFI (SIS589)
% ==============================================================================

\chapter*{Avant-propos}
\addcontentsline{toc}{chapter}{Avant-propos}
\markboth{Avant-propos}{Avant-propos}

Ce livre est le support de l'unité d'enseignement \textbf{SIS589, Supervision, Analyse Forensique et Gestion des
Incidents}, du Master
Professionnel en Sécurité des Systèmes Informatiques et des Données de
l'Université de Garoua (Faculté des Sciences, Département de Mathématiques et
Informatiques).

\medskip
Il prolonge, sur le terrain opérationnel, les enseignements de sécurité et
d'audit~: là où l'on apprend ailleurs à apprécier les risques et à simuler les
attaques, SIS589 apprend à \textbf{détecter, investiguer et maîtriser} ce qui
survient en temps réel. Le fil conducteur est unique~: une entreprise
fictive, \textbf{LiZbiZ}, transporteur de colis présent à Garoua, subit un
rançongiciel. Le cours suit l'incident sur trois temps (la supervision qui
aurait dû le voir venir, l'analyse forensique qui l'établit, la réponse qui
le maîtrise), à travers trois opérations situées dans le Nord~: \textbf{BÉNOUÉ},
\textbf{FARO} et \textbf{MAYO}.

\medskip
Chaque notion est ancrée à la fois dans les standards internationaux (ISO/IEC
27035 et 27037, NIST SP 800-86, RFC 3227, MITRE ATT\&CK) et dans le droit
camerounais et africain (lois n°~2010/012, 2010/013 et 2024/017, Convention
de Malabo, Convention de Budapest). Un fil théorique parcourt le livre~: le
\textbf{Trilemme CRO} (Confidentialité, Fiabilité, Opposabilité), par lequel
se juge toute trace et tout acte d'enquête.

\medskip
Les cas sont fictifs~: les noms, les dates et les montants le sont~; les
mécanismes techniques et les références normatives sont exacts. Le livre
s'accompagne de l'environnement \textbf{NDAP SUKULU}, qui fait pratiquer
chaque chapitre.

\bigskip
\noindent \IfFileExists{../../Logos/signature_avant_propos.png}{\includegraphics[height=1.6cm]{../../Logos/signature_avant_propos.png}}{\textit{MINKA MI NGUIDJOÏ Thierry Emmanuel}}

    \cleardoublepage

    % ==============================================================================
% Fichier : front/02_mode_emploi.tex : SAFI (SIS589)
% ==============================================================================

\chapter*{Mode d'emploi}
\addcontentsline{toc}{chapter}{Mode d'emploi}
\markboth{Mode d'emploi}{Mode d'emploi}

\section*{À qui s'adresse ce cours}

Aux étudiants du Master Professionnel en Sécurité des Systèmes Informatiques
et des Données, et aux professionnels qui exercent ou se destinent aux
métiers d'\textbf{analyste SOC}, d'\textbf{analyste forensique}, de
\textbf{responsable de la réponse à incident} et de \textbf{RSSI}. Des
notions de réseaux (TCP/IP), d'administration Windows et Linux et de droit
numérique sont utiles.

\section*{Comment lire ce livre}

Le livre comprend un prologue et six parties. Chaque chapitre suit le même
plan~: une épigraphe, le niveau et les prérequis, les objectifs, une question
directrice, cinq à neuf sections, un encadré « à retenir », dix exercices
gradués et une transition. Les niveaux des exercices sont signalés par des
étoiles~: \niveaubase{} fondamental, \niveaumoyen{} intermédiaire,
\niveauexpert{} expert.

\section*{Les encadrés}

\begin{itemize}
    \item \textbf{Scénario opérationnel}~: un épisode des opérations BÉNOUÉ,
          FARO ou MAYO, avec \emph{constatation} (faits) et \emph{analyse}
          (interprétation, niveau de confiance).
    \item \textbf{Artefact}~: la fiche d'une pièce d'enquête (nature, source,
          pertinence, fiabilité, opposabilité).
    \item \textbf{Chaîne de custody}~: une procédure de préservation de la
          preuve.
    \item \textbf{Analyse CRO}~: l'arbitrage Confidentialité, Fiabilité,
          Opposabilité.
    \item \textbf{Application métier}~: ce que le chapitre change pour un
          métier.
    \item \textbf{Cadre juridique}~: le texte de droit applicable.
\end{itemize}

\section*{Constatation et analyse}

Le cours distingue systématiquement ce que l'on \textbf{constate} (un fait
vérifiable) de ce que l'on en \textbf{déduit} (une interprétation, assortie
d'un niveau de confiance). C'est la discipline première de la supervision et
du forensique.

\section*{NDAP SUKULU et évaluation}

L'environnement compagnon NDAP SUKULU (page suivante) propose, pour chaque
chapitre, une fiche d'auto-évaluation et, pour les chapitres pratiques, un
lab sur pièces. Fiches et labs se font hors des 30~heures de cours. Le code
du cours est \textbf{\codeUE}.

    \cleardoublepage

    % ==============================================================================
% Fichier : front/03_ndap_sukulu.tex : SAFI (SIS589)
% ==============================================================================

\chapter*{NDAP SUKULU, l'environnement compagnon}
\addcontentsline{toc}{chapter}{NDAP SUKULU, l'environnement compagnon}
\markboth{NDAP SUKULU}{NDAP SUKULU}

\begin{center}
    \includegraphics[width=0.52\textwidth]{figures/logos/logo_ndap_sukulu.png}
\end{center}

\medskip
\textit{Ndap sukulu} signifie « la salle de classe » en bassa, et
\textit{Malet Yon}, « ton enseignant ». \textbf{NDAP SUKULU} est
l'environnement compagnon du cours~: pour chaque chapitre, une fiche
d'auto-évaluation de 40~questions, corrigée automatiquement~; pour les
chapitres pratiques, un lab sur pièces (journaux, images, flux exportés)
dans un cadre entièrement fermé~; et un tableau de bord de progression.

\medskip
On y accède par le code de cours \textbf{\codeUE}. Les fiches et les labs se
font hors des 30~heures de cours. Les pièces des labs sont fictives et
établies pour l'enseignement.

% ==============================================================================
% Fichier : front/03b_labs.tex : SAFI (SIS589)
% ==============================================================================

\section*{Les dix labs de NDAP SUKULU}

\begin{tableaucours}{L{2.4cm}L{8.4cm}L{1.1cm}}{%
    \tableauHeader{Lab} &
    \tableauHeader{Sujet} &
    \tableauHeader{Ch.}
}
SAFI-LAB-01 & Lire un tableau de bord de supervision & 3 \\
SAFI-LAB-02 & Analyser des journaux et repérer les angles morts & 4 \\
SAFI-LAB-03 & Écrire et éprouver une règle de détection (SIEM) & 7 \\
SAFI-LAB-04 & Reconnaître une intrusion dans des alertes corrélées & 5 \\
SAFI-LAB-05 & Chaîne de custody et qualification d'artefacts & 11 \\
SAFI-LAB-06 & Forensique système~: reconstituer une chronologie & 12 \\
SAFI-LAB-07 & Forensique réseau~: exfiltration et pilotage (flux exportés) & 14 \\
SAFI-LAB-08 & Classer et traiter des incidents (gravité, délais) & 15 \\
SAFI-LAB-09 & Mapping d'un incident à MITRE ATT\&CK & 17 \\
SAFI-LAB-10 & Couverture de détection et verdict OUI/NON/NEK & 20 \\
\end{tableaucours}

\medskip
Chaque lab dure environ une heure et se corrige automatiquement. Les pièces
se trouvent dans le dossier \texttt{preuves} (lecture seule) et une copie de
travail dans \texttt{travail}.


    \cleardoublepage

    \tableofcontents
    \cleardoublepage

    \listoffigures
    \listoftables
    \cleardoublepage

% --------------------------------------------------------------------------
% CORPS DU COURS (pages arabes)
% Numérotation des chapitres CONTINUE sur tout le cours : le prologue est le
% chapitre 0, puis C01 à C21, un numéro par fichier.
% --------------------------------------------------------------------------
\mainmatter

    \part*{Prologue Opérationnel}
    \addcontentsline{toc}{part}{Prologue Opérationnel}
        \setcounter{chapter}{-1}
        \IfFileExists{partie0/P00_decouverte_rancongiciel.tex}{% ============================================================================
% Fichier  : partie0/P00_decouverte_rancongiciel.tex
% Titre    : Prologue : la Découverte
% Auteur   : MINKA MI NGUIDJOÏ Thierry Emmanuel
% ============================================================================

\chapter{Prologue : la Découverte}
\label{chap:prologue}
\chaptermark{Prologue}

\epigraph{%
    « Quand l'écran de rançon s'affiche, l'attaque est déjà vieille~;\\
    la seule question est de savoir si quelqu'un l'a vue venir. »%
}{Principe du cours}

\niveauChapitre{Introductif}{Aucun prérequis. Les notions sont reprises en
détail dans les chapitres.}

\begin{boxobjectifs}
\begin{itemize}
    \item Situer la supervision, le forensique et la réponse à incident dans
          une même journée réelle.
    \item Distinguer ce qui relève de la détection, de l'investigation et de
          la décision.
    \item Découvrir le fil rouge LiZbiZ et le Trilemme CRO.
    \item Formuler les questions auxquelles le cours répond.
\end{itemize}
\end{boxobjectifs}

\bigskip

% ============================================================================
\section*{6 mars 2026, 6 h 40, site de LiZbiZ à Garoua}
% ============================================================================

Le vigile de nuit a cru à une coupure~: les écrans de l'entrepôt affichent
tous le même fond noir et un texte en anglais réclamant un paiement. À
6~h~40, l'astreinte informatique est jointe. À 7~h~15, M.~KOUNOU, responsable
informatique, constate que le serveur de fichiers \texttt{FS-01}, la base de
suivi des colis et deux postes de la comptabilité sont chiffrés. Les
sauvegardes de la semaine sont, elles aussi, inaccessibles.

\medskip
Ce que personne ne sait encore ce matin-là\footnote{Cas fictif établi pour le
cours. Les noms (LiZbiZ, NYONG SÉCURITÉ, BANQUE SANAGA), les dates et les
montants sont fictifs~; les mécanismes techniques et les références
normatives sont exacts.}~: l'attaque n'a pas commencé cette nuit. Les
premières traces, dans les journaux, remontent à onze jours plus tôt. Entre
les deux, un compte de service compromis, un déplacement discret d'une
machine à l'autre, une exfiltration de 40~Go, puis le chiffrement. Pendant
onze jours, tout était dans les journaux. Personne ne les lisait.

% ============================================================================
\section{Trois métiers, une seule journée}
\label{sec:p00-metiers}
% ============================================================================

Cette journée met en scène les trois compétences du cours, dans l'ordre où
elles interviennent.

\begin{tableaucours}{L{3.2cm}L{4.4cm}L{4.6cm}}{%
    \tableauHeader{Moment} &
    \tableauHeader{Compétence} &
    \tableauHeader{Question posée}
}
Avant (onze jours) & Supervision, détection & L'a-t-on vu venir~? \\
Pendant (le constat) & Réponse à incident & Comment limiter les dégâts~? \\
Après (l'enquête) & Analyse forensique & Que s'est-il exactement passé, et peut-on le prouver~? \\
\end{tableaucours}
\vspace{-0.8em}
\begin{center}{\small\itshape Tableau~0.1. Les trois temps d'un incident}\end{center}

\begin{boxscenario}{Opération LOGONE. La journée du 6 mars}
\textbf{Constatation.} À 7~h~15, quatre systèmes chiffrés, sauvegardes
hebdomadaires inaccessibles, un message de rançon. La production de
l'entrepôt de Garoua est arrêtée.

\medskip
\textbf{Analyse} (confiance moyenne au matin, élevée après enquête).
L'incident n'est pas un simple chiffrement~: c'est l'aboutissement d'une
intrusion de onze jours. La distinction entre ce que l'on \emph{constate}
(les écrans noirs) et ce que l'on \emph{suppose} puis \emph{établit} (le
chemin de l'attaquant) est le cœur du métier.
\end{boxscenario}

% ============================================================================
\section{Le fil rouge : LiZbiZ et ses opérations}
\label{sec:p00-filrouge}
% ============================================================================

Tout le cours suit la même entreprise, \sigle{LiZbiZ}, transporteur de colis
présent à Yaoundé, Douala et Garoua, et le cabinet \sigle{NYONG SÉCURITÉ}
qui l'assiste. Trois opérations, situées dans le Nord, structurent la
pratique~:

\begin{tableaucours}{L{3.0cm}L{9.2cm}}{%
    \tableauHeader{Opération} &
    \tableauHeader{Ce qu'elle illustre}
}
BÉNOUÉ & La supervision et la détection~: ce que le SIEM a vu, et pas vu, pendant les onze jours \\
FARO & L'analyse forensique~: disque, mémoire, réseau, et la preuve opposable \\
MAYO & La réponse à incident~: confinement, éradication, rétablissement, retour d'expérience \\
\end{tableaucours}
\vspace{-0.8em}
\begin{center}{\small\itshape Tableau~0.2. Les trois opérations du fil rouge}\end{center}

\begin{tableaucours}{L{3.4cm}L{8.8cm}}{%
    \tableauHeader{Acteur} &
    \tableauHeader{Rôle}
}
M.~KOUNOU & Responsable informatique de LiZbiZ \\
Mme~NKENG & Analyste SOC, cabinet NYONG SÉCURITÉ \\
M.~BIKORO & Analyste forensique et chef de mission \\
Mme~MOUAFO & Responsable de la réponse à incident \\
M.~ETEME & Directeur général de LiZbiZ \\
\end{tableaucours}
\vspace{-0.8em}
\begin{center}{\small\itshape Tableau~0.3. Les acteurs de l'opération LOGONE}\end{center}

% ============================================================================
\section{Le Trilemme CRO}
\label{sec:p00-cro}
% ============================================================================

Ce cours partage avec l'investigation numérique un fil conducteur théorique~:
toute trace, tout acte d'enquête se juge sur trois dimensions à la fois.

\begin{boxCRO}
\textbf{Le Trilemme CRO.} \termecle{Confidentialité}~: la trace et l'enquête
ne doivent pas exposer plus de données que nécessaire. \termecle{Fiabilité}~:
la trace est-elle fidèle, complète, non altérée~? \termecle{Opposabilité}~:
pourra-t-on s'en prévaloir devant un tiers (direction, ANTIC, juge)~?

\smallskip
Une action d'enquête améliore souvent une dimension au détriment d'une autre.
On la note \cro{+}{-}{0}~: ici, confidentialité préservée, fiabilité
dégradée, opposabilité inchangée. Le professionnel arbitre en conscience et
\emph{documente} son arbitrage.
\end{boxCRO}

% ============================================================================
\section{Les questions du cours}
\label{sec:p00-questions}
% ============================================================================

\begin{boxinfo}
\textbf{Ce que le cours apprend à répondre.}
\begin{itemize}
    \item Comment \textbf{superviser} un système pour voir l'attaque pendant
          qu'elle se déroule (Parties~I et~II)~?
    \item Que disent les \textbf{normes et le droit} sur la preuve et le
          signalement (Partie~III)~?
    \item Comment \textbf{investiguer} proprement disque, mémoire et réseau
          (Partie~IV)~?
    \item Comment \textbf{répondre} à l'incident et s'en relever (Partie~V)~?
    \item Comment \textbf{apprendre} et mesurer, pour que la prochaine fois
          quelqu'un voie venir l'attaque (Partie~VI)~?
\end{itemize}
\end{boxinfo}

\separateur

\begin{boxresume}
\textbf{Ce qu'il faut retenir de ce prologue~:}
\begin{itemize}
    \item \textbf{L'écran de rançon} est la fin d'une histoire, pas le début.
    \item \textbf{Trois temps}~: voir venir (supervision), limiter (réponse),
          établir (forensique).
    \item \textbf{Le fil rouge} LiZbiZ et les opérations BÉNOUÉ, FARO, MAYO.
    \item \textbf{Le Trilemme CRO}~: confidentialité, fiabilité, opposabilité.
    \item \textbf{Constater} n'est pas \textbf{supposer}~; on sépare toujours
          les deux.
\end{itemize}
\end{boxresume}

\begin{boxinfo}
\textbf{Transition.} Pour voir venir l'attaque, il faut d'abord comprendre
pourquoi et comment on supervise~: c'est l'objet du
Chapitre~\ref{chap:pourquoi}.
\end{boxinfo}
}{}

    \part{Superviser : le Centre d'Opérations de Sécurité}
    \label{part:superviser}

        \IfFileExists{partie1/P01_C01_pourquoi_superviser.tex}{% ============================================================================
% Fichier  : partie1/P01_C01_pourquoi_superviser.tex
% Titre    : Pourquoi Superviser
% Auteur   : MINKA MI NGUIDJOÏ Thierry Emmanuel
% ============================================================================
% Sources : ISO/IEC 27035 ; NIST SP 800-61r2 ; CIS Controls v8. Opération
% BÉNOUÉ (LiZbiZ) : cas fictif.
% ============================================================================

\chapter{Pourquoi Superviser}
\label{chap:pourquoi}
\chaptermark{Pourquoi superviser}

\epigraph{%
    « On ne défend pas ce que l'on ne voit pas. »%
}{Principe du cours}

\niveauChapitre{Fondamental}{Le prologue~; notions de réseaux et de systèmes.}

\begin{boxobjectifs}
\begin{itemize}
    \item Définir la supervision de sécurité et la distinguer de la
          supervision d'exploitation.
    \item Situer la détection dans la chaîne prévention, détection, réaction.
    \item Comprendre le délai de présence de l'attaquant et pourquoi il
          compte.
    \item Relier la supervision aux obligations de l'entreprise.
    \item Poser les indicateurs fondamentaux de la détection.
\end{itemize}
\end{boxobjectifs}

\bigskip

% ============================================================================
\section*{Question directrice : si tout était dans les journaux, pourquoi personne n'a-t-il rien vu~?}
% ============================================================================

Onze jours de présence, des traces à chaque étape, et pourtant la découverte
par un écran de rançon. Ce chapitre explique pourquoi la prévention ne suffit
pas, ce qu'est superviser, et comment se mesure la capacité à voir venir.

\medskip
\textbf{Plan du chapitre.}
\begin{enumerate}
    \item Prévenir ne suffit pas (section~\ref{sec:c01-prevenir}).
    \item Qu'est-ce que superviser (section~\ref{sec:c01-definir}).
    \item Le délai de présence (section~\ref{sec:c01-delai}).
    \item Superviser, une obligation (section~\ref{sec:c01-obligation}).
    \item Les indicateurs fondamentaux (section~\ref{sec:c01-indicateurs}).
\end{enumerate}

\separateur

% ============================================================================
\section{Prévenir ne Suffit Pas}
\label{sec:c01-prevenir}
% ============================================================================

La sécurité se joue sur trois temps~: \termecle{prévention} (empêcher),
\termecle{détection} (voir ce qui est passé malgré la prévention) et
\termecle{réaction} (limiter et corriger). Aucune prévention n'est parfaite~:
un correctif manquant, un mot de passe faible, un courriel piégé suffisent.
La détection est le filet qui rattrape ce que la prévention a laissé passer.

% ============================================================================
\section{Qu'est-ce que Superviser}
\label{sec:c01-definir}
% ============================================================================

\begin{boxdef}
\termecle{Supervision de sécurité}~: collecte continue des traces d'un
système, leur corrélation et leur analyse pour détecter, le plus tôt
possible, les activités malveillantes ou anormales. Elle se distingue de la
\termecle{supervision d'exploitation}, qui surveille la disponibilité et la
performance, non les attaques.
\end{boxdef}

% ============================================================================
\section{Le Délai de Présence}
\label{sec:c01-delai}
% ============================================================================

\begin{boxdef}
\termecle{Délai de présence} (temps de séjour)~: durée entre la première
compromission et sa détection. Plus il est long, plus l'attaquant s'étend,
exfiltre et s'installe.
\end{boxdef}

\begin{boxscenario}{Opération BÉNOUÉ. Onze jours}
\textbf{Constatation.} Première trace le 23~février (connexion d'un compte de
service à une heure inhabituelle)~; chiffrement dans la nuit du 5 au 6 mars. Délai de présence~:
onze jours.

\medskip
\textbf{Analyse} (confiance élevée). Réduire ce délai de onze jours à
quelques heures aurait permis de contenir l'attaque avant l'exfiltration et
le chiffrement. C'est l'objectif premier de la supervision.
\end{boxscenario}

% ============================================================================
\section{Superviser, une Obligation}
\label{sec:c01-obligation}
% ============================================================================

\begin{boxjuris}
Au Cameroun, la loi n°~2010/012 du 21~décembre~2010 relative à la
cybersécurité et à la cybercriminalité impose la conservation des traces de
connexion et la coopération avec l'\sigle{ANTIC} en cas d'incident
(Chapitre~\ref{chap:droit}). Ne pas disposer de journaux exploitables n'est
pas seulement une faiblesse technique~: c'est un manquement opposable.
\end{boxjuris}

% ============================================================================
\section{Les Indicateurs Fondamentaux}
\label{sec:c01-indicateurs}
% ============================================================================

\begin{tableaucours}{L{4.4cm}L{7.8cm}}{%
    \tableauHeader{Indicateur} &
    \tableauHeader{Ce qu'il mesure}
}
Délai moyen de détection & Temps entre compromission et alerte \\
Délai moyen de réaction & Temps entre alerte et confinement \\
Couverture de journalisation & Part des systèmes dont les traces sont collectées \\
Taux de faux positifs & Alertes sans incident réel \\
Taux de détection & Part des attaques (ou techniques) effectivement vues \\
\end{tableaucours}
\vspace{-0.8em}
\begin{center}{\small\itshape Tableau~1.1. Indicateurs fondamentaux de la supervision}\end{center}

\begin{boxmetier}{Analyste SOC}
\textbf{Mesurer pour progresser.} L'analyste ne se contente pas de traiter
les alertes~: il suit le délai de détection d'un mois sur l'autre et
l'améliore. Un SOC qui ne mesure pas son délai de détection ne sait pas s'il
s'améliore.
\end{boxmetier}

\separateur

\begin{boxresume}
\textbf{Ce qu'il faut retenir de ce chapitre~:}
\begin{itemize}
    \item \textbf{Prévention, détection, réaction}~: la détection rattrape ce
          que la prévention laisse passer.
    \item \textbf{Superviser} la sécurité n'est pas superviser l'exploitation.
    \item \textbf{Le délai de présence} est l'ennemi~; le réduire est
          l'objectif.
    \item \textbf{Superviser est une obligation} (loi 2010/012).
    \item \textbf{On mesure} délai de détection, de réaction, couverture.
\end{itemize}
\end{boxresume}

\begin{boxinfo}
\textbf{Points de contrôle de l'élève.} Je sais~: définir la supervision de
sécurité~; situer la détection dans la chaîne~; expliquer le délai de
présence~; rappeler l'obligation légale~; citer les indicateurs.
\end{boxinfo}

\bigskip

\begin{boxexo}[title={\ding{45}\quad Exercice~1.1. Trois temps \niveaubase}]
Classez~: installer un correctif~; lire les journaux chaque matin~; isoler un
poste infecté. À quel temps (prévention, détection, réaction) chacun
appartient-il~?
\end{boxexo}

\begin{boxexo}[title={\ding{45}\quad Exercice~1.2. Deux supervisions \niveaubase}]
Distinguez supervision de sécurité et supervision d'exploitation par un
exemple de chacune.
\end{boxexo}

\begin{boxexo}[title={\ding{45}\quad Exercice~1.3. Le délai \niveaubase}]
Calculez le délai de présence de l'opération BÉNOUÉ. Pourquoi le réduire~?
\end{boxexo}

\begin{boxexo}[title={\ding{45}\quad Exercice~1.4. Le filet \niveaumoyen}]
« Puisque nous avons un pare-feu et un antivirus, la détection est inutile. »
Répondez.
\end{boxexo}

\begin{boxexo}[title={\ding{45}\quad Exercice~1.5. Obligation \niveaumoyen}]
En quoi l'absence de journaux exploitables est-elle un manquement opposable
chez LiZbiZ~?
\end{boxexo}

\begin{boxexo}[title={\ding{45}\quad Exercice~1.6. Indicateurs \niveaumoyen}]
Pour LiZbiZ, proposez une valeur cible de délai de détection et justifiez-la.
\end{boxexo}

\begin{boxexo}[title={\ding{45}\quad Exercice~1.7. Couverture \niveaumoyen}]
Pourquoi une couverture de journalisation partielle crée-t-elle des angles
morts~?
\end{boxexo}

\begin{boxexo}[title={\ding{45}\quad Exercice~1.8. Faux positifs \niveaumoyen}]
Quel risque un taux de faux positifs trop élevé fait-il courir à la
détection~?
\end{boxexo}

\begin{boxexo}[title={\ding{45}\quad Exercice~1.9. Tableau de bord \niveauexpert}]
Construisez le tableau de bord de supervision que LiZbiZ devrait suivre après
l'opération GAROUA.
\end{boxexo}

\begin{boxexo}[title={\ding{45}\quad Exercice~1.10. Note à la direction \niveauexpert}]
Rédigez une note qui explique à M.~ETEME pourquoi investir dans la détection,
et non seulement dans la prévention.
\end{boxexo}

\bigskip

\begin{boxinfo}
\textbf{Pour aller plus loin.}
\begin{itemize}
    \item ISO/IEC 27035~\cite{iso27035}~; NIST SP~800-61r2~\cite{nist80061r2}.
\end{itemize}
\end{boxinfo}

\begin{boxinfo}
\textbf{Transition.} Superviser suppose une organisation~: le
Chapitre~\ref{chap:soc} présente l'architecture d'un centre d'opérations de
sécurité.
\end{boxinfo}
}{}
        \IfFileExists{partie1/P01_C02_architecture_soc.tex}{% ============================================================================
% Fichier  : partie1/P01_C02_architecture_soc.tex
% Titre    : L'Architecture d'un SOC
% Auteur   : MINKA MI NGUIDJOÏ Thierry Emmanuel
% ============================================================================
% Sources : NIST SP 800-61r2 ; modèles de SOC (niveaux N1-N3) ; CIS Controls v8.
% ============================================================================

\chapter{L'Architecture d'un SOC}
\label{chap:soc}
\chaptermark{L'architecture d'un SOC}

\epigraph{%
    « Un SOC n'est pas une salle pleine d'écrans~;\\
    c'est un processus qui transforme des traces en décisions. »%
}{Principe du cours}

\niveauChapitre{Fondamental}{Chapitre~\ref{chap:pourquoi}.}

\begin{boxobjectifs}
\begin{itemize}
    \item Définir le centre d'opérations de sécurité (SOC) et sa mission.
    \item Décrire les niveaux d'analyse (N1, N2, N3) et leurs rôles.
    \item Comprendre les modèles d'organisation (interne, infogéré, hybride).
    \item Situer les fonctions chasse, ingénierie de détection, réponse.
    \item Dimensionner un SOC réaliste pour une PME.
\end{itemize}
\end{boxobjectifs}

\bigskip

% ============================================================================
\section*{Question directrice : de quelle équipe LiZbiZ a-t-elle besoin pour voir venir la prochaine attaque~?}
% ============================================================================

LiZbiZ n'a ni SOC ni analyste dédié~; M.~KOUNOU fait tout. Ce chapitre montre
ce qu'est un SOC, ses rôles et ses modèles, et comment en bâtir un
proportionné à une PME, seul ou avec le cabinet NYONG SÉCURITÉ.

\medskip
\textbf{Plan du chapitre.}
\begin{enumerate}
    \item La mission du SOC (section~\ref{sec:c02-mission}).
    \item Les niveaux d'analyse (section~\ref{sec:c02-niveaux}).
    \item Les modèles d'organisation (section~\ref{sec:c02-modeles}).
    \item Les fonctions transverses (section~\ref{sec:c02-fonctions}).
    \item Dimensionner pour une PME (section~\ref{sec:c02-pme}).
\end{enumerate}

\separateur

% ============================================================================
\section{La Mission du SOC}
\label{sec:c02-mission}
% ============================================================================

\begin{boxdef}
\termecle{SOC} (\textit{Security Operations Center})~: fonction, humaine et
technique, chargée de détecter, analyser et traiter les événements de
sécurité en continu, à partir des traces collectées par le \sigle{SIEM} et
les sondes.
\end{boxdef}

% ============================================================================
\section{Les Niveaux d'Analyse}
\label{sec:c02-niveaux}
% ============================================================================

\begin{tableaucours}{L{2.2cm}L{4.6cm}L{5.4cm}}{%
    \tableauHeader{Niveau} &
    \tableauHeader{Rôle} &
    \tableauHeader{Ce qu'il fait}
}
N1 & Surveillance, tri & Reçoit les alertes, écarte les faux positifs, qualifie, escalade \\
N2 & Investigation & Analyse l'alerte qualifiée, reconstitue, décide du traitement \\
N3 & Expertise, chasse & Forensique, chasse proactive, ingénierie de détection \\
\end{tableaucours}
\vspace{-0.8em}
\begin{center}{\small\itshape Tableau~2.1. Les niveaux d'un SOC}\end{center}

\begin{boxalert}
\textbf{Erreur fréquente.} Vouloir trois niveaux dès le premier jour.

\smallskip
\textbf{Règle~:} une PME commence souvent par un N1 outillé et un N2 à temps
partiel, l'expertise N3 étant externalisée. L'organisation suit la maturité,
pas l'inverse.
\end{boxalert}

% ============================================================================
\section{Les Modèles d'Organisation}
\label{sec:c02-modeles}
% ============================================================================

\begin{tableaucours}{L{3.0cm}L{4.8cm}L{4.4cm}}{%
    \tableauHeader{Modèle} &
    \tableauHeader{Principe} &
    \tableauHeader{Pour qui}
}
Interne & Équipe propre, 24/7 ou heures ouvrées & Grandes organisations \\
Infogéré (MSSP) & Un prestataire opère le SOC & PME sans équipe \\
Hybride & N1/N2 infogéré, N3 et décision internes & Cas fréquent et équilibré \\
\end{tableaucours}
\vspace{-0.8em}
\begin{center}{\small\itshape Tableau~2.2. Modèles d'organisation d'un SOC}\end{center}

% ============================================================================
\section{Les Fonctions Transverses}
\label{sec:c02-fonctions}
% ============================================================================

Au-delà des niveaux, un SOC mûr distingue trois fonctions~: la
\termecle{chasse} (rechercher activement ce qui n'a pas déclenché d'alerte,
Chapitre~\ref{chap:renseignement}), l'\termecle{ingénierie de détection}
(écrire et maintenir les règles, Chapitre~\ref{chap:regles}) et la
\termecle{réponse} (traiter l'incident, Partie~V). Elles se nourrissent l'une
l'autre~: chaque incident traité améliore les règles.

% ============================================================================
\section{Dimensionner pour une PME}
\label{sec:c02-pme}
% ============================================================================

\begin{boxscenario}{Opération BÉNOUÉ. Le SOC qui n'existait pas}
\textbf{Constatation.} LiZbiZ n'avait aucune fonction de détection~: pas de
SIEM, pas d'analyste, des journaux dispersés et non lus.

\medskip
\textbf{Analyse} (confiance élevée). Un modèle hybride convient~: le cabinet
NYONG SÉCURITÉ opère la surveillance (N1/N2) sur un SIEM mutualisé, LiZbiZ
garde la décision et le lien avec l'ANTIC. Le coût est compatible avec une
PME, et le délai de présence tombe de onze jours à quelques heures.
\end{boxscenario}

\begin{boxmetier}{RSSI}
\textbf{Choisir un modèle soutenable.} Le RSSI ne vise pas le SOC le plus
complet, mais celui que l'organisation peut financer et faire vivre dans la
durée. Un SOC interne sous-doté détecte moins bien qu'un infogéré bien piloté.
\end{boxmetier}

\separateur

\begin{boxresume}
\textbf{Ce qu'il faut retenir de ce chapitre~:}
\begin{itemize}
    \item \textbf{Le SOC} transforme des traces en décisions, en continu.
    \item \textbf{N1, N2, N3}~: tri, investigation, expertise.
    \item \textbf{Interne, infogéré, hybride}~: le modèle suit la maturité et
          le budget.
    \item \textbf{Chasse, ingénierie de détection, réponse}~: les fonctions
          transverses.
    \item \textbf{Pour une PME}, l'hybride est souvent le bon départ.
\end{itemize}
\end{boxresume}

\begin{boxinfo}
\textbf{Points de contrôle de l'élève.} Je sais~: définir un SOC~; décrire
N1/N2/N3~; comparer les modèles~; situer chasse, ingénierie, réponse~;
proposer un SOC pour une PME.
\end{boxinfo}

\bigskip

\begin{boxexo}[title={\ding{45}\quad Exercice~2.1. Les niveaux \niveaubase}]
Associez à N1, N2 ou N3~: écarter un faux positif~; faire du forensique~;
décider d'isoler un serveur.
\end{boxexo}

\begin{boxexo}[title={\ding{45}\quad Exercice~2.2. SOC ou exploitation \niveaubase}]
Pourquoi un SOC n'est-il pas « une salle pleine d'écrans »~?
\end{boxexo}

\begin{boxexo}[title={\ding{45}\quad Exercice~2.3. Modèles \niveaubase}]
Citez un avantage et une limite de chaque modèle (interne, infogéré,
hybride).
\end{boxexo}

\begin{boxexo}[title={\ding{45}\quad Exercice~2.4. Maturité \niveaumoyen}]
Pourquoi déconseille-t-on à une PME de monter d'emblée trois niveaux
internes~?
\end{boxexo}

\begin{boxexo}[title={\ding{45}\quad Exercice~2.5. Hybride \niveaumoyen}]
Dans le modèle hybride de LiZbiZ, qui garde la décision~? Pourquoi~?
\end{boxexo}

\begin{boxexo}[title={\ding{45}\quad Exercice~2.6. Fonctions \niveaumoyen}]
Expliquez comment un incident traité améliore les règles de détection.
\end{boxexo}

\begin{boxexo}[title={\ding{45}\quad Exercice~2.7. Escalade \niveaumoyen}]
Décrivez le parcours d'une alerte de N1 à N2, avec les critères d'escalade.
\end{boxexo}

\begin{boxexo}[title={\ding{45}\quad Exercice~2.8. Chasse \niveaumoyen}]
En quoi la chasse diffère-t-elle du traitement d'alerte~?
\end{boxexo}

\begin{boxexo}[title={\ding{45}\quad Exercice~2.9. Proposition \niveauexpert}]
Rédigez une proposition de SOC hybride pour LiZbiZ~: périmètre, rôles,
horaires, responsabilités partagées avec NYONG SÉCURITÉ.
\end{boxexo}

\begin{boxexo}[title={\ding{45}\quad Exercice~2.10. Convention de service \niveauexpert}]
Quels engagements de service (délais, couverture, rapports) inscririez-vous
au contrat avec un SOC infogéré~?
\end{boxexo}

\bigskip

\begin{boxinfo}
\textbf{Pour aller plus loin.}
\begin{itemize}
    \item NIST SP~800-61r2~\cite{nist80061r2}~; CIS Controls v8~\cite{cis2021controls}.
\end{itemize}
\end{boxinfo}

\begin{boxinfo}
\textbf{Transition.} Un SOC se pilote~: le Chapitre~\ref{chap:gouvernance}
traite de sa gouvernance et de ses indicateurs.
\end{boxinfo}
}{}
        \IfFileExists{partie1/P01_C03_gouvernance_indicateurs.tex}{% ============================================================================
% Fichier  : partie1/P01_C03_gouvernance_indicateurs.tex
% Titre    : Gouvernance et Indicateurs de la Supervision
% Auteur   : MINKA MI NGUIDJOÏ Thierry Emmanuel
% ============================================================================
% Sources : ISO/IEC 27004 (mesure) ; ISO/IEC 27035 ; CIS Controls v8.
% ============================================================================

\chapter{Gouvernance et Indicateurs de la Supervision}
\label{chap:gouvernance}
\chaptermark{Gouvernance et indicateurs}

\epigraph{%
    « Ce qui ne se mesure pas ne se pilote pas~;\\
    ce qui ne se pilote pas se dégrade. »%
}{Principe du cours}

\niveauChapitre{Fondamental}{Chapitres~\ref{chap:pourquoi} et~\ref{chap:soc}.}

\begin{boxobjectifs}
\begin{itemize}
    \item Définir la gouvernance d'un dispositif de supervision.
    \item Distinguer indicateurs d'activité, de performance et de résultat.
    \item Construire un tableau de bord utile et honnête.
    \item Fixer des objectifs de délai (détection, réaction).
    \item Rendre compte à la direction sans noyer sous la technique.
\end{itemize}
\end{boxobjectifs}

\bigskip

% ============================================================================
\section*{Question directrice : comment savoir si la supervision de LiZbiZ s'améliore vraiment~?}
% ============================================================================

Après l'incident, LiZbiZ met en place une détection. Six mois plus tard,
M.~ETEME demande~: « Sommes-nous mieux protégés~? » La réponse ne peut être
ni « oui » vague ni une avalanche de courbes~: elle tient dans quelques
indicateurs choisis et suivis dans le temps.

\medskip
\textbf{Plan du chapitre.}
\begin{enumerate}
    \item Gouverner la supervision (section~\ref{sec:c03-gouverner}).
    \item Trois familles d'indicateurs (section~\ref{sec:c03-familles}).
    \item Le tableau de bord (section~\ref{sec:c03-tdb}).
    \item Les objectifs de délai (section~\ref{sec:c03-delais}).
    \item Rendre compte (section~\ref{sec:c03-rendre}).
\end{enumerate}

\separateur

% ============================================================================
\section{Gouverner la Supervision}
\label{sec:c03-gouverner}
% ============================================================================

\begin{boxdef}
\termecle{Gouvernance de la supervision}~: ensemble des décisions qui fixent
le périmètre surveillé, les priorités, les moyens, les responsabilités et les
indicateurs, et qui en rendent compte à la direction. Elle relie la
technique à la stratégie.
\end{boxdef}

% ============================================================================
\section{Trois Familles d'Indicateurs}
\label{sec:c03-familles}
% ============================================================================

\begin{tableaucours}{L{3.2cm}L{4.6cm}L{4.4cm}}{%
    \tableauHeader{Famille} &
    \tableauHeader{Question} &
    \tableauHeader{Exemple}
}
Activité & Que fait-on~? & Nombre d'alertes traitées \\
Performance & Le fait-on bien et vite~? & Délai moyen de détection \\
Résultat & Est-on plus sûr~? & Part des techniques détectées, incidents évités \\
\end{tableaucours}
\vspace{-0.8em}
\begin{center}{\small\itshape Tableau~3.1. Trois familles d'indicateurs~\cite{iso27004}}\end{center}

\begin{boxalert}
\textbf{Erreur fréquente.} Ne suivre que l'activité (« nous avons traité
3~000 alertes »).

\smallskip
\textbf{Règle~:} l'activité ne dit rien de l'efficacité. Un volume d'alertes
élevé peut signaler de mauvaises règles, non une bonne détection. On
privilégie la performance et le résultat.
\end{boxalert}

% ============================================================================
\section{Le Tableau de Bord}
\label{sec:c03-tdb}
% ============================================================================

\begin{boxprocedure}
\textbf{Construire un tableau de bord de supervision}
\begin{enumerate}
    \item Choisir cinq à sept indicateurs, pas plus.
    \item Au moins un de chaque famille (activité, performance, résultat).
    \item Pour chacun~: définition, source, cible, tendance.
    \item Montrer l'évolution, pas seulement la valeur du mois.
    \item Dire honnêtement ce qui n'est pas couvert.
\end{enumerate}
\end{boxprocedure}

% ============================================================================
\section{Les Objectifs de Délai}
\label{sec:c03-delais}
% ============================================================================

Les deux délais cardinaux sont le \termecle{délai de détection} et le
\termecle{délai de réaction}. On les fixe par type d'incident~: un
rançongiciel exige une détection en minutes et un confinement en moins d'une
heure~; une anomalie mineure tolère un traitement en jours. Ces cibles
guident le dimensionnement du SOC (Chapitre~\ref{chap:soc}).

% ============================================================================
\section{Rendre Compte}
\label{sec:c03-rendre}
% ============================================================================

\begin{boxscenario}{Opération BÉNOUÉ. Six mois après}
\textbf{Constatation.} Nouveau tableau de bord~: délai de détection médian
ramené à 2~heures, couverture de journalisation portée à 85~\%, 40~\% des
techniques ATT\&CK jouées en test désormais détectées.

\medskip
\textbf{Analyse} (confiance moyenne). Les chiffres montrent un progrès réel,
mais la couverture reste partielle (15~\% des systèmes non journalisés) et la
détection en direct limitée~: ce sont les NEK à traiter
(Chapitre~\ref{chap:mesurer}). Rendre compte, c'est aussi dire cela.
\end{boxscenario}

\begin{boxmetier}{RSSI et analyste SOC}
\textbf{Un tableau de bord par lecteur.} Le RSSI présente à la direction trois
indicateurs de résultat~; l'analyste suit en interne le détail de la
performance. Les deux disent la même vérité, à deux niveaux.
\end{boxmetier}

\separateur

\begin{boxresume}
\textbf{Ce qu'il faut retenir de ce chapitre~:}
\begin{itemize}
    \item \textbf{Gouverner}, c'est relier la supervision à la stratégie.
    \item \textbf{Activité, performance, résultat}~: trois familles.
    \item \textbf{L'activité seule} trompe~; on vise performance et résultat.
    \item \textbf{Cinq à sept indicateurs}, avec tendance et angles morts.
    \item \textbf{Des objectifs de délai} par type d'incident.
\end{itemize}
\end{boxresume}

\begin{boxinfo}
\textbf{Points de contrôle de l'élève.} Je sais~: définir la gouvernance de
la supervision~; classer un indicateur~; bâtir un tableau de bord~; fixer des
délais cibles~; rendre compte à deux niveaux.
\end{boxinfo}

\bigskip

\begin{boxexo}[title={\ding{45}\quad Exercice~3.1. Familles \niveaubase}]
Classez~: nombre d'alertes~; délai de détection~; part des techniques
détectées.
\end{boxexo}

\begin{boxexo}[title={\ding{45}\quad Exercice~3.2. Activité trompeuse \niveaubase}]
Pourquoi « 3~000 alertes traitées » ne prouve-t-il pas une bonne détection~?
\end{boxexo}

\begin{boxexo}[title={\ding{45}\quad Exercice~3.3. Combien \niveaubase}]
Combien d'indicateurs un tableau de bord devrait-il comporter, et pourquoi
pas plus~?
\end{boxexo}

\begin{boxexo}[title={\ding{45}\quad Exercice~3.4. Cibles \niveaumoyen}]
Proposez des délais cibles de détection et de réaction pour un rançongiciel,
et justifiez.
\end{boxexo}

\begin{boxexo}[title={\ding{45}\quad Exercice~3.5. Lire BÉNOUÉ \niveaumoyen}]
Commentez les chiffres de l'opération BÉNOUÉ à six mois~: progrès et limites.
\end{boxexo}

\begin{boxexo}[title={\ding{45}\quad Exercice~3.6. Angles morts \niveaumoyen}]
Pourquoi faut-il afficher ce qui n'est pas couvert dans un tableau de bord~?
\end{boxexo}

\begin{boxexo}[title={\ding{45}\quad Exercice~3.7. Deux lecteurs \niveaumoyen}]
Quels indicateurs montrez-vous à la direction~? Lesquels gardez-vous en
interne~?
\end{boxexo}

\begin{boxexo}[title={\ding{45}\quad Exercice~3.8. Tendance \niveaumoyen}]
Pourquoi la tendance importe-t-elle plus que la valeur d'un mois isolé~?
\end{boxexo}

\begin{boxexo}[title={\ding{45}\quad Exercice~3.9. Tableau de bord \niveauexpert}]
Construisez le tableau de bord complet de supervision de LiZbiZ~: sept
indicateurs, définition, source, cible.
\end{boxexo}

\begin{boxexo}[title={\ding{45}\quad Exercice~3.10. Présentation au comité \niveauexpert}]
Rédigez la page que le RSSI présente au comité de direction pour répondre à
« Sommes-nous mieux protégés~? ».
\end{boxexo}

\bigskip

\begin{boxinfo}
\textbf{Pour aller plus loin.}
\begin{itemize}
    \item ISO/IEC 27004~\cite{iso27004}~; ISO/IEC 27035~\cite{iso27035}.
\end{itemize}
\end{boxinfo}

\begin{boxinfo}
\textbf{Transition.} La supervision commence par les traces~: la Partie~II
s'ouvre sur les journaux (Chapitre~\ref{chap:journaux}).
\end{boxinfo}
}{}

    \part{Collecter et Détecter}
    \label{part:detecter}

        \IfFileExists{partie2/P02_C04_journaux.tex}{% ============================================================================
% Fichier  : partie2/P02_C04_journaux.tex
% Titre    : Les Journaux
% Auteur   : MINKA MI NGUIDJOÏ Thierry Emmanuel
% ============================================================================
% Sources : RFC 3227 ; NIST SP 800-92 (gestion des journaux) ; loi 2010/012.
% ============================================================================

\chapter{Les Journaux}
\label{chap:journaux}
\chaptermark{Les journaux}

\epigraph{%
    « Un journal que l'on ne collecte pas est une mémoire que l'on efface. »%
}{Principe du cours}

\niveauChapitre{Fondamental}{Chapitre~\ref{chap:pourquoi}~; notions de
systèmes et de réseaux.}

\begin{boxobjectifs}
\begin{itemize}
    \item Identifier les sources de journaux utiles à la détection.
    \item Comprendre la normalisation et l'horodatage.
    \item Fixer une politique de collecte et de conservation.
    \item Relier la journalisation aux obligations légales.
    \item Repérer les angles morts d'une journalisation partielle.
\end{itemize}
\end{boxobjectifs}

\bigskip

% ============================================================================
\section*{Question directrice : parmi tous les journaux de LiZbiZ, lesquels auraient révélé l'intrusion~?}
% ============================================================================

L'attaque a laissé des traces~: une connexion anormale d'un compte de
service, une création de compte, un service installé à distance. Encore
fallait-il que ces journaux existent, soient collectés et conservés. Ce
chapitre recense les sources, leur normalisation et leur conservation.

\medskip
\textbf{Plan du chapitre.}
\begin{enumerate}
    \item Les sources de journaux (section~\ref{sec:c04-sources}).
    \item Normaliser et horodater (section~\ref{sec:c04-normaliser}).
    \item Collecter et conserver (section~\ref{sec:c04-conserver}).
    \item L'obligation de conservation (section~\ref{sec:c04-obligation}).
    \item Les angles morts (section~\ref{sec:c04-angles}).
\end{enumerate}

\separateur

% ============================================================================
\section{Les Sources de Journaux}
\label{sec:c04-sources}
% ============================================================================

\begin{tableaucours}{L{3.6cm}L{8.6cm}}{%
    \tableauHeader{Source} &
    \tableauHeader{Ce qu'elle révèle}
}
Journaux du domaine (Windows) & Connexions, créations de comptes, services, tâches \\
Journaux système (Linux) & Authentifications, sudo, démons, cron \\
Pare-feu et mandataire & Connexions, flux refusés, destinations \\
Serveurs d'application & Accès, erreurs, requêtes \\
Messagerie & Courriels reçus, pièces jointes, signalements \\
Équipements (imprimantes, bornes) & Accès d'administration \\
\end{tableaucours}
\vspace{-0.8em}
\begin{center}{\small\itshape Tableau~4.1. Sources de journaux utiles à la détection}\end{center}

% ============================================================================
\section{Normaliser et Horodater}
\label{sec:c04-normaliser}
% ============================================================================

Des journaux hétérogènes ne se corrèlent que s'ils sont \termecle{normalisés}
(ramenés à un format commun~: qui, quoi, quand, où) et surtout
\termecle{horodatés} de façon cohérente. Une source d'horloge commune (NTP)
est indispensable~: sans elle, impossible de reconstituer une chronologie,
et la preuve perd sa fiabilité.

\begin{boxCRO}
\textbf{Analyse CRO de l'horodatage.} Un horodatage fiable et synchronisé
sert la \cro{0}{+}{+}~: il ne change rien à la confidentialité, améliore la
fiabilité de la trace et son opposabilité devant un tiers. Un horodatage
incohérent ruine l'opposabilité d'une chronologie, même exacte.
\end{boxCRO}

% ============================================================================
\section{Collecter et Conserver}
\label{sec:c04-conserver}
% ============================================================================

\begin{boxprocedure}
\textbf{Politique de collecte des journaux}
\begin{enumerate}
    \item Lister les sources et leur criticité.
    \item Centraliser les journaux hors des machines surveillées (un
          attaquant efface les journaux locaux).
    \item Protéger l'intégrité (accès restreint, écriture seule, condensats).
    \item Synchroniser les horloges (NTP).
    \item Fixer une durée de conservation conforme au droit et aux besoins.
\end{enumerate}
\end{boxprocedure}

% ============================================================================
\section{L'Obligation de Conservation}
\label{sec:c04-obligation}
% ============================================================================

\begin{boxjuris}
La loi n°~2010/012 et la loi n°~2010/013 sur les communications électroniques
imposent aux opérateurs et aux entreprises la conservation des données de
connexion et leur mise à disposition des autorités. Les durées et modalités
précises se lisent dans le texte publié~; le cours en retient le principe~:
des journaux fiables, conservés, sont une obligation autant qu'un outil.
\end{boxjuris}

% ============================================================================
\section{Les Angles Morts}
\label{sec:c04-angles}
% ============================================================================

\begin{boxscenario}{Opération BÉNOUÉ. Ce qui n'était pas journalisé}
\textbf{Constatation.} Les contrôleurs de domaine journalisaient les
connexions, mais le serveur de fichiers \texttt{FS-01} n'envoyait rien au
central, et l'horloge d'un poste de Garoua avançait de 40~minutes.

\medskip
\textbf{Analyse} (confiance élevée). L'intrusion a laissé des traces
visibles sur le domaine, mais l'exfiltration depuis \texttt{FS-01} est restée
invisible faute de collecte, et l'écart d'horloge a brouillé la chronologie.
Deux angles morts à combler en priorité.
\end{boxscenario}

\begin{boxmetier}{Analyste SOC}
\textbf{Cartographier avant de corréler.} L'analyste tient à jour la carte
des sources journalisées et de leurs manques~: une détection n'est fiable que
sur un périmètre connu. Ce qui n'est pas collecté doit être dit, non ignoré.
\end{boxmetier}

\separateur

\begin{boxresume}
\textbf{Ce qu'il faut retenir de ce chapitre~:}
\begin{itemize}
    \item \textbf{Les sources}~: domaine, système, pare-feu, applications,
          messagerie, équipements.
    \item \textbf{Normaliser et horodater} conditionnent la corrélation.
    \item \textbf{Centraliser et protéger} les journaux hors des machines
          surveillées.
    \item \textbf{Conserver} est une obligation légale.
    \item \textbf{Les angles morts} (sources non collectées, horloges
          dérivées) sont des priorités.
\end{itemize}
\end{boxresume}

\begin{boxinfo}
\textbf{Points de contrôle de l'élève.} Je sais~: citer les sources de
journaux~; expliquer normalisation et horodatage~; rédiger une politique de
collecte~; rappeler l'obligation~; repérer un angle mort.
\end{boxinfo}

\bigskip

\begin{boxexo}[title={\ding{45}\quad Exercice~4.1. Sources \niveaubase}]
Quelle source journalise la création d'un compte~? L'exfiltration depuis un
partage de fichiers~?
\end{boxexo}

\begin{boxexo}[title={\ding{45}\quad Exercice~4.2. Horloge \niveaubase}]
Pourquoi une horloge commune (NTP) est-elle indispensable~?
\end{boxexo}

\begin{boxexo}[title={\ding{45}\quad Exercice~4.3. Centraliser \niveaubase}]
Pourquoi centraliser les journaux hors des machines surveillées~?
\end{boxexo}

\begin{boxexo}[title={\ding{45}\quad Exercice~4.4. Normaliser \niveaumoyen}]
Qu'apporte la normalisation à la corrélation de journaux hétérogènes~?
\end{boxexo}

\begin{boxexo}[title={\ding{45}\quad Exercice~4.5. CRO \niveaumoyen}]
Expliquez l'analyse CRO de l'horodatage de l'opération BÉNOUÉ.
\end{boxexo}

\begin{boxexo}[title={\ding{45}\quad Exercice~4.6. Angle mort \niveaumoyen}]
Pourquoi l'exfiltration depuis \texttt{FS-01} est-elle restée invisible~?
\end{boxexo}

\begin{boxexo}[title={\ding{45}\quad Exercice~4.7. Conservation \niveaumoyen}]
Quels critères fixent la durée de conservation des journaux~?
\end{boxexo}

\begin{boxexo}[title={\ding{45}\quad Exercice~4.8. Intégrité \niveaumoyen}]
Comment garantir qu'un attaquant ne puisse pas effacer les journaux
centralisés~?
\end{boxexo}

\begin{boxexo}[title={\ding{45}\quad Exercice~4.9. Politique \niveauexpert}]
Rédigez la politique de journalisation de LiZbiZ~: sources, centralisation,
intégrité, conservation, responsabilités.
\end{boxexo}

\begin{boxexo}[title={\ding{45}\quad Exercice~4.10. Plan de comblement \niveauexpert}]
À partir des angles morts de BÉNOUÉ, proposez un plan de comblement priorisé.
\end{boxexo}

\bigskip

\begin{boxinfo}
\textbf{Pour aller plus loin.}
\begin{itemize}
    \item RFC~3227~\cite{rfc3227}~; loi 2010/012~\cite{loi2010012} et
          2010/013~\cite{loi2010013}.
\end{itemize}
\end{boxinfo}

\begin{boxinfo}
\textbf{Transition.} Collecter ne suffit pas~; il faut corréler~: le
Chapitre~\ref{chap:siem} traite du SIEM.
\end{boxinfo}
}{}
        \IfFileExists{partie2/P02_C05_siem.tex}{% ============================================================================
% Fichier  : partie2/P02_C05_siem.tex
% Titre    : Le SIEM
% Auteur   : MINKA MI NGUIDJOÏ Thierry Emmanuel
% ============================================================================
% Sources : NIST SP 800-92 ; CIS Controls v8 ; MITRE ATT&CK.
% ============================================================================

\chapter{Le SIEM}
\label{chap:siem}
\chaptermark{Le SIEM}

\epigraph{%
    « Le SIEM ne détecte pas~; il rapproche. Ce sont vos règles qui détectent. »%
}{Principe du cours}

\niveauChapitre{Intermédiaire}{Chapitre~\ref{chap:journaux}.}

\begin{boxobjectifs}
\begin{itemize}
    \item Définir le SIEM et son rôle dans le SOC.
    \item Décrire la chaîne collecte, normalisation, corrélation, alerte.
    \item Comprendre la notion de cas d'usage de détection.
    \item Distinguer corrélation et simple recherche.
    \item Éviter les pièges du volume et du bruit.
\end{itemize}
\end{boxobjectifs}

\bigskip

% ============================================================================
\section*{Question directrice : un SIEM aurait-il, à lui seul, arrêté l'attaque~?}
% ============================================================================

Beaucoup croient qu'installer un SIEM « rend visible » l'attaque. En réalité,
un SIEM sans cas d'usage est un entrepôt de journaux muet. Ce chapitre montre
ce qu'un SIEM fait, et ce qu'il ne fait pas sans travail humain.

\medskip
\textbf{Plan du chapitre.}
\begin{enumerate}
    \item Qu'est-ce qu'un SIEM (section~\ref{sec:c05-definir}).
    \item La chaîne de traitement (section~\ref{sec:c05-chaine}).
    \item Les cas d'usage (section~\ref{sec:c05-cas}).
    \item Corréler, pas seulement chercher (section~\ref{sec:c05-correler}).
    \item Les pièges (section~\ref{sec:c05-pieges}).
\end{enumerate}

\separateur

% ============================================================================
\section{Qu'est-ce qu'un SIEM}
\label{sec:c05-definir}
% ============================================================================

\begin{boxdef}
\termecle{SIEM} (\textit{Security Information and Event Management})~:
plateforme qui collecte les journaux de sources multiples, les normalise, les
conserve, et déclenche des alertes lorsque des règles de corrélation
reconnaissent un motif suspect.
\end{boxdef}

% ============================================================================
\section{La Chaîne de Traitement}
\label{sec:c05-chaine}
% ============================================================================

\begin{tableaucours}{L{3.0cm}L{9.2cm}}{%
    \tableauHeader{Étape} &
    \tableauHeader{Rôle}
}
Collecte & Recevoir les journaux des sources (agents, syslog, API) \\
Normalisation & Ramener à un format commun (qui, quoi, quand, où) \\
Enrichissement & Ajouter contexte (géolocalisation, criticité de l'actif) \\
Corrélation & Rapprocher plusieurs événements selon une règle \\
Alerte & Signaler à l'analyste un motif suspect, avec son contexte \\
\end{tableaucours}
\vspace{-0.8em}
\begin{center}{\small\itshape Tableau~5.1. La chaîne de traitement d'un SIEM}\end{center}

% ============================================================================
\section{Les Cas d'Usage}
\label{sec:c05-cas}
% ============================================================================

\begin{boxdef}
\termecle{Cas d'usage de détection}~: situation précise que l'on veut
détecter, traduite en une ou plusieurs règles sur des sources identifiées.
Exemple~: « un compte de service se connecte à un poste de travail hors
plage horaire habituelle ».
\end{boxdef}

Un SIEM se construit cas d'usage par cas d'usage, en priorité sur les
techniques d'attaque les plus probables (ATT\&CK). Sans cas d'usage, il ne
produit aucune alerte utile.

% ============================================================================
\section{Corréler, pas seulement Chercher}
\label{sec:c05-correler}
% ============================================================================

Chercher dans les journaux (« trouve-moi les connexions de ce compte »)
répond à une question que l'on se pose déjà. \termecle{Corréler}, c'est
rapprocher automatiquement des événements de sources différentes pour faire
émerger une alerte que personne n'a pensé à chercher~: une connexion, suivie
d'une création de compte, suivie d'un accès à un partage sensible, en
quelques minutes.

% ============================================================================
\section{Les Pièges}
\label{sec:c05-pieges}
% ============================================================================

\begin{boxalert}
\textbf{Erreur fréquente.} Tout collecter et activer toutes les règles par
défaut.

\smallskip
\textbf{Règle~:} le volume coûte cher et le bruit noie les vraies alertes. On
collecte ce qui sert des cas d'usage, on règle les seuils, on mesure le taux
de faux positifs (Chapitre~\ref{chap:gouvernance}).
\end{boxalert}

\begin{boxscenario}{Opération BÉNOUÉ. Le SIEM arrivé trop tard}
\textbf{Constatation.} Après l'incident, LiZbiZ déploie un SIEM mutualisé
chez NYONG SÉCURITÉ. Rejouées sur les journaux de février, trois règles
simples (connexion anormale d'un compte de service, création de compte hors
maintenance, service installé à distance) auraient alerté dès le 23~février.

\medskip
\textbf{Analyse} (confiance élevée). Le SIEM n'aurait rien fait seul~: ce sont
ces trois cas d'usage, écrits à la main, qui auraient réduit le délai de
présence de onze jours à quelques heures.
\end{boxscenario}

\begin{boxmetier}{Analyste SOC}
\textbf{Un SIEM se nourrit de cas d'usage.} L'analyste ne « regarde » pas le
SIEM~: il y écrit des cas d'usage, les éprouve, mesure leur bruit et les
affine. La valeur est dans les règles, pas dans l'outil.
\end{boxmetier}

\separateur

\begin{boxresume}
\textbf{Ce qu'il faut retenir de ce chapitre~:}
\begin{itemize}
    \item \textbf{Le SIEM} collecte, normalise, corrèle, alerte.
    \item \textbf{Sans cas d'usage}, il est muet.
    \item \textbf{Corréler} fait émerger ce qu'on n'aurait pas cherché.
    \item \textbf{Le volume et le bruit} sont les pièges.
    \item \textbf{La valeur} est dans les règles, pas dans l'outil.
\end{itemize}
\end{boxresume}

\begin{boxinfo}
\textbf{Points de contrôle de l'élève.} Je sais~: définir un SIEM~; décrire
sa chaîne~; écrire un cas d'usage~; distinguer corréler et chercher~; éviter
le piège du volume.
\end{boxinfo}

\bigskip

\begin{boxexo}[title={\ding{45}\quad Exercice~5.1. Chaîne \niveaubase}]
Remettez dans l'ordre~: alerte, collecte, corrélation, normalisation.
\end{boxexo}

\begin{boxexo}[title={\ding{45}\quad Exercice~5.2. Muet \niveaubase}]
Pourquoi un SIEM sans cas d'usage ne produit-il aucune alerte utile~?
\end{boxexo}

\begin{boxexo}[title={\ding{45}\quad Exercice~5.3. Cas d'usage \niveaubase}]
Rédigez un cas d'usage en une phrase pour « compte de service sur un poste
ordinaire ».
\end{boxexo}

\begin{boxexo}[title={\ding{45}\quad Exercice~5.4. Corréler \niveaumoyen}]
Donnez un exemple de corrélation à trois événements qui révèle une intrusion.
\end{boxexo}

\begin{boxexo}[title={\ding{45}\quad Exercice~5.5. Bruit \niveaumoyen}]
Pourquoi activer toutes les règles par défaut nuit-il à la détection~?
\end{boxexo}

\begin{boxexo}[title={\ding{45}\quad Exercice~5.6. BÉNOUÉ \niveaumoyen}]
Quelles trois règles auraient alerté dès le 23~février~?
\end{boxexo}

\begin{boxexo}[title={\ding{45}\quad Exercice~5.7. Enrichissement \niveaumoyen}]
Qu'apporte l'enrichissement (criticité de l'actif) à une alerte~?
\end{boxexo}

\begin{boxexo}[title={\ding{45}\quad Exercice~5.8. Chercher ou corréler \niveaumoyen}]
Classez~: « liste les connexions de svc\_sauvegarde »~; « alerte si un compte
de service accède à un partage sensible après une connexion anormale ».
\end{boxexo}

\begin{boxexo}[title={\ding{45}\quad Exercice~5.9. Catalogue \niveauexpert}]
Proposez un catalogue de dix cas d'usage prioritaires pour LiZbiZ, reliés à
des techniques ATT\&CK.
\end{boxexo}

\begin{boxexo}[title={\ding{45}\quad Exercice~5.10. Dimensionnement \niveauexpert}]
Comment arbitrer entre volume collecté, coût et couverture de détection~?
\end{boxexo}

\bigskip

\begin{boxinfo}
\textbf{Pour aller plus loin.}
\begin{itemize}
    \item CIS Controls v8~\cite{cis2021controls}~; MITRE ATT\&CK~\cite{mitreattack}.
\end{itemize}
\end{boxinfo}

\begin{boxinfo}
\textbf{Transition.} À côté du SIEM, des sondes détectent au fil du réseau et
des hôtes~: le Chapitre~\ref{chap:idsips} traite des IDS/IPS.
\end{boxinfo}
}{}
        \IfFileExists{partie2/P02_C06_ids_ips.tex}{% ============================================================================
% Fichier  : partie2/P02_C06_ids_ips.tex
% Titre    : La Détection d'Intrusion (IDS/IPS)
% Auteur   : MINKA MI NGUIDJOÏ Thierry Emmanuel
% ============================================================================
% Sources : NIST SP 800-94 (IDS/IPS) ; MITRE ATT&CK.
% ============================================================================

\chapter{La Détection d'Intrusion (IDS/IPS)}
\label{chap:idsips}
\chaptermark{La détection d'intrusion}

\epigraph{%
    « Détecter laisse passer pour voir~; prévenir bloque au risque de couper. »%
}{Principe du cours}

\niveauChapitre{Intermédiaire}{Chapitre~\ref{chap:siem}~; notions de réseaux.}

\begin{boxobjectifs}
\begin{itemize}
    \item Distinguer IDS et IPS, réseau et hôte.
    \item Comparer détection par signatures et par anomalies.
    \item Comprendre le placement des sondes.
    \item Relier les sondes au SIEM.
    \item Peser le risque d'un blocage automatique.
\end{itemize}
\end{boxobjectifs}

\bigskip

% ============================================================================
\section*{Question directrice : fallait-il bloquer automatiquement, au risque de couper la production de Garoua~?}
% ============================================================================

Un IPS aurait pu bloquer la connexion suspecte du 23~février. Mais un blocage
mal réglé coupe aussi des flux légitimes. Ce chapitre distingue les sondes,
leurs méthodes et leur placement, et pose l'arbitrage entre détecter et
bloquer.

\medskip
\textbf{Plan du chapitre.}
\begin{enumerate}
    \item IDS et IPS, réseau et hôte (section~\ref{sec:c06-types}).
    \item Signatures ou anomalies (section~\ref{sec:c06-methodes}).
    \item Le placement des sondes (section~\ref{sec:c06-placement}).
    \item Détecter ou bloquer (section~\ref{sec:c06-arbitrage}).
\end{enumerate}

\separateur

% ============================================================================
\section{IDS et IPS, Réseau et Hôte}
\label{sec:c06-types}
% ============================================================================

\begin{tableaucours}{L{2.4cm}L{4.8cm}L{5.0cm}}{%
    \tableauHeader{Sonde} &
    \tableauHeader{Agit} &
    \tableauHeader{Observe}
}
IDS & Détecte et alerte, sans bloquer & Le trafic ou l'activité \\
IPS & Détecte et bloque en ligne & Le trafic ou l'activité \\
Réseau (NIDS/NIPS) & Sur un point de passage du réseau & Les flux \\
Hôte (HIDS/HIPS) & Sur une machine & Les processus, fichiers, journaux \\
\end{tableaucours}
\vspace{-0.8em}
\begin{center}{\small\itshape Tableau~6.1. Familles de sondes de détection}\end{center}

% ============================================================================
\section{Signatures ou Anomalies}
\label{sec:c06-methodes}
% ============================================================================

\begin{tableaucours}{L{3.2cm}L{4.6cm}L{4.4cm}}{%
    \tableauHeader{Méthode} &
    \tableauHeader{Principe} &
    \tableauHeader{Limite}
}
Signatures & Reconnaît des motifs connus & Aveugle à l'inconnu \\
Anomalies & Signale l'écart à un profil normal & Faux positifs, apprentissage délicat \\
\end{tableaucours}
\vspace{-0.8em}
\begin{center}{\small\itshape Tableau~6.2. Deux méthodes de détection}\end{center}

Les deux sont complémentaires~: les signatures attrapent le connu à faible
bruit, les anomalies ouvrent une chance sur l'inconnu, au prix de faux
positifs à régler.

% ============================================================================
\section{Le Placement des Sondes}
\label{sec:c06-placement}
% ============================================================================

Une sonde réseau ne voit que le trafic qui passe par elle~: on la place aux
points stratégiques (sortie Internet, cœur de réseau, segment serveurs). Une
sonde hôte voit tout de sa machine mais rien des autres. Les alertes des
sondes alimentent le SIEM, qui les corrèle avec les journaux
(Chapitre~\ref{chap:siem}).

% ============================================================================
\section{Détecter ou Bloquer}
\label{sec:c06-arbitrage}
% ============================================================================

\begin{boxalert}
\textbf{Erreur fréquente.} Mettre tous les IPS en blocage automatique « pour
être tranquille ».

\smallskip
\textbf{Règle~:} un blocage mal réglé coupe des flux légitimes et peut arrêter
la production. On bloque automatiquement le connu à très forte confiance~; on
détecte et on fait valider le reste par un analyste.
\end{boxalert}

\begin{boxscenario}{Opération BÉNOUÉ. Détecter d'abord}
\textbf{Constatation.} LiZbiZ installe une sonde réseau en sortie et des
sondes hôte sur les serveurs. Les sondes sont d'abord en détection, non en
blocage.

\medskip
\textbf{Analyse} (confiance moyenne). Sur un parc sans historique de réglage,
commencer en détection évite de couper la production~; le passage au blocage
se fait règle par règle, une fois le bruit maîtrisé. Prudence d'abord, au nom
de la disponibilité.
\end{boxscenario}

\begin{boxmetier}{Analyste SOC et administrateur réseau}
\textbf{Régler avant de bloquer.} Le passage d'un IPS en blocage se décide
sur mesure~: on observe en détection, on mesure les faux positifs, puis on
bloque les règles sûres. La disponibilité de la production est un critère de
sécurité.
\end{boxmetier}

\separateur

\begin{boxresume}
\textbf{Ce qu'il faut retenir de ce chapitre~:}
\begin{itemize}
    \item \textbf{IDS} alerte, \textbf{IPS} bloque~; \textbf{réseau} ou
          \textbf{hôte}.
    \item \textbf{Signatures} (le connu) et \textbf{anomalies} (l'inconnu)
          se complètent.
    \item \textbf{Le placement} décide de ce qu'une sonde voit.
    \item \textbf{Les sondes alimentent le SIEM}.
    \item \textbf{Bloquer} se fait avec prudence, règle par règle.
\end{itemize}
\end{boxresume}

\begin{boxinfo}
\textbf{Points de contrôle de l'élève.} Je sais~: distinguer IDS/IPS et
réseau/hôte~; comparer signatures et anomalies~; placer une sonde~; arbitrer
entre détecter et bloquer.
\end{boxinfo}

\bigskip

\begin{boxexo}[title={\ding{45}\quad Exercice~6.1. IDS ou IPS \niveaubase}]
Lequel bloque en ligne~? Lequel se contente d'alerter~?
\end{boxexo}

\begin{boxexo}[title={\ding{45}\quad Exercice~6.2. Réseau ou hôte \niveaubase}]
Quelle sonde voit les processus d'une machine~? Laquelle voit les flux~?
\end{boxexo}

\begin{boxexo}[title={\ding{45}\quad Exercice~6.3. Méthodes \niveaubase}]
Quelle méthode est aveugle à l'inconnu~? Laquelle produit des faux positifs~?
\end{boxexo}

\begin{boxexo}[title={\ding{45}\quad Exercice~6.4. Placement \niveaumoyen}]
Où placer une sonde réseau pour voir l'exfiltration vers l'Internet~?
\end{boxexo}

\begin{boxexo}[title={\ding{45}\quad Exercice~6.5. Complémentarité \niveaumoyen}]
Pourquoi combiner signatures et anomalies~?
\end{boxexo}

\begin{boxexo}[title={\ding{45}\quad Exercice~6.6. Blocage \niveaumoyen}]
Pourquoi ne pas tout mettre en blocage automatique chez LiZbiZ~?
\end{boxexo}

\begin{boxexo}[title={\ding{45}\quad Exercice~6.7. SIEM \niveaumoyen}]
Comment les alertes des sondes s'articulent-elles avec le SIEM~?
\end{boxexo}

\begin{boxexo}[title={\ding{45}\quad Exercice~6.8. Angle mort \niveaumoyen}]
Quel trafic une sonde réseau mal placée laisse-t-elle passer~?
\end{boxexo}

\begin{boxexo}[title={\ding{45}\quad Exercice~6.9. Plan de déploiement \niveauexpert}]
Proposez le plan de déploiement des sondes pour LiZbiZ (points, détection
puis blocage).
\end{boxexo}

\begin{boxexo}[title={\ding{45}\quad Exercice~6.10. Passage au blocage \niveauexpert}]
Décrivez la démarche pour faire passer une règle d'IPS de la détection au
blocage sans risque pour la production.
\end{boxexo}

\bigskip

\begin{boxinfo}
\textbf{Pour aller plus loin.}
\begin{itemize}
    \item NIST SP~800-61r2~\cite{nist80061r2}~; MITRE ATT\&CK~\cite{mitreattack}.
\end{itemize}
\end{boxinfo}

\begin{boxinfo}
\textbf{Transition.} Sondes et SIEM ne valent que par leurs règles~: le
Chapitre~\ref{chap:regles} apprend à en écrire et à les éprouver.
\end{boxinfo}
}{}
        \IfFileExists{partie2/P02_C07_regles_detection.tex}{% ============================================================================
% Fichier  : partie2/P02_C07_regles_detection.tex
% Titre    : Écrire et Éprouver une Règle de Détection
% Auteur   : MINKA MI NGUIDJOÏ Thierry Emmanuel
% ============================================================================
% Sources : MITRE ATT&CK ; pratiques d'ingénierie de détection.
% ============================================================================

\chapter{Écrire et Éprouver une Règle de Détection}
\label{chap:regles}
\chaptermark{Écrire une règle de détection}

\epigraph{%
    « Une bonne règle se juge à deux nombres~:\\
    ce qu'elle attrape, et ce qu'elle dérange pour rien. »%
}{Principe du cours}

\niveauChapitre{Intermédiaire à expert}{Chapitres~\ref{chap:siem}
et~\ref{chap:idsips}.}

\begin{boxobjectifs}
\begin{itemize}
    \item Formuler une règle de détection en français avant de la coder.
    \item Relier chaque règle à une technique ATT\&CK et à sa trace.
    \item Éprouver une règle~: vrais positifs, faux positifs, faux négatifs.
    \item Maintenir un catalogue de règles dans le temps.
    \item Rédiger la fiche d'une règle.
\end{itemize}
\end{boxobjectifs}

\bigskip

% ============================================================================
\section*{Question directrice : comment transformer « un compte de service anormal » en une alerte fiable~?}
% ============================================================================

L'idée d'une règle est simple~; la faire fonctionner sans noyer l'analyste
l'est moins. Ce chapitre apprend à écrire une règle, à l'éprouver et à la
maintenir, à partir des trois cas d'usage de l'opération BÉNOUÉ.

\medskip
\textbf{Plan du chapitre.}
\begin{enumerate}
    \item Énoncer avant de coder (section~\ref{sec:c07-enoncer}).
    \item Relier à une technique et à sa trace (section~\ref{sec:c07-attack}).
    \item Éprouver une règle (section~\ref{sec:c07-eprouver}).
    \item Maintenir le catalogue (section~\ref{sec:c07-maintenir}).
\end{enumerate}

\separateur

% ============================================================================
\section{Énoncer avant de Coder}
\label{sec:c07-enoncer}
% ============================================================================

\begin{boxprocedure}
\textbf{Énoncer une règle}
\begin{enumerate}
    \item \textbf{Quoi}~: le comportement à détecter, en une phrase.
    \item \textbf{Source}~: le ou les journaux concernés.
    \item \textbf{Condition}~: ce qui déclenche (seuil, enchaînement, contexte).
    \item \textbf{Exceptions}~: les cas légitimes à exclure.
    \item \textbf{Réponse}~: priorité, action attendue de l'analyste.
\end{enumerate}
\end{boxprocedure}

Exemple~: « Alerter si un compte de service ouvre une session interactive sur
un poste de travail (source~: journaux du domaine, événement de connexion de
type~2), hors des comptes et plages de maintenance déclarés. »

% ============================================================================
\section{Relier à une Technique et à sa Trace}
\label{sec:c07-attack}
% ============================================================================

\begin{tableaucours}{L{4.4cm}L{4.2cm}L{3.6cm}}{%
    \tableauHeader{Cas d'usage (BÉNOUÉ)} &
    \tableauHeader{Technique ATT\&CK} &
    \tableauHeader{Trace}
}
Compte de service sur un poste & Identifiants valides (usage) & Connexion du domaine \\
Création de compte hors maintenance & Persistance & Événement de création de compte \\
Service installé à distance & Mouvement latéral & Événement d'installation de service \\
\end{tableaucours}
\vspace{-0.8em}
\begin{center}{\small\itshape Tableau~7.1. Trois règles reliées à ATT\&CK~\cite{mitreattack}}\end{center}

% ============================================================================
\section{Éprouver une Règle}
\label{sec:c07-eprouver}
% ============================================================================

\begin{boxdef}
\termecle{Vrai positif}~: la règle alerte sur un vrai comportement visé.
\termecle{Faux positif}~: elle alerte à tort. \termecle{Faux négatif}~: elle
laisse passer le comportement visé. On éprouve une règle en la rejouant sur
des journaux réels et sur un jeu de test (la technique jouée en laboratoire).
\end{boxdef}

\begin{boxCRO}
\textbf{Analyse CRO d'une règle trop large.} Une règle qui collecte et
inspecte un trafic très étendu « pour ne rien rater » dégrade la
\cro{-}{+}{0}~: elle augmente la détection (fiabilité) mais expose davantage
de données (confidentialité). On cible la règle sur ce qui est nécessaire.
\end{boxCRO}

% ============================================================================
\section{Maintenir le Catalogue}
\label{sec:c07-maintenir}
% ============================================================================

\begin{boxscenario}{Opération BÉNOUÉ. Du cas d'usage à la règle éprouvée}
\textbf{Constatation.} La règle « compte de service sur un poste » alerte
24~fois en une semaine de test~: 3~vrais positifs (comptes réellement mal
utilisés), 21~faux positifs (un script de sauvegarde légitime).

\medskip
\textbf{Analyse} (confiance élevée). En excluant le compte du script de
sauvegarde déclaré, les faux positifs tombent à zéro sans perdre les vrais
positifs. Une règle n'est utile qu'une fois éprouvée et affinée.
\end{boxscenario}

\begin{boxmetier}{Analyste SOC / ingénierie de détection}
\textbf{Une règle vit.} L'ingénieur de détection versionne ses règles, note
leur technique ATT\&CK, mesure leur bruit et les révise après chaque
incident. Un catalogue de règles est un actif à entretenir, pas une
installation unique.
\end{boxmetier}

\separateur

\begin{boxresume}
\textbf{Ce qu'il faut retenir de ce chapitre~:}
\begin{itemize}
    \item \textbf{Énoncer} en français~: quoi, source, condition, exceptions,
          réponse.
    \item \textbf{Relier} chaque règle à une technique ATT\&CK et à sa trace.
    \item \textbf{Éprouver}~: vrais positifs, faux positifs, faux négatifs.
    \item \textbf{Affiner} par les exceptions, sans perdre les vrais positifs.
    \item \textbf{Maintenir} le catalogue dans le temps.
\end{itemize}
\end{boxresume}

\begin{boxinfo}
\textbf{Points de contrôle de l'élève.} Je sais~: énoncer une règle~; la
relier à ATT\&CK~; distinguer les trois cas (VP, FP, FN)~; affiner par
exceptions~; entretenir un catalogue.
\end{boxinfo}

\bigskip

\begin{boxexo}[title={\ding{45}\quad Exercice~7.1. Énoncer \niveaubase}]
Énoncez en français la règle « service installé à distance hors maintenance ».
\end{boxexo}

\begin{boxexo}[title={\ding{45}\quad Exercice~7.2. Trois cas \niveaubase}]
Définissez vrai positif, faux positif, faux négatif.
\end{boxexo}

\begin{boxexo}[title={\ding{45}\quad Exercice~7.3. Technique \niveaubase}]
À quelle tactique ATT\&CK relier « création de compte hors maintenance »~?
\end{boxexo}

\begin{boxexo}[title={\ding{45}\quad Exercice~7.4. Affiner \niveaumoyen}]
Comment ramener à zéro les 21~faux positifs sans perdre les 3~vrais positifs~?
\end{boxexo}

\begin{boxexo}[title={\ding{45}\quad Exercice~7.5. Faux négatif \niveaumoyen}]
Pourquoi un faux négatif est-il plus grave qu'un faux positif~?
\end{boxexo}

\begin{boxexo}[title={\ding{45}\quad Exercice~7.6. CRO \niveaumoyen}]
Expliquez l'analyse CRO d'une règle trop large.
\end{boxexo}

\begin{boxexo}[title={\ding{45}\quad Exercice~7.7. Jeu de test \niveaumoyen}]
Comment constituer un jeu de test pour éprouver une règle~?
\end{boxexo}

\begin{boxexo}[title={\ding{45}\quad Exercice~7.8. Fiche de règle \niveaumoyen}]
Rédigez la fiche complète d'une des trois règles de BÉNOUÉ.
\end{boxexo}

\begin{boxexo}[title={\ding{45}\quad Exercice~7.9. Catalogue \niveauexpert}]
Organisez un catalogue de règles~: versionnage, métadonnées, révision.
\end{boxexo}

\begin{boxexo}[title={\ding{45}\quad Exercice~7.10. Couverture \niveauexpert}]
À partir d'ATT\&CK, proposez cinq règles qui augmenteraient la couverture de
détection de LiZbiZ.
\end{boxexo}

\bigskip

\begin{boxinfo}
\textbf{Pour aller plus loin.}
\begin{itemize}
    \item MITRE ATT\&CK~\cite{mitreattack}~; NIST SP~800-61r2~\cite{nist80061r2}.
\end{itemize}
\end{boxinfo}

\begin{boxinfo}
\textbf{Transition.} La Partie~III pose le cadre normatif et juridique de
tout ce qui précède~: le Chapitre~\ref{chap:normes} ouvre sur les normes de
l'incident et de la preuve.
\end{boxinfo}
}{}

    \part{Cadre Normatif et Juridique}
    \label{part:cadre}

        \IfFileExists{partie3/P03_C08_normes_incident_preuve.tex}{% ============================================================================
% Fichier  : partie3/P03_C08_normes_incident_preuve.tex
% Titre    : Normes de la Gestion d'Incident et de la Preuve
% Auteur   : MINKA MI NGUIDJOÏ Thierry Emmanuel
% ============================================================================
% Sources : ISO/IEC 27035 ; ISO/IEC 27037 ; NIST SP 800-86 ; RFC 3227.
% ============================================================================

\chapter{Normes de la Gestion d'Incident et de la Preuve}
\label{chap:normes}
\chaptermark{Normes de l'incident et de la preuve}

\epigraph{%
    « Une norme ne remplace pas le jugement~;\\
    elle en garantit la méthode et la trace. »%
}{Principe du cours}

\niveauChapitre{Fondamental}{Partie~II~; notions de gestion de la sécurité.}

\begin{boxobjectifs}
\begin{itemize}
    \item Décrire le cycle de gestion d'incident de l'ISO/IEC 27035.
    \item Connaître les principes de collecte de preuve de l'ISO/IEC 27037.
    \item Situer le guide NIST SP 800-86 et la RFC 3227.
    \item Articuler ces référentiels entre eux.
    \item Relier chaque norme à une pratique du cours.
\end{itemize}
\end{boxobjectifs}

\bigskip

% ============================================================================
\section*{Question directrice : quelles normes encadrent ce que LiZbiZ doit faire, du constat à la preuve~?}
% ============================================================================

Face à l'incident, LiZbiZ et NYONG SÉCURITÉ ne partent pas d'une page
blanche~: des normes internationales fixent le cycle de gestion et les
exigences de la preuve. Ce chapitre les présente et montre comment elles
s'articulent.

\medskip
\textbf{Plan du chapitre.}
\begin{enumerate}
    \item Le cycle ISO/IEC 27035 (section~\ref{sec:c08-27035}).
    \item La preuve selon l'ISO/IEC 27037 (section~\ref{sec:c08-27037}).
    \item NIST SP 800-86 et RFC 3227 (section~\ref{sec:c08-nist}).
    \item Articuler les référentiels (section~\ref{sec:c08-articuler}).
\end{enumerate}

\separateur

% ============================================================================
\section{Le Cycle ISO/IEC 27035}
\label{sec:c08-27035}
% ============================================================================

\begin{boxdef}
\termecle{ISO/IEC 27035}~: norme de gestion des incidents de sécurité de
l'information. Elle décrit un cycle en cinq phases~: \emph{planifier et
préparer}~; \emph{détecter et signaler}~; \emph{évaluer et décider}~;
\emph{répondre}~; \emph{tirer les leçons}.
\end{boxdef}

\begin{tableaucours}{L{4.0cm}L{8.2cm}}{%
    \tableauHeader{Phase (ISO 27035)} &
    \tableauHeader{Ce qu'elle recouvre}
}
Planifier et préparer & Politique, rôles, outils, playbooks \\
Détecter et signaler & Supervision, alertes, point de contact \\
Évaluer et décider & Qualification, gravité, décision de réponse \\
Répondre & Confinement, éradication, rétablissement \\
Tirer les leçons & Retour d'expérience, amélioration \\
\end{tableaucours}
\vspace{-0.8em}
\begin{center}{\small\itshape Tableau~8.1. Le cycle en cinq phases de l'ISO/IEC 27035~\cite{iso27035}}\end{center}

% ============================================================================
\section{La Preuve selon l'ISO/IEC 27037}
\label{sec:c08-27037}
% ============================================================================

\begin{boxdef}
\termecle{ISO/IEC 27037}~: lignes directrices pour l'identification, la
collecte, l'acquisition et la préservation des preuves numériques. Elle
repose sur trois exigences~: \termecle{pertinence}, \termecle{fiabilité} et
\termecle{suffisance}, garanties par la traçabilité des actes.
\end{boxdef}

\begin{boxprocedure}
\textbf{Préserver une preuve (ISO/IEC 27037)}
\begin{enumerate}
    \item Identifier le support et sa volatilité.
    \item Acquérir selon l'ordre de volatilité (du plus volatil au plus
          stable).
    \item Calculer un condensat (empreinte) à l'acquisition.
    \item Documenter chaque acte~: qui, quoi, quand, comment.
    \item Conserver en garantissant l'intégrité (chaîne de custody,
          Chapitre~\ref{chap:fondements}).
\end{enumerate}
\end{boxprocedure}

% ============================================================================
\section{NIST SP 800-86 et RFC 3227}
\label{sec:c08-nist}
% ============================================================================

Le \termecle{NIST SP 800-86} intègre la forensique à la réponse à incident en
quatre temps~: \emph{collecte}, \emph{examen}, \emph{analyse},
\emph{rapport}. La \termecle{RFC 3227} donne les lignes directrices de
collecte et d'archivage des preuves, dont le fameux \termecle{ordre de
volatilité} (mémoire avant disque, disque avant archives).

\begin{boxCRO}
\textbf{Analyse CRO de l'ordre de volatilité.} Capturer la mémoire d'abord
sert la \cro{0}{+}{+}~: on fige une trace fugace (fiabilité) exploitable
devant un tiers (opposabilité), sans exposer plus que nécessaire. Attendre,
c'est perdre la trace et son opposabilité.
\end{boxCRO}

% ============================================================================
\section{Articuler les Référentiels}
\label{sec:c08-articuler}
% ============================================================================

\begin{boxscenario}{Opération FARO. Quelle norme pour quel acte}
\textbf{Constatation.} Pour l'enquête, M.~BIKORO suit l'ISO 27035 (cycle),
l'ISO 27037 (préservation), le NIST 800-86 (démarche forensique) et la
RFC 3227 (ordre de volatilité).

\medskip
\textbf{Analyse} (confiance élevée). Les référentiels ne se concurrencent
pas~: 27035 cadre la gestion, 27037 et RFC 3227 encadrent la preuve, 800-86
relie forensique et réponse. Ensemble, ils garantissent une méthode
opposable.
\end{boxscenario}

\begin{boxmetier}{Analyste forensique et RSSI}
\textbf{Citer la norme suivie.} Dans un rapport, l'analyste indique quelle
norme a guidé chaque acte~: c'est ce qui rend la démarche vérifiable et
opposable. « J'ai acquis selon l'ISO 27037 » vaut mieux que « j'ai copié le
disque ».
\end{boxmetier}

\separateur

\begin{boxresume}
\textbf{Ce qu'il faut retenir de ce chapitre~:}
\begin{itemize}
    \item \textbf{ISO/IEC 27035}~: cycle d'incident en cinq phases.
    \item \textbf{ISO/IEC 27037}~: pertinence, fiabilité, suffisance de la
          preuve.
    \item \textbf{NIST 800-86}~: collecte, examen, analyse, rapport.
    \item \textbf{RFC 3227}~: ordre de volatilité.
    \item \textbf{Les référentiels s'articulent}~; on cite celui que l'on
          suit.
\end{itemize}
\end{boxresume}

\begin{boxinfo}
\textbf{Points de contrôle de l'élève.} Je sais~: énoncer les cinq phases de
27035~; citer les exigences de 27037~; décrire la démarche 800-86~; énoncer
l'ordre de volatilité~; articuler les référentiels.
\end{boxinfo}

\bigskip

\begin{boxexo}[title={\ding{45}\quad Exercice~8.1. Cinq phases \niveaubase}]
Citez les cinq phases du cycle ISO/IEC 27035 dans l'ordre.
\end{boxexo}

\begin{boxexo}[title={\ding{45}\quad Exercice~8.2. Exigences \niveaubase}]
Quelles sont les trois exigences de la preuve selon l'ISO/IEC 27037~?
\end{boxexo}

\begin{boxexo}[title={\ding{45}\quad Exercice~8.3. Volatilité \niveaubase}]
Dans quel ordre acquérir~: disque, mémoire, archives~?
\end{boxexo}

\begin{boxexo}[title={\ding{45}\quad Exercice~8.4. 800-86 \niveaumoyen}]
Décrivez les quatre temps du NIST SP 800-86 sur l'exemple de FARO.
\end{boxexo}

\begin{boxexo}[title={\ding{45}\quad Exercice~8.5. CRO \niveaumoyen}]
Expliquez l'analyse CRO de l'ordre de volatilité.
\end{boxexo}

\begin{boxexo}[title={\ding{45}\quad Exercice~8.6. Articulation \niveaumoyen}]
Quelle norme cadre la gestion~? Laquelle encadre la preuve~?
\end{boxexo}

\begin{boxexo}[title={\ding{45}\quad Exercice~8.7. Condensat \niveaumoyen}]
Pourquoi calculer un condensat dès l'acquisition~?
\end{boxexo}

\begin{boxexo}[title={\ding{45}\quad Exercice~8.8. Citer la norme \niveaumoyen}]
Pourquoi « acquis selon l'ISO 27037 » vaut-il mieux que « copié le disque »~?
\end{boxexo}

\begin{boxexo}[title={\ding{45}\quad Exercice~8.9. Démarche \niveauexpert}]
Rédigez, pour FARO, la démarche d'enquête en reliant chaque acte à sa norme.
\end{boxexo}

\begin{boxexo}[title={\ding{45}\quad Exercice~8.10. Playbook \niveauexpert}]
Intégrez l'ISO 27035 dans un playbook de réponse au rançongiciel de LiZbiZ.
\end{boxexo}

\bigskip

\begin{boxinfo}
\textbf{Pour aller plus loin.}
\begin{itemize}
    \item ISO/IEC 27035~\cite{iso27035}~; ISO/IEC 27037~\cite{iso27037}~;
          NIST SP~800-86~\cite{nist80086}~; RFC~3227~\cite{rfc3227}.
\end{itemize}
\end{boxinfo}

\begin{boxinfo}
\textbf{Transition.} Aux normes s'ajoute le droit~: le
Chapitre~\ref{chap:droit} traite du cadre juridique camerounais et africain.
\end{boxinfo}
}{}
        \IfFileExists{partie3/P03_C09_cadre_juridique.tex}{% ============================================================================
% Fichier  : partie3/P03_C09_cadre_juridique.tex
% Titre    : Le Cadre Juridique
% Auteur   : MINKA MI NGUIDJOÏ Thierry Emmanuel
% ============================================================================
% Sources : loi 2010/012, loi 2010/013, loi 2024/017 (Cameroun) ; Convention
% de Malabo (UA, 2014) ; Convention de Budapest (2001).
% ============================================================================

\chapter{Le Cadre Juridique}
\label{chap:droit}
\chaptermark{Le cadre juridique}

\epigraph{%
    « Les journaux tracent la frontière entre la faute professionnelle\\
    et la bonne foi prouvée. »%
}{Principe du cours}

\niveauChapitre{Fondamental}{Chapitre~\ref{chap:normes}~; notions de droit
numérique.}

\begin{boxobjectifs}
\begin{itemize}
    \item Connaître la loi camerounaise 2010/012 et la loi 2010/013.
    \item Situer la loi de 2024 sur les données personnelles.
    \item Comprendre le rôle et les pouvoirs de l'ANTIC.
    \item Articuler le droit national avec la Convention de Malabo et celle
          de Budapest.
    \item Déduire les obligations de LiZbiZ en cas d'incident.
\end{itemize}
\end{boxobjectifs}

\bigskip

% ============================================================================
\section*{Question directrice : qu'est-ce que la loi impose à LiZbiZ le 6 mars au matin~?}
% ============================================================================

L'incident n'est pas qu'un problème technique~: il déclenche des obligations.
Conserver les traces, signaler à l'ANTIC, notifier une fuite de données~:
autant de devoirs dont le non-respect engage la responsabilité. Ce chapitre
pose le cadre.

\medskip
\textbf{Plan du chapitre.}
\begin{enumerate}
    \item La loi 2010/012 et la loi 2010/013 (section~\ref{sec:c09-lois}).
    \item Les données personnelles (section~\ref{sec:c09-donnees}).
    \item L'ANTIC (section~\ref{sec:c09-antic}).
    \item Malabo et Budapest (section~\ref{sec:c09-international}).
    \item Les obligations de LiZbiZ (section~\ref{sec:c09-obligations}).
\end{enumerate}

\separateur

% ============================================================================
\section{La Loi 2010/012 et la Loi 2010/013}
\label{sec:c09-lois}
% ============================================================================

\begin{boxjuris}
La \loicam{loi n°~2010/012 du 21~décembre~2010} relative à la cybersécurité et
à la cybercriminalité est le texte fondateur~: elle réprime l'accès et le
maintien frauduleux dans un système, impose la conservation des traces et la
coopération avec l'ANTIC. La \loicam{loi n°~2010/013} régit les
communications électroniques et les obligations des opérateurs. Les articles
précis se lisent dans le texte publié au Journal officiel~; le cours n'en
cite que des éléments vérifiés.
\end{boxjuris}

% ============================================================================
\section{Les Données Personnelles}
\label{sec:c09-donnees}
% ============================================================================

\begin{boxjuris}
La \loicam{loi n°~2024/017 du 23~décembre~2024} régit la protection des
données à caractère personnel. Une fuite de données clients, comme celle
redoutée dans l'opération GAROUA, relève de ce texte~: information de
l'autorité, des personnes concernées, et documentation de la violation.
\end{boxjuris}

% ============================================================================
\section{L'ANTIC}
\label{sec:c09-antic}
% ============================================================================

\begin{boxdef}
\termecle{ANTIC} (Agence Nationale des Technologies de l'Information et de la
Communication)~: autorité camerounaise de la sécurité électronique. Elle
reçoit les signalements d'incidents, peut auditer, coordonne la réponse
nationale et coopère avec le parquet.
\end{boxdef}

% ============================================================================
\section{Malabo et Budapest}
\label{sec:c09-international}
% ============================================================================

\begin{tableaucours}{L{3.6cm}L{8.6cm}}{%
    \tableauHeader{Instrument} &
    \tableauHeader{Portée}
}
Convention de Malabo (UA, 2014) & Cadre continental~: cybersécurité, cybercriminalité, données personnelles \\
Convention de Budapest (2001) & Cadre international~: infractions, preuve électronique, coopération \\
\end{tableaucours}
\vspace{-0.8em}
\begin{center}{\small\itshape Tableau~9.1. Deux instruments supranationaux~\cite{malabo2014,budapest2001}}\end{center}

% ============================================================================
\section{Les Obligations de LiZbiZ}
\label{sec:c09-obligations}
% ============================================================================

\begin{boxscenario}{Opération GAROUA. Les devoirs du 6 mars}
\textbf{Constatation.} LiZbiZ subit un rançongiciel avec exfiltration
probable de données clients. Elle doit~: conserver les journaux, signaler
l'incident à l'ANTIC, documenter la violation, et, si des données
personnelles sont touchées, en informer l'autorité et les personnes.

\medskip
\textbf{Analyse} (confiance élevée). Ne pas disposer de journaux exploitables
affaiblirait la position de LiZbiZ~: la preuve de sa diligence repose sur les
traces. Le droit transforme la supervision en obligation, et le forensique en
moyen de preuve.
\end{boxscenario}

\begin{boxmetier}{RSSI et juriste TIC}
\textbf{Connaître ses obligations avant l'incident.} Le RSSI prépare, à
froid, la liste des signalements et notifications dus, avec leurs
destinataires et délais. En pleine crise, on applique une procédure, on ne
découvre pas le droit.
\end{boxmetier}

\separateur

\begin{boxresume}
\textbf{Ce qu'il faut retenir de ce chapitre~:}
\begin{itemize}
    \item \textbf{Loi 2010/012}~: accès frauduleux, conservation, coopération.
    \item \textbf{Loi 2024/017}~: protection des données personnelles.
    \item \textbf{L'ANTIC} reçoit les signalements et coordonne.
    \item \textbf{Malabo et Budapest} encadrent au niveau continental et
          international.
    \item \textbf{Les journaux} fondent la preuve de la diligence.
\end{itemize}
\end{boxresume}

\begin{boxinfo}
\textbf{Points de contrôle de l'élève.} Je sais~: citer les lois 2010/012,
2010/013, 2024/017~; décrire l'ANTIC~; situer Malabo et Budapest~; énoncer
les obligations de LiZbiZ.
\end{boxinfo}

\bigskip

\begin{boxexo}[title={\ding{45}\quad Exercice~9.1. La loi fondatrice \niveaubase}]
Quel texte camerounais fonde le droit de la cybersécurité~? De quelle année~?
\end{boxexo}

\begin{boxexo}[title={\ding{45}\quad Exercice~9.2. L'ANTIC \niveaubase}]
Citez trois pouvoirs de l'ANTIC.
\end{boxexo}

\begin{boxexo}[title={\ding{45}\quad Exercice~9.3. Données \niveaubase}]
Quel texte protège les données personnelles au Cameroun~?
\end{boxexo}

\begin{boxexo}[title={\ding{45}\quad Exercice~9.4. Malabo \niveaumoyen}]
Qu'apporte la Convention de Malabo par rapport au droit national~?
\end{boxexo}

\begin{boxexo}[title={\ding{45}\quad Exercice~9.5. Preuve de diligence \niveaumoyen}]
Pourquoi les journaux tracent-ils la frontière entre faute et bonne foi~?
\end{boxexo}

\begin{boxexo}[title={\ding{45}\quad Exercice~9.6. Notifications \niveaumoyen}]
Quelles notifications LiZbiZ doit-elle envisager le 6~mars~?
\end{boxexo}

\begin{boxexo}[title={\ding{45}\quad Exercice~9.7. Budapest \niveaumoyen}]
En quoi la Convention de Budapest facilite-t-elle la coopération
internationale~?
\end{boxexo}

\begin{boxexo}[title={\ding{45}\quad Exercice~9.8. À froid \niveaumoyen}]
Pourquoi préparer la liste des obligations avant l'incident~?
\end{boxexo}

\begin{boxexo}[title={\ding{45}\quad Exercice~9.9. Procédure de signalement \niveauexpert}]
Rédigez la procédure de signalement à l'ANTIC pour LiZbiZ~: déclencheur,
contenu, délai, destinataire.
\end{boxexo}

\begin{boxexo}[title={\ding{45}\quad Exercice~9.10. Note juridique \niveauexpert}]
Rédigez une note qui récapitule les obligations légales de LiZbiZ en cas
d'incident majeur.
\end{boxexo}

\bigskip

\begin{boxinfo}
\textbf{Pour aller plus loin.}
\begin{itemize}
    \item Lois 2010/012~\cite{loi2010012}, 2010/013~\cite{loi2010013},
          2024/017~\cite{loi2024017}.
    \item Convention de Malabo~\cite{malabo2014}~; Convention de
          Budapest~\cite{budapest2001}.
\end{itemize}
\end{boxinfo}

\begin{boxinfo}
\textbf{Transition.} Normes et droit convergent sur une question~: une trace
est-elle recevable~? C'est l'objet du Chapitre~\ref{chap:preuve}.
\end{boxinfo}
}{}
        \IfFileExists{partie3/P03_C10_admissibilite_preuve.tex}{% ============================================================================
% Fichier  : partie3/P03_C10_admissibilite_preuve.tex
% Titre    : L'Admissibilité de la Preuve Numérique
% Auteur   : MINKA MI NGUIDJOÏ Thierry Emmanuel
% ============================================================================
% Sources : ISO/IEC 27037 ; RFC 3227 ; loi 2010/012. Opération FARO : fictif.
% ============================================================================

\chapter{L'Admissibilité de la Preuve Numérique}
\label{chap:preuve}
\chaptermark{L'admissibilité de la preuve}

\epigraph{%
    « Une preuve que l'on ne peut pas défendre\\
    ne vaut pas mieux qu'une preuve que l'on n'a pas. »%
}{Principe du cours}

\niveauChapitre{Intermédiaire}{Chapitres~\ref{chap:normes} et~\ref{chap:droit}.}

\begin{boxobjectifs}
\begin{itemize}
    \item Définir l'admissibilité et ses conditions.
    \item Comprendre l'intégrité, l'authenticité et la traçabilité.
    \item Relier l'admissibilité au Trilemme CRO.
    \item Reconnaître ce qui détruit une preuve.
    \item Préparer une preuve opposable dès l'acquisition.
\end{itemize}
\end{boxobjectifs}

\bigskip

% ============================================================================
\section*{Question directrice : la preuve recueillie chez LiZbiZ tiendrait-elle devant un tiers~?}
% ============================================================================

Recueillir une trace est une chose~; pouvoir s'en prévaloir devant la
direction, l'ANTIC ou un juge en est une autre. Ce chapitre pose les
conditions de l'admissibilité et montre comment une manipulation imprudente
détruit une preuve pourtant exacte.

\medskip
\textbf{Plan du chapitre.}
\begin{enumerate}
    \item Qu'est-ce que l'admissibilité (section~\ref{sec:c10-definir}).
    \item Intégrité, authenticité, traçabilité (section~\ref{sec:c10-conditions}).
    \item Le Trilemme CRO appliqué (section~\ref{sec:c10-cro}).
    \item Ce qui détruit une preuve (section~\ref{sec:c10-detruire}).
\end{enumerate}

\separateur

% ============================================================================
\section{Qu'est-ce que l'Admissibilité}
\label{sec:c10-definir}
% ============================================================================

\begin{boxdef}
\termecle{Admissibilité}~: aptitude d'une preuve à être reçue et prise en
compte par un tiers (direction, autorité, juge). Elle ne tient pas à la seule
exactitude de la trace, mais à la \emph{manière} dont elle a été recueillie,
conservée et documentée.
\end{boxdef}

% ============================================================================
\section{Intégrité, Authenticité, Traçabilité}
\label{sec:c10-conditions}
% ============================================================================

\begin{tableaucours}{L{3.2cm}L{9.0cm}}{%
    \tableauHeader{Condition} &
    \tableauHeader{Ce qu'elle exige}
}
Intégrité & La preuve n'a pas été altérée depuis l'acquisition (condensat). \\
Authenticité & La preuve est bien ce qu'elle prétend être, de la source déclarée. \\
Traçabilité & Chaque manipulation est documentée (chaîne de custody). \\
Licéité & La collecte respecte le droit (autorisation, proportionnalité). \\
\end{tableaucours}
\vspace{-0.8em}
\begin{center}{\small\itshape Tableau~10.1. Conditions de l'admissibilité}\end{center}

% ============================================================================
\section{Le Trilemme CRO Appliqué}
\label{sec:c10-cro}
% ============================================================================

\begin{boxCRO}
\textbf{L'admissibilité est l'opposabilité.} La dimension O du Trilemme est
précisément l'admissibilité~: une preuve opposable est une preuve admissible.
La F (fiabilité) recouvre l'intégrité et l'authenticité~; la C
(confidentialité) impose de ne saisir que le nécessaire. Une bonne pratique
vise \cro{0}{+}{+} à chaque acte.
\end{boxCRO}

% ============================================================================
\section{Ce qui Détruit une Preuve}
\label{sec:c10-detruire}
% ============================================================================

\begin{boxalert}
\textbf{Violation typique.} Travailler directement sur l'original, sans
condensat ni copie.

\smallskip
\textbf{Règle~:} on travaille sur une copie vérifiée par condensat, on
documente chaque acte, on ne modifie jamais l'original. Une preuve manipulée
sans trace est perdue, même si son contenu est exact.
\end{boxalert}

\begin{boxscenario}{Opération FARO. La preuve sauvée de justesse}
\textbf{Constatation.} Un technicien de LiZbiZ, de bonne foi, avait
redémarré le serveur \texttt{FS-01} et commencé à « nettoyer » avant
l'arrivée de M.~BIKORO. La mémoire vive était perdue~; le disque, lui, n'avait
pas encore été écrasé.

\medskip
\textbf{Analyse} (confiance élevée). Le redémarrage a détruit une preuve
volatile (mémoire) et a failli compromettre le disque. M.~BIKORO a figé
l'état restant, calculé les condensats et documenté la chaîne~: le disque
reste opposable, la mémoire est perdue. La leçon~: ne rien toucher avant
d'avoir figé.
\end{boxscenario}

\begin{boxmetier}{Analyste forensique}
\textbf{Figer avant de comprendre.} L'analyste résiste à la tentation
d'« aller voir » sur le système vivant~: chaque action non documentée dégrade
l'opposabilité. On fige, on copie, on vérifie, puis on analyse la copie.
\end{boxmetier}

\separateur

\begin{boxresume}
\textbf{Ce qu'il faut retenir de ce chapitre~:}
\begin{itemize}
    \item \textbf{L'admissibilité} tient à la manière, pas à la seule
          exactitude.
    \item \textbf{Intégrité, authenticité, traçabilité, licéité}~: les
          conditions.
    \item \textbf{L'opposabilité} du CRO \emph{est} l'admissibilité.
    \item \textbf{Travailler sur l'original} détruit la preuve.
    \item \textbf{Figer avant de comprendre.}
\end{itemize}
\end{boxresume}

\begin{boxinfo}
\textbf{Points de contrôle de l'élève.} Je sais~: définir l'admissibilité~;
citer ses conditions~; la relier au CRO~; reconnaître ce qui détruit une
preuve~; préparer une preuve opposable.
\end{boxinfo}

\bigskip

\begin{boxexo}[title={\ding{45}\quad Exercice~10.1. Définir \niveaubase}]
L'admissibilité tient-elle à l'exactitude de la trace ou à la manière de la
recueillir~?
\end{boxexo}

\begin{boxexo}[title={\ding{45}\quad Exercice~10.2. Conditions \niveaubase}]
Citez les quatre conditions de l'admissibilité.
\end{boxexo}

\begin{boxexo}[title={\ding{45}\quad Exercice~10.3. Condensat \niveaubase}]
À quoi sert le condensat (empreinte) d'une preuve~?
\end{boxexo}

\begin{boxexo}[title={\ding{45}\quad Exercice~10.4. CRO \niveaumoyen}]
Pourquoi dit-on que l'opposabilité du CRO est l'admissibilité~?
\end{boxexo}

\begin{boxexo}[title={\ding{45}\quad Exercice~10.5. FARO \niveaumoyen}]
Quelle preuve le redémarrage a-t-il détruite~? Laquelle a été sauvée~?
\end{boxexo}

\begin{boxexo}[title={\ding{45}\quad Exercice~10.6. L'original \niveaumoyen}]
Pourquoi ne jamais travailler sur l'original~?
\end{boxexo}

\begin{boxexo}[title={\ding{45}\quad Exercice~10.7. Bonne foi \niveaumoyen}]
Le technicien était de bonne foi. Pourquoi son geste était-il néanmoins
dommageable~?
\end{boxexo}

\begin{boxexo}[title={\ding{45}\quad Exercice~10.8. Licéité \niveaumoyen}]
En quoi une collecte illicite compromet-elle l'admissibilité~?
\end{boxexo}

\begin{boxexo}[title={\ding{45}\quad Exercice~10.9. Note de première intervention \niveauexpert}]
Rédigez la consigne de première intervention à afficher près des serveurs de
LiZbiZ (« que faire, que ne pas faire »).
\end{boxexo}

\begin{boxexo}[title={\ding{45}\quad Exercice~10.10. Dossier de preuve \niveauexpert}]
Décrivez le dossier qui accompagne une preuve pour la rendre opposable.
\end{boxexo}

\bigskip

\begin{boxinfo}
\textbf{Pour aller plus loin.}
\begin{itemize}
    \item ISO/IEC 27037~\cite{iso27037}~; RFC~3227~\cite{rfc3227}.
\end{itemize}
\end{boxinfo}

\begin{boxinfo}
\textbf{Transition.} La Partie~IV entre dans la pratique de l'enquête~: le
Chapitre~\ref{chap:fondements} pose les fondements du forensique et la chaîne
de custody.
\end{boxinfo}
}{}

    \part{L'Analyse Forensique}
    \label{part:forensique}

        \IfFileExists{partie4/P04_C11_fondements_custody.tex}{% ============================================================================
% Fichier  : partie4/P04_C11_fondements_custody.tex
% Titre    : Fondements du Forensique et Chaîne de Custody
% Auteur   : MINKA MI NGUIDJOÏ Thierry Emmanuel
% ============================================================================
% Sources : ISO/IEC 27037 ; RFC 3227 ; NIST SP 800-86. Opération FARO : fictif.
% ============================================================================

\chapter{Fondements du Forensique et Chaîne de Custody}
\label{chap:fondements}
\chaptermark{Fondements du forensique}

\epigraph{%
    « L'enquêteur ne cherche pas à avoir raison~;\\
    il cherche à pouvoir le prouver. »%
}{Principe du cours}

\niveauChapitre{Fondamental}{Partie~III.}

\begin{boxobjectifs}
\begin{itemize}
    \item Définir l'analyse forensique et ses principes.
    \item Comprendre l'ordre de volatilité et la copie vérifiée.
    \item Tenir une chaîne de custody.
    \item Qualifier un artefact.
    \item Préserver l'opposabilité à chaque étape.
\end{itemize}
\end{boxobjectifs}

\bigskip

% ============================================================================
\section*{Question directrice : par où commence une enquête propre sur FS-01~?}
% ============================================================================

Devant un serveur compromis, la tentation est d'« aller voir ». Le forensique
impose l'inverse~: figer, copier, documenter, puis analyser une copie. Ce
chapitre pose les fondements et la chaîne de custody.

\medskip
\textbf{Plan du chapitre.}
\begin{enumerate}
    \item Qu'est-ce que le forensique (section~\ref{sec:c11-definir}).
    \item L'ordre de volatilité (section~\ref{sec:c11-volatilite}).
    \item La copie vérifiée (section~\ref{sec:c11-copie}).
    \item La chaîne de custody (section~\ref{sec:c11-custody}).
    \item Qualifier un artefact (section~\ref{sec:c11-artefact}).
\end{enumerate}

\separateur

% ============================================================================
\section{Qu'est-ce que le Forensique}
\label{sec:c11-definir}
% ============================================================================

\begin{boxdef}
\termecle{Analyse forensique}~: recherche méthodique, documentée et opposable
des traces laissées par un événement sur les systèmes, afin d'établir ce qui
s'est passé, comment, et par qui, dans le respect de l'intégrité des preuves.
\end{boxdef}

% ============================================================================
\section{L'Ordre de Volatilité}
\label{sec:c11-volatilite}
% ============================================================================

\begin{boxprocedure}
\textbf{Acquérir selon l'ordre de volatilité (RFC 3227)}
\begin{enumerate}
    \item Mémoire vive et état du système (processus, connexions).
    \item Fichiers temporaires et données de session.
    \item Disque (systèmes de fichiers).
    \item Journaux distants et configuration.
    \item Supports d'archives.
\end{enumerate}
\end{boxprocedure}

On capture d'abord ce qui disparaît le plus vite. Un redémarrage détruit la
mémoire~: d'où la règle « ne rien toucher avant d'avoir figé »
(Chapitre~\ref{chap:preuve}).

% ============================================================================
\section{La Copie Vérifiée}
\label{sec:c11-copie}
% ============================================================================

On n'analyse jamais l'original~: on en fait une \termecle{copie bit à bit},
dont on calcule un \termecle{condensat} (empreinte). Le même condensat sur
l'original et la copie prouve que la copie est fidèle et que l'original n'a
pas changé. Toute l'analyse se fait sur la copie.

% ============================================================================
\section{La Chaîne de Custody}
\label{sec:c11-custody}
% ============================================================================

\begin{boxprocedure}
\textbf{Tenir la chaîne de custody}
\begin{enumerate}
    \item Identifier chaque pièce (numéro, description, condensat).
    \item Consigner chaque transfert~: qui, à qui, quand, pourquoi.
    \item Protéger physiquement et logiquement les pièces.
    \item Ne jamais laisser de trou dans la continuité de la détention.
    \item Joindre la chaîne au dossier de preuve.
\end{enumerate}
\end{boxprocedure}

\begin{boxCRO}
\textbf{La chaîne de custody sert l'opposabilité.} Une pièce dont la
détention présente un trou devient contestable~: \cro{0}{0}{-}. La continuité
documentée est ce qui rend la preuve opposable devant un tiers.
\end{boxCRO}

% ============================================================================
\section{Qualifier un Artefact}
\label{sec:c11-artefact}
% ============================================================================

\begin{boxobjet}{Journal de connexion de FS-01}
\textbf{Nature~:} journal d'authentification extrait de la copie du disque de
\texttt{FS-01}. \textbf{Source~:} copie vérifiée (condensat au dossier).
\textbf{Pertinence~:} contient les connexions du compte de service suspect.
\textbf{Fiabilité~:} intègre (condensat), horodatage à recaler (écart connu
de 40~min). \textbf{Opposabilité~:} chaîne de custody continue.
\end{boxobjet}

\begin{boxmetier}{Analyste forensique}
\textbf{Documenter est enquêter.} Pour l'analyste, la documentation n'est pas
une formalité postérieure~: c'est une partie de l'acte d'enquête. Un geste
non documenté est un geste perdu.
\end{boxmetier}

\separateur

\begin{boxresume}
\textbf{Ce qu'il faut retenir de ce chapitre~:}
\begin{itemize}
    \item \textbf{Le forensique} établit les faits de façon opposable.
    \item \textbf{L'ordre de volatilité}~: du plus fugace au plus stable.
    \item \textbf{La copie vérifiée} par condensat remplace l'original.
    \item \textbf{La chaîne de custody} assure la continuité de la détention.
    \item \textbf{Un artefact} se qualifie (nature, source, pertinence,
          fiabilité, opposabilité).
\end{itemize}
\end{boxresume}

\begin{boxinfo}
\textbf{Points de contrôle de l'élève.} Je sais~: définir le forensique~;
énoncer l'ordre de volatilité~; faire une copie vérifiée~; tenir une chaîne de
custody~; qualifier un artefact.
\end{boxinfo}

\bigskip

\begin{boxexo}[title={\ding{45}\quad Exercice~11.1. Volatilité \niveaubase}]
Classez du plus volatil au plus stable~: disque, mémoire vive, archives.
\end{boxexo}

\begin{boxexo}[title={\ding{45}\quad Exercice~11.2. Copie \niveaubase}]
Pourquoi analyse-t-on une copie et non l'original~?
\end{boxexo}

\begin{boxexo}[title={\ding{45}\quad Exercice~11.3. Custody \niveaubase}]
Qu'est-ce qu'un « trou » dans la chaîne de custody~?
\end{boxexo}

\begin{boxexo}[title={\ding{45}\quad Exercice~11.4. Condensat \niveaumoyen}]
Comment un condensat prouve-t-il qu'une copie est fidèle~?
\end{boxexo}

\begin{boxexo}[title={\ding{45}\quad Exercice~11.5. CRO \niveaumoyen}]
Pourquoi un trou de custody dégrade-t-il l'opposabilité~?
\end{boxexo}

\begin{boxexo}[title={\ding{45}\quad Exercice~11.6. Artefact \niveaumoyen}]
Qualifiez un artefact de votre choix de l'opération FARO (cinq rubriques).
\end{boxexo}

\begin{boxexo}[title={\ding{45}\quad Exercice~11.7. Horloge \niveaumoyen}]
Comment traiter l'écart d'horloge de 40~min dans l'analyse~?
\end{boxexo}

\begin{boxexo}[title={\ding{45}\quad Exercice~11.8. Documentation \niveaumoyen}]
Pourquoi un geste non documenté est-il « perdu »~?
\end{boxexo}

\begin{boxexo}[title={\ding{45}\quad Exercice~11.9. Procédure d'acquisition \niveauexpert}]
Rédigez la procédure d'acquisition du serveur \texttt{FS-01}, dans l'ordre de
volatilité.
\end{boxexo}

\begin{boxexo}[title={\ding{45}\quad Exercice~11.10. Formulaire de custody \niveauexpert}]
Concevez le formulaire de chaîne de custody de LiZbiZ.
\end{boxexo}

\bigskip

\begin{boxinfo}
\textbf{Pour aller plus loin.}
\begin{itemize}
    \item ISO/IEC 27037~\cite{iso27037}~; RFC~3227~\cite{rfc3227}~;
          NIST SP~800-86~\cite{nist80086}.
\end{itemize}
\end{boxinfo}

\begin{boxinfo}
\textbf{Transition.} Le premier terrain d'enquête est le disque~: le
Chapitre~\ref{chap:systeme} traite du forensique système.
\end{boxinfo}
}{}
        \IfFileExists{partie4/P04_C12_forensique_systeme.tex}{% ============================================================================
% Fichier  : partie4/P04_C12_forensique_systeme.tex
% Titre    : Forensique Système
% Auteur   : MINKA MI NGUIDJOÏ Thierry Emmanuel
% ============================================================================
% Sources : NIST SP 800-86 ; artefacts Windows et Linux. Opération FARO.
% ============================================================================

\chapter{Forensique Système}
\label{chap:systeme}
\chaptermark{Forensique système}

\epigraph{%
    « Le disque garde la mémoire de ce que l'utilisateur croyait avoir effacé. »%
}{Principe du cours}

\niveauChapitre{Intermédiaire à expert}{Chapitre~\ref{chap:fondements}~;
administration Windows et Linux.}

\begin{boxobjectifs}
\begin{itemize}
    \item Identifier les artefacts système utiles sous Windows et Linux.
    \item Reconstituer une chronologie à partir des horodatages.
    \item Repérer les traces de persistance et d'exécution.
    \item Distinguer suppression logique et effacement réel.
    \item Rédiger un constat système.
\end{itemize}
\end{boxobjectifs}

\bigskip

% ============================================================================
\section*{Question directrice : que révèle le disque de FS-01 sur les onze jours~?}
% ============================================================================

Le disque de \texttt{FS-01} contient les traces du passage de l'attaquant~:
comptes créés, services installés, programmes exécutés, fichiers consultés.
Ce chapitre apprend à les lire et à en tirer une chronologie opposable.

\medskip
\textbf{Plan du chapitre.}
\begin{enumerate}
    \item Les artefacts Windows (section~\ref{sec:c12-windows}).
    \item Les artefacts Linux (section~\ref{sec:c12-linux}).
    \item La chronologie (section~\ref{sec:c12-chrono}).
    \item Suppression et récupération (section~\ref{sec:c12-suppression}).
\end{enumerate}

\separateur

% ============================================================================
\section{Les Artefacts Windows}
\label{sec:c12-windows}
% ============================================================================

\begin{tableaucours}{L{4.0cm}L{8.2cm}}{%
    \tableauHeader{Artefact} &
    \tableauHeader{Ce qu'il révèle}
}
Journaux d'événements & Connexions, comptes, services, tâches \\
Base de registre & Programmes au démarrage, périphériques branchés \\
Traces d'exécution & Programmes lancés, fréquence, dates \\
Fichiers récents, raccourcis & Documents ouverts, chemins \\
Fichier d'échange, clichés & Fragments de mémoire, versions antérieures \\
\end{tableaucours}
\vspace{-0.8em}
\begin{center}{\small\itshape Tableau~12.1. Artefacts forensiques Windows}\end{center}

% ============================================================================
\section{Les Artefacts Linux}
\label{sec:c12-linux}
% ============================================================================

\begin{tableaucours}{L{4.0cm}L{8.2cm}}{%
    \tableauHeader{Artefact} &
    \tableauHeader{Ce qu'il révèle}
}
Journaux système & Authentifications, sudo, démons \\
Historique des interpréteurs & Commandes saisies \\
Tâches planifiées (cron) & Persistance, exécutions régulières \\
Horodatages des fichiers & Accès, modification, changement de métadonnées \\
Comptes et clés & Comptes ajoutés, clés d'accès déposées \\
\end{tableaucours}
\vspace{-0.8em}
\begin{center}{\small\itshape Tableau~12.2. Artefacts forensiques Linux}\end{center}

% ============================================================================
\section{La Chronologie}
\label{sec:c12-chrono}
% ============================================================================

\begin{boxprocedure}
\textbf{Construire une chronologie (\textit{timeline})}
\begin{enumerate}
    \item Rassembler les horodatages de tous les artefacts.
    \item Recaler les horloges (écarts connus).
    \item Ordonner les événements et repérer les enchaînements.
    \item Distinguer le certain (trace) du probable (interprétation).
    \item Dater chaque étape de l'attaque.
\end{enumerate}
\end{boxprocedure}

\begin{boxscenario}{Opération FARO. Onze jours reconstitués}
\textbf{Constatation.} La chronologie du disque de \texttt{FS-01} établit~:
23~février, connexion du compte \texttt{svc\_sauvegarde} à 2~h~14~; 24~février,
création du compte \texttt{svc\_maj2}~; 1\textsuperscript{er}~mars,
installation du service \texttt{WinUpd32}~; 4~mars, accès massif aux
partages~; 6~mars, chiffrement.

\medskip
\textbf{Analyse} (confiance élevée pour les traces, moyenne pour les
intentions). Les horodatages, recalés, dessinent une progression cohérente.
On distingue ce qui est tracé (les événements) de ce qu'on en déduit (le but
de l'attaquant).
\end{boxscenario}

% ============================================================================
\section{Suppression et Récupération}
\label{sec:c12-suppression}
% ============================================================================

Supprimer un fichier ne l'efface pas~: le système marque l'espace comme
réutilisable, mais les données subsistent jusqu'à réécriture. D'où la
possibilité de \termecle{récupérer} des fichiers supprimés, tant qu'ils n'ont
pas été écrasés. L'attaquant qui « nettoie » laisse souvent des traces
récupérables.

\begin{boxmetier}{Analyste forensique}
\textbf{Séparer le tracé du déduit.} Dans un constat système, l'analyste
écrit d'abord ce que la trace montre, puis, à part, ce qu'il en déduit, avec
un niveau de confiance. C'est le cœur de la méthode constatation/analyse.
\end{boxmetier}

\separateur

\begin{boxresume}
\textbf{Ce qu'il faut retenir de ce chapitre~:}
\begin{itemize}
    \item \textbf{Windows et Linux} laissent des artefacts riches.
    \item \textbf{La chronologie} relie les horodatages recalés.
    \item \textbf{On date} chaque étape de l'attaque.
    \item \textbf{Supprimer n'est pas effacer}~: récupération possible.
    \item \textbf{On sépare} le tracé du déduit.
\end{itemize}
\end{boxresume}

\begin{boxinfo}
\textbf{Points de contrôle de l'élève.} Je sais~: citer les artefacts Windows
et Linux~; construire une chronologie~; dater une attaque~; expliquer la
récupération~; séparer trace et interprétation.
\end{boxinfo}

\bigskip

\begin{boxexo}[title={\ding{45}\quad Exercice~12.1. Artefact Windows \niveaubase}]
Quel artefact Windows révèle les programmes lancés au démarrage~?
\end{boxexo}

\begin{boxexo}[title={\ding{45}\quad Exercice~12.2. Artefact Linux \niveaubase}]
Quel artefact Linux révèle les commandes saisies~?
\end{boxexo}

\begin{boxexo}[title={\ding{45}\quad Exercice~12.3. Suppression \niveaubase}]
Pourquoi un fichier supprimé peut-il être récupéré~?
\end{boxexo}

\begin{boxexo}[title={\ding{45}\quad Exercice~12.4. Chronologie \niveaumoyen}]
À partir des faits de FARO, dressez la chronologie des onze jours.
\end{boxexo}

\begin{boxexo}[title={\ding{45}\quad Exercice~12.5. Recaler \niveaumoyen}]
Comment l'écart d'horloge de 40~min modifie-t-il la chronologie~?
\end{boxexo}

\begin{boxexo}[title={\ding{45}\quad Exercice~12.6. Tracé ou déduit \niveaumoyen}]
Classez~: « le service WinUpd32 a été installé le 1\textsuperscript{er}~mars »~;
« l'attaquant cherchait la persistance ».
\end{boxexo}

\begin{boxexo}[title={\ding{45}\quad Exercice~12.7. Persistance \niveaumoyen}]
Quels artefacts trahissent une persistance sous Windows~? Sous Linux~?
\end{boxexo}

\begin{boxexo}[title={\ding{45}\quad Exercice~12.8. Nettoyage \niveaumoyen}]
Pourquoi le « nettoyage » de l'attaquant laisse-t-il souvent des traces~?
\end{boxexo}

\begin{boxexo}[title={\ding{45}\quad Exercice~12.9. Constat système \niveauexpert}]
Rédigez un constat système complet pour l'installation de \texttt{WinUpd32}.
\end{boxexo}

\begin{boxexo}[title={\ding{45}\quad Exercice~12.10. Plan d'analyse \niveauexpert}]
Proposez un plan d'analyse du disque de \texttt{FS-01}, des artefacts à la
chronologie.
\end{boxexo}

\bigskip

\begin{boxinfo}
\textbf{Pour aller plus loin.}
\begin{itemize}
    \item NIST SP~800-86~\cite{nist80086}~; ISO/IEC 27037~\cite{iso27037}.
\end{itemize}
\end{boxinfo}

\begin{boxinfo}
\textbf{Transition.} Le disque ne dit pas tout~: la mémoire vive garde ce qui
n'a jamais été écrit sur disque. Le Chapitre~\ref{chap:memoire} traite du
forensique mémoire.
\end{boxinfo}
}{}
        \IfFileExists{partie4/P04_C13_forensique_memoire.tex}{% ============================================================================
% Fichier  : partie4/P04_C13_forensique_memoire.tex
% Titre    : Forensique Mémoire
% Auteur   : MINKA MI NGUIDJOÏ Thierry Emmanuel
% ============================================================================
% Sources : NIST SP 800-86 ; RFC 3227 ; MITRE ATT&CK. Opération FARO.
% ============================================================================

\chapter{Forensique Mémoire}
\label{chap:memoire}
\chaptermark{Forensique mémoire}

\epigraph{%
    « La mémoire vive est le seul témoin qui dit la vérité du moment présent\\
    et meurt dès qu'on éteint. »%
}{Principe du cours}

\niveauChapitre{Expert}{Chapitre~\ref{chap:systeme}.}

\begin{boxobjectifs}
\begin{itemize}
    \item Comprendre pourquoi et quand capturer la mémoire vive.
    \item Identifier ce que la mémoire révèle et que le disque ignore.
    \item Décrire l'acquisition et ses précautions.
    \item Repérer les traces d'exécution en mémoire.
    \item Mesurer le coût d'un redémarrage imprudent.
\end{itemize}
\end{boxobjectifs}

\bigskip

% ============================================================================
\section*{Question directrice : qu'a-t-on perdu en redémarrant FS-01~?}
% ============================================================================

En redémarrant le serveur, le technicien a effacé la mémoire vive. Ce
chapitre montre ce qu'elle contenait, pourquoi elle prime dans l'ordre de
volatilité, et comment on la capture quand elle est encore là.

\medskip
\textbf{Plan du chapitre.}
\begin{enumerate}
    \item Pourquoi la mémoire (section~\ref{sec:c13-pourquoi}).
    \item Ce qu'elle révèle (section~\ref{sec:c13-revele}).
    \item L'acquisition (section~\ref{sec:c13-acquisition}).
    \item Le coût du redémarrage (section~\ref{sec:c13-cout}).
\end{enumerate}

\separateur

% ============================================================================
\section{Pourquoi la Mémoire}
\label{sec:c13-pourquoi}
% ============================================================================

La mémoire vive contient l'état réel du système au moment présent~: processus
en cours, connexions réseau ouvertes, clés et secrets déchiffrés, code
malveillant qui ne touche jamais le disque. C'est pourquoi elle vient en tête
de l'ordre de volatilité (Chapitre~\ref{chap:fondements})~: elle disparaît à
l'extinction.

% ============================================================================
\section{Ce qu'elle Révèle}
\label{sec:c13-revele}
% ============================================================================

\begin{tableaucours}{L{4.0cm}L{8.2cm}}{%
    \tableauHeader{Élément en mémoire} &
    \tableauHeader{Intérêt pour l'enquête}
}
Processus en cours & Code malveillant actif, parent suspect \\
Connexions réseau & Vers où le système parlait (exfiltration, pilotage) \\
Secrets déchiffrés & Clés, mots de passe, jetons en clair \\
Code injecté & Programme présent en mémoire mais absent du disque \\
Commandes récentes & Ce que l'attaquant faisait à l'instant \\
\end{tableaucours}
\vspace{-0.8em}
\begin{center}{\small\itshape Tableau~13.1. Ce que révèle la mémoire vive}\end{center}

% ============================================================================
\section{L'Acquisition}
\label{sec:c13-acquisition}
% ============================================================================

\begin{boxprocedure}
\textbf{Acquérir la mémoire vive}
\begin{enumerate}
    \item Agir sur le système vivant, avant toute extinction.
    \item Utiliser un outil d'acquisition à empreinte minimale, documenté.
    \item Écrire la capture sur un support externe, jamais sur le disque
          analysé.
    \item Calculer le condensat de la capture.
    \item Consigner l'acte dans la chaîne de custody.
\end{enumerate}
\end{boxprocedure}

\begin{boxCRO}
\textbf{Capturer la mémoire, un arbitrage CRO.} L'acquisition touche
nécessairement le système vivant~: elle laisse une empreinte \cro{0}{-}{+}.
On accepte une légère altération (fiabilité) pour sauver une trace unique et
opposable (opposabilité), et on la documente pour la rendre défendable.
\end{boxCRO}

% ============================================================================
\section{Le Coût du Redémarrage}
\label{sec:c13-cout}
% ============================================================================

\begin{boxscenario}{Opération FARO. Ce que la mémoire aurait dit}
\textbf{Constatation.} La mémoire de \texttt{FS-01}, perdue au redémarrage,
aurait vraisemblablement contenu~: le processus de chiffrement actif, la
connexion d'exfiltration ouverte, et peut-être la clé de chiffrement en clair.

\medskip
\textbf{Analyse} (confiance moyenne). Sans la mémoire, l'enquête a dû
reconstituer ces éléments indirectement, à partir du disque et du réseau~:
plus long, moins complet. La clé en clair, si elle avait été capturée, aurait
pu éviter de payer la rançon. Le redémarrage a coûté cher.
\end{boxscenario}

\begin{boxmetier}{Analyste forensique et analyste SOC}
\textbf{Capturer la mémoire d'abord.} Face à un système compromis encore
allumé, le réflexe est de capturer la mémoire avant tout autre geste. Éteindre
« pour arrêter l'attaque » peut détruire la preuve qui aurait permis de la
comprendre.
\end{boxmetier}

\separateur

\begin{boxresume}
\textbf{Ce qu'il faut retenir de ce chapitre~:}
\begin{itemize}
    \item \textbf{La mémoire vive} contient l'état réel, et meurt à
          l'extinction.
    \item \textbf{Elle révèle} processus, connexions, secrets, code injecté.
    \item \textbf{On la capture} sur le système vivant, vers un support
          externe.
    \item \textbf{L'acquisition} est un arbitrage CRO documenté.
    \item \textbf{Un redémarrage} détruit cette preuve unique.
\end{itemize}
\end{boxresume}

\begin{boxinfo}
\textbf{Points de contrôle de l'élève.} Je sais~: expliquer la primauté de la
mémoire~; citer ce qu'elle révèle~; décrire son acquisition~; analyser le CRO
de l'acquisition~; mesurer le coût d'un redémarrage.
\end{boxinfo}

\bigskip

\begin{boxexo}[title={\ding{45}\quad Exercice~13.1. Pourquoi d'abord \niveaubase}]
Pourquoi la mémoire vient-elle en tête de l'ordre de volatilité~?
\end{boxexo}

\begin{boxexo}[title={\ding{45}\quad Exercice~13.2. Révélations \niveaubase}]
Citez trois éléments que la mémoire révèle et que le disque ignore.
\end{boxexo}

\begin{boxexo}[title={\ding{45}\quad Exercice~13.3. Support \niveaubase}]
Sur quel support écrit-on la capture mémoire~? Pourquoi pas le disque
analysé~?
\end{boxexo}

\begin{boxexo}[title={\ding{45}\quad Exercice~13.4. CRO \niveaumoyen}]
Expliquez l'arbitrage CRO de l'acquisition mémoire.
\end{boxexo}

\begin{boxexo}[title={\ding{45}\quad Exercice~13.5. FARO \niveaumoyen}]
Qu'aurait pu contenir la mémoire de \texttt{FS-01}~? Qu'a coûté sa perte~?
\end{boxexo}

\begin{boxexo}[title={\ding{45}\quad Exercice~13.6. Code injecté \niveaumoyen}]
Pourquoi un code injecté peut-il être invisible sur le disque~?
\end{boxexo}

\begin{boxexo}[title={\ding{45}\quad Exercice~13.7. Éteindre \niveaumoyen}]
« Il faut éteindre pour arrêter l'attaque. » Discutez.
\end{boxexo}

\begin{boxexo}[title={\ding{45}\quad Exercice~13.8. Empreinte \niveaumoyen}]
Pourquoi documente-t-on l'empreinte laissée par l'outil d'acquisition~?
\end{boxexo}

\begin{boxexo}[title={\ding{45}\quad Exercice~13.9. Procédure \niveauexpert}]
Rédigez la procédure de capture mémoire de première intervention pour LiZbiZ.
\end{boxexo}

\begin{boxexo}[title={\ding{45}\quad Exercice~13.10. Consigne \niveauexpert}]
Rédigez la consigne « système compromis allumé~: que faire en premier~? ».
\end{boxexo}

\bigskip

\begin{boxinfo}
\textbf{Pour aller plus loin.}
\begin{itemize}
    \item NIST SP~800-86~\cite{nist80086}~; RFC~3227~\cite{rfc3227}.
\end{itemize}
\end{boxinfo}

\begin{boxinfo}
\textbf{Transition.} Reste la troisième scène~: le réseau, qui garde la trace
des échanges. Le Chapitre~\ref{chap:reseau} traite du forensique réseau.
\end{boxinfo}
}{}
        \IfFileExists{partie4/P04_C14_forensique_reseau.tex}{% ============================================================================
% Fichier  : partie4/P04_C14_forensique_reseau.tex
% Titre    : Forensique Réseau
% Auteur   : MINKA MI NGUIDJOÏ Thierry Emmanuel
% ============================================================================
% Sources : RFC 3227 ; NIST SP 800-86 ; MITRE ATT&CK. Opération FARO.
% Le bureau d'analyse n'a pas d'outil de capture graphique : on travaille sur
% des journaux de flux exportés (CSV).
% ============================================================================

\chapter{Forensique Réseau}
\label{chap:reseau}
\chaptermark{Forensique réseau}

\epigraph{%
    « Les fichiers mentent parfois~; les flux, rarement. »%
}{Principe du cours}

\niveauChapitre{Intermédiaire à expert}{Chapitre~\ref{chap:memoire}~; notions
de réseaux (TCP/IP).}

\begin{boxobjectifs}
\begin{itemize}
    \item Identifier les sources de traces réseau (flux, journaux, DNS).
    \item Repérer une exfiltration et un canal de pilotage.
    \item Lire des enregistrements de flux exportés.
    \item Relier les traces réseau à la chronologie système.
    \item Rédiger un constat réseau.
\end{itemize}
\end{boxobjectifs}

\bigskip

% ============================================================================
\section*{Question directrice : où sont partis les 40 Go, et par quel canal l'attaquant pilotait-il~?}
% ============================================================================

L'exfiltration et le pilotage laissent des traces réseau~: connexions
sortantes inhabituelles, requêtes DNS répétées, volumes anormaux. Ce chapitre
apprend à les lire, à partir des journaux de flux exportés (le bureau
d'analyse n'a pas d'outil de capture graphique).

\medskip
\textbf{Plan du chapitre.}
\begin{enumerate}
    \item Les sources de traces réseau (section~\ref{sec:c14-sources}).
    \item Exfiltration et pilotage (section~\ref{sec:c14-exfil}).
    \item Lire des flux exportés (section~\ref{sec:c14-flux}).
    \item Recouper avec le système (section~\ref{sec:c14-recouper}).
\end{enumerate}

\separateur

% ============================================================================
\section{Les Sources de Traces Réseau}
\label{sec:c14-sources}
% ============================================================================

\begin{tableaucours}{L{3.6cm}L{8.6cm}}{%
    \tableauHeader{Source} &
    \tableauHeader{Ce qu'elle révèle}
}
Journaux de flux (pare-feu) & Qui parle à qui, quand, combien \\
Journaux du mandataire web & Destinations web, volumes, agents \\
Journaux DNS & Noms résolus, requêtes répétées ou suspectes \\
Journaux d'authentification distante & Accès VPN, bureau à distance \\
\end{tableaucours}
\vspace{-0.8em}
\begin{center}{\small\itshape Tableau~14.1. Sources de traces réseau}\end{center}

\begin{boxalert}
\textbf{Contrainte du cours.} Le bureau d'analyse de NDAP SUKULU ne dispose
pas d'outil de capture graphique.

\smallskip
\textbf{Règle~:} on travaille sur des \emph{journaux de flux exportés} (au
format tableur), ce qui suffit à établir les faits et reste parfaitement
opposable. L'essentiel n'est pas l'outil, mais la lecture méthodique.
\end{boxalert}

% ============================================================================
\section{Exfiltration et Pilotage}
\label{sec:c14-exfil}
% ============================================================================

Deux signatures réseau trahissent une intrusion active~: l'\termecle{exfiltration}
(un volume sortant anormal vers une destination inhabituelle) et le
\termecle{canal de pilotage} (de petites connexions régulières vers un serveur
distant, qui commandent le code malveillant). Les deux se repèrent dans les
journaux de flux.

% ============================================================================
\section{Lire des Flux Exportés}
\label{sec:c14-flux}
% ============================================================================

\begin{boxprocedure}
\textbf{Analyser des journaux de flux}
\begin{enumerate}
    \item Agréger par couple source-destination et par volume.
    \item Repérer les destinations externes inhabituelles.
    \item Repérer les connexions régulières de faible volume (pilotage).
    \item Relever les pics de volume sortant (exfiltration).
    \item Recaler les horodatages sur la chronologie système.
\end{enumerate}
\end{boxprocedure}

\begin{boxscenario}{Opération FARO. Les 40 Go de la nuit du 4 mars}
\textbf{Constatation.} Les journaux de flux montrent, dans la nuit du 4~mars,
un transfert sortant de 40~Go depuis \texttt{FS-01} vers une adresse externe,
précédé, depuis le 24~février, de petites connexions régulières (toutes les
10~minutes) vers la même adresse.

\medskip
\textbf{Analyse} (confiance élevée). Les petites connexions régulières
constituent le canal de pilotage~; le transfert de 40~Go, l'exfiltration. Les
dates recoupent la chronologie système (création du compte le 24, accès aux
partages le 4). Le faisceau est cohérent.
\end{boxscenario}

% ============================================================================
\section{Recouper avec le Système}
\label{sec:c14-recouper}
% ============================================================================

Une trace réseau isolée prouve peu~; recoupée avec le disque et la mémoire,
elle devient un faisceau solide. La force de l'enquête tient à la convergence
des trois scènes (système, mémoire, réseau) sur une même chronologie.

\begin{boxobjet}{Flux d'exfiltration du 4 mars}
\textbf{Nature~:} enregistrement de flux (journal du pare-feu, exporté).
\textbf{Source~:} copie du journal (condensat au dossier). \textbf{Pertinence~:}
transfert de 40~Go de \texttt{FS-01} vers une adresse externe.
\textbf{Fiabilité~:} intègre, horodatage recalé. \textbf{Opposabilité~:}
chaîne de custody continue, recoupé avec la chronologie système.
\end{boxobjet}

\begin{boxmetier}{Analyste forensique et analyste SOC}
\textbf{Faire converger les scènes.} L'analyste ne se contente pas d'une
source~: il recoupe réseau, système et mémoire. Une conclusion tenue par une
seule trace est fragile~; tenue par trois convergentes, elle est opposable.
\end{boxmetier}

\separateur

\begin{boxresume}
\textbf{Ce qu'il faut retenir de ce chapitre~:}
\begin{itemize}
    \item \textbf{Flux, mandataire, DNS}~: les sources réseau.
    \item \textbf{Exfiltration} (gros volume sortant) et \textbf{pilotage}
          (petites connexions régulières).
    \item \textbf{Des journaux de flux exportés} suffisent à établir les
          faits.
    \item \textbf{On recoupe} réseau, système et mémoire.
    \item \textbf{La convergence} fait la force de la preuve.
\end{itemize}
\end{boxresume}

\begin{boxinfo}
\textbf{Points de contrôle de l'élève.} Je sais~: citer les sources réseau~;
reconnaître exfiltration et pilotage~; lire des flux exportés~; recouper avec
la chronologie~; rédiger un constat réseau.
\end{boxinfo}

\bigskip

\begin{boxexo}[title={\ding{45}\quad Exercice~14.1. Sources \niveaubase}]
Quelle source révèle le volume d'un transfert sortant~? Les noms résolus~?
\end{boxexo}

\begin{boxexo}[title={\ding{45}\quad Exercice~14.2. Exfiltration \niveaubase}]
Quelle signature réseau trahit une exfiltration~?
\end{boxexo}

\begin{boxexo}[title={\ding{45}\quad Exercice~14.3. Pilotage \niveaubase}]
Comment se repère un canal de pilotage dans les flux~?
\end{boxexo}

\begin{boxexo}[title={\ding{45}\quad Exercice~14.4. Sans capture \niveaumoyen}]
Pourquoi des journaux de flux exportés suffisent-ils, sans capture graphique~?
\end{boxexo}

\begin{boxexo}[title={\ding{45}\quad Exercice~14.5. FARO \niveaumoyen}]
Dans FARO, distinguez le canal de pilotage de l'exfiltration.
\end{boxexo}

\begin{boxexo}[title={\ding{45}\quad Exercice~14.6. Recouper \niveaumoyen}]
Comment les flux du 4~mars recoupent-ils la chronologie système~?
\end{boxexo}

\begin{boxexo}[title={\ding{45}\quad Exercice~14.7. Faisceau \niveaumoyen}]
Pourquoi une trace réseau isolée prouve-t-elle peu~?
\end{boxexo}

\begin{boxexo}[title={\ding{45}\quad Exercice~14.8. Constat réseau \niveaumoyen}]
Rédigez le constat du flux d'exfiltration du 4~mars.
\end{boxexo}

\begin{boxexo}[title={\ding{45}\quad Exercice~14.9. Analyse de flux \niveauexpert}]
À partir d'un journal de flux exporté, proposez la démarche pour isoler
l'exfiltration.
\end{boxexo}

\begin{boxexo}[title={\ding{45}\quad Exercice~14.10. Convergence \niveauexpert}]
Construisez le faisceau de preuves de FARO en faisant converger les trois
scènes.
\end{boxexo}

\bigskip

\begin{boxinfo}
\textbf{Pour aller plus loin.}
\begin{itemize}
    \item RFC~3227~\cite{rfc3227}~; NIST SP~800-86~\cite{nist80086}~;
          MITRE ATT\&CK~\cite{mitreattack}.
\end{itemize}
\end{boxinfo}

\begin{boxinfo}
\textbf{Transition.} L'enquête établie, il faut agir~: la Partie~V traite de
la réponse à incident, à commencer par son cycle
(Chapitre~\ref{chap:cycle}).
\end{boxinfo}
}{}

    \part{La Réponse à Incident}
    \label{part:reponse}

        \IfFileExists{partie5/P05_C15_cycle_reponse.tex}{% ============================================================================
% Fichier  : partie5/P05_C15_cycle_reponse.tex
% Titre    : Le Cycle de Réponse à Incident
% Auteur   : MINKA MI NGUIDJOÏ Thierry Emmanuel
% ============================================================================
% Sources : ISO/IEC 27035 ; NIST SP 800-61r2. Opération MAYO : fictif.
% ============================================================================

\chapter{Le Cycle de Réponse à Incident}
\label{chap:cycle}
\chaptermark{Le cycle de réponse}

\epigraph{%
    « La réponse se prépare à froid~;\\
    en pleine crise, on applique, on n'improvise pas. »%
}{Principe du cours}

\niveauChapitre{Fondamental}{Chapitre~\ref{chap:normes}.}

\begin{boxobjectifs}
\begin{itemize}
    \item Décrire le cycle de réponse (ISO 27035, NIST 800-61r2).
    \item Classer un incident par gravité et fixer les délais.
    \item Comprendre le rôle de la préparation.
    \item Distinguer confinement, éradication, rétablissement.
    \item Articuler réponse et forensique sans les opposer.
\end{itemize}
\end{boxobjectifs}

\bigskip

% ============================================================================
\section*{Question directrice : le 6 mars à 7 h 15, que faut-il faire, et dans quel ordre~?}
% ============================================================================

Devant l'incident, l'urgence pousse à tout faire en même temps. Le cycle de
réponse impose un ordre~: qualifier, décider, contenir, éradiquer, rétablir,
apprendre. Ce chapitre le pose, avec les délais par gravité.

\medskip
\textbf{Plan du chapitre.}
\begin{enumerate}
    \item Le cycle (section~\ref{sec:c15-cycle}).
    \item Classer par gravité (section~\ref{sec:c15-gravite}).
    \item La préparation (section~\ref{sec:c15-preparation}).
    \item Réponse et forensique (section~\ref{sec:c15-forensique}).
\end{enumerate}

\separateur

% ============================================================================
\section{Le Cycle}
\label{sec:c15-cycle}
% ============================================================================

\begin{tableaucours}{L{3.6cm}L{8.6cm}}{%
    \tableauHeader{Étape} &
    \tableauHeader{Ce qu'on y fait}
}
Préparer & Politique, rôles, playbooks, outils, astreinte \\
Détecter et qualifier & Confirmer l'incident, évaluer la gravité \\
Contenir & Limiter la propagation sans détruire la preuve \\
Éradiquer & Supprimer la cause et la présence de l'attaquant \\
Rétablir & Restaurer le service à partir de sauvegardes vérifiées \\
Apprendre & Retour d'expérience, amélioration (Partie~VI) \\
\end{tableaucours}
\vspace{-0.8em}
\begin{center}{\small\itshape Tableau~15.1. Le cycle de réponse~\cite{iso27035,nist80061r2}}\end{center}

% ============================================================================
\section{Classer par Gravité}
\label{sec:c15-gravite}
% ============================================================================

\begin{tableaucours}{L{2.4cm}L{5.2cm}L{4.2cm}}{%
    \tableauHeader{Gravité} &
    \tableauHeader{Exemple} &
    \tableauHeader{Délai de prise en charge}
}
Critique & Rançongiciel, arrêt de production & Moins d'une heure, astreinte \\
Élevée & Compromission d'un compte à privilèges & Moins de 4 heures \\
Moyenne & Poste infecté isolé & Moins de 24 heures \\
Faible & Tentative bloquée & Moins de 72 heures \\
\end{tableaucours}
\vspace{-0.8em}
\begin{center}{\small\itshape Tableau~15.2. Grille de gravité et délais (LiZbiZ)}\end{center}

% ============================================================================
\section{La Préparation}
\label{sec:c15-preparation}
% ============================================================================

La préparation est la phase la plus déterminante~: la qualité de la réponse
se joue avant l'incident. Elle comprend les rôles, les coordonnées, les
playbooks (rançongiciel, hameçonnage, fuite de données), les sauvegardes
vérifiées et les droits d'accès d'urgence. Un incident révèle toujours l'état
de la préparation.

\begin{boxalert}
\textbf{Erreur fréquente.} Rétablir avant d'éradiquer.

\smallskip
\textbf{Règle~:} restaurer un service alors que l'attaquant est toujours
présent le ramène aussitôt. On éradique (et on vérifie l'éradication) avant
de rétablir.
\end{boxalert}

% ============================================================================
\section{Réponse et Forensique}
\label{sec:c15-forensique}
% ============================================================================

\begin{boxscenario}{Opération MAYO. L'ordre tenu}
\textbf{Constatation.} Le 6~mars, après qualification (critique), Mme~MOUAFO
décide~: isoler les systèmes atteints (confinement), capturer les preuves
avant toute remise en service, éradiquer le compte \texttt{svc\_maj2} et le
service \texttt{WinUpd32}, puis rétablir depuis des sauvegardes vérifiées hors
ligne.

\medskip
\textbf{Analyse} (confiance élevée). Le confinement a précédé le
rétablissement~; la capture des preuves (Partie~IV) a été insérée avant
l'éradication. Réponse et forensique ne s'opposent pas~: on contient, on fige,
puis on nettoie.
\end{boxscenario}

\begin{boxCRO}
\textbf{Contenir sans détruire la preuve.} Isoler une machine du réseau
préserve la preuve (on ne l'éteint pas) et limite la propagation~:
\cro{+}{+}{+}. Éteindre brutalement contiendrait aussi, mais détruirait la
mémoire~: \cro{+}{-}{-}. On préfère isoler.
\end{boxCRO}

\begin{boxmetier}{Responsable de la réponse à incident}
\textbf{Tenir l'ordre sous pression.} Le RRI résiste à la pression de
« tout rallumer tout de suite »~: il fait tenir l'ordre confiner, figer,
éradiquer, rétablir. C'est la discipline qui évite la rechute.
\end{boxmetier}

\separateur

\begin{boxresume}
\textbf{Ce qu'il faut retenir de ce chapitre~:}
\begin{itemize}
    \item \textbf{Préparer, qualifier, contenir, éradiquer, rétablir,
          apprendre.}
    \item \textbf{La gravité} fixe les délais.
    \item \textbf{La préparation} décide de la qualité de la réponse.
    \item \textbf{Éradiquer avant de rétablir}, sous peine de rechute.
    \item \textbf{Contenir en isolant}, pour préserver la preuve.
\end{itemize}
\end{boxresume}

\begin{boxinfo}
\textbf{Points de contrôle de l'élève.} Je sais~: décrire le cycle~; classer
par gravité~; expliquer le rôle de la préparation~; ordonner éradication et
rétablissement~; articuler réponse et forensique.
\end{boxinfo}

\bigskip

\begin{boxexo}[title={\ding{45}\quad Exercice~15.1. Le cycle \niveaubase}]
Citez les six étapes du cycle de réponse.
\end{boxexo}

\begin{boxexo}[title={\ding{45}\quad Exercice~15.2. Gravité \niveaubase}]
À quelle gravité classer un rançongiciel qui arrête la production~?
\end{boxexo}

\begin{boxexo}[title={\ding{45}\quad Exercice~15.3. Ordre \niveaubase}]
Pourquoi ne pas rétablir avant d'éradiquer~?
\end{boxexo}

\begin{boxexo}[title={\ding{45}\quad Exercice~15.4. Préparation \niveaumoyen}]
Pourquoi la préparation est-elle la phase la plus déterminante~?
\end{boxexo}

\begin{boxexo}[title={\ding{45}\quad Exercice~15.5. Isoler \niveaumoyen}]
Pourquoi isoler plutôt qu'éteindre~? (Analyse CRO)
\end{boxexo}

\begin{boxexo}[title={\ding{45}\quad Exercice~15.6. MAYO \niveaumoyen}]
Reconstituez l'ordre des décisions de Mme~MOUAFO le 6~mars.
\end{boxexo}

\begin{boxexo}[title={\ding{45}\quad Exercice~15.7. Figer \niveaumoyen}]
Où insère-t-on la capture des preuves dans le cycle~?
\end{boxexo}

\begin{boxexo}[title={\ding{45}\quad Exercice~15.8. Sauvegardes \niveaumoyen}]
Pourquoi rétablir à partir de sauvegardes vérifiées hors ligne~?
\end{boxexo}

\begin{boxexo}[title={\ding{45}\quad Exercice~15.9. Playbook \niveauexpert}]
Rédigez le playbook rançongiciel de LiZbiZ, étape par étape, avec les
décisions.
\end{boxexo}

\begin{boxexo}[title={\ding{45}\quad Exercice~15.10. Grille \niveauexpert}]
Construisez la grille de gravité de LiZbiZ avec délais, astreinte et
escalades.
\end{boxexo}

\bigskip

\begin{boxinfo}
\textbf{Pour aller plus loin.}
\begin{itemize}
    \item ISO/IEC 27035~\cite{iso27035}~; NIST SP~800-61r2~\cite{nist80061r2}.
\end{itemize}
\end{boxinfo}

\begin{boxinfo}
\textbf{Transition.} Répondre suppose une équipe organisée~: le
Chapitre~\ref{chap:csirt} traite du CSIRT et du CERT.
\end{boxinfo}
}{}
        \IfFileExists{partie5/P05_C16_csirt_cert.tex}{% ============================================================================
% Fichier  : partie5/P05_C16_csirt_cert.tex
% Titre    : Le CSIRT et le CERT
% Auteur   : MINKA MI NGUIDJOÏ Thierry Emmanuel
% ============================================================================
% Sources : ISO/IEC 27035 ; modèles de CSIRT/CERT ; ANTIC. Opération MAYO.
% ============================================================================

\chapter{Le CSIRT et le CERT}
\label{chap:csirt}
\chaptermark{Le CSIRT et le CERT}

\epigraph{%
    « En crise, ce qui n'est pas écrit n'a pas été décidé. »%
}{Principe du cours}

\niveauChapitre{Fondamental}{Chapitre~\ref{chap:cycle}.}

\begin{boxobjectifs}
\begin{itemize}
    \item Définir le CSIRT et le CERT, et leur différence.
    \item Décrire les rôles d'une cellule de réponse.
    \item Comprendre la main courante et la communication de crise.
    \item Situer la coopération avec l'ANTIC et les tiers.
    \item Dimensionner une cellule pour une PME.
\end{itemize}
\end{boxobjectifs}

\bigskip

% ============================================================================
\section*{Question directrice : qui décide, qui agit, qui parle pendant l'incident de LiZbiZ~?}
% ============================================================================

Sans organisation, un incident se gère dans la confusion~: tout le monde
agit, personne ne décide, et rien n'est tracé. Ce chapitre présente la
cellule de réponse, ses rôles, sa main courante et ses relations
extérieures.

\medskip
\textbf{Plan du chapitre.}
\begin{enumerate}
    \item CSIRT et CERT (section~\ref{sec:c16-definir}).
    \item Les rôles de la cellule (section~\ref{sec:c16-roles}).
    \item La main courante (section~\ref{sec:c16-maincourante}).
    \item Coopérer avec l'extérieur (section~\ref{sec:c16-cooperer}).
\end{enumerate}

\separateur

% ============================================================================
\section{CSIRT et CERT}
\label{sec:c16-definir}
% ============================================================================

\begin{boxdef}
\termecle{CSIRT} (\textit{Computer Security Incident Response Team})~: équipe
chargée de recevoir, analyser et traiter les incidents de sécurité d'un
périmètre donné. \termecle{CERT} est une appellation voisine, protégée~; tout
CERT est un CSIRT. L'\sigle{ANTIC} opère le CERT national camerounais.
\end{boxdef}

% ============================================================================
\section{Les Rôles de la Cellule}
\label{sec:c16-roles}
% ============================================================================

\begin{tableaucours}{L{3.6cm}L{8.6cm}}{%
    \tableauHeader{Rôle} &
    \tableauHeader{Responsabilité}
}
Responsable de la cellule & Décide, arbitre, rend compte à la direction \\
Analystes & Investiguent, contiennent, éradiquent \\
Référent technique métier & Connaît les systèmes atteints \\
Communication & Messages internes et externes, maîtrisés \\
Juridique & Obligations, signalements, ANTIC, assurance \\
Main courante & Consigne chaque décision et chaque action \\
\end{tableaucours}
\vspace{-0.8em}
\begin{center}{\small\itshape Tableau~16.1. Les rôles d'une cellule de réponse}\end{center}

% ============================================================================
\section{La Main Courante}
\label{sec:c16-maincourante}
% ============================================================================

\begin{boxprocedure}
\textbf{Tenir la main courante}
\begin{enumerate}
    \item Horodater chaque entrée.
    \item Noter la décision, qui l'a prise, sur quelle base.
    \item Noter l'action, qui l'a faite, son résultat.
    \item Ne rien reconstruire après coup~: on écrit sur le moment.
    \item Conserver la main courante au dossier (preuve et RETEX).
\end{enumerate}
\end{boxprocedure}

\begin{boxobjet}{Main courante de l'opération MAYO}
\textbf{Nature~:} journal horodaté des décisions et actions du 6~mars.
\textbf{Source~:} tenue en direct par le secrétaire de cellule.
\textbf{Pertinence~:} établit qui a décidé quoi, et quand.
\textbf{Fiabilité~:} entrées horodatées, non modifiées. \textbf{Opposabilité~:}
preuve de la diligence de LiZbiZ.
\end{boxobjet}

% ============================================================================
\section{Coopérer avec l'Extérieur}
\label{sec:c16-cooperer}
% ============================================================================

\begin{boxscenario}{Opération MAYO. La cellule de LiZbiZ}
\textbf{Constatation.} LiZbiZ n'avait pas de CSIRT. Le 6~mars, une cellule est
constituée en urgence~: M.~KOUNOU (technique), Mme~MOUAFO (responsable),
NYONG SÉCURITÉ (analystes), un juriste, un chargé de communication. L'ANTIC
est informée dans la journée.

\medskip
\textbf{Analyse} (confiance élevée). L'improvisation a coûté du temps.
Après l'incident, LiZbiZ formalise un CSIRT minimal et des playbooks~: la
prochaine cellule se réunira en minutes, pas en heures.
\end{boxscenario}

\begin{boxmetier}{Responsable de la réponse à incident}
\textbf{Désigner avant la crise.} Le responsable constitue la cellule et
nomme les rôles à froid, avec suppléants et coordonnées. En crise, on active
une liste, on ne cherche pas qui appeler.
\end{boxmetier}

\separateur

\begin{boxresume}
\textbf{Ce qu'il faut retenir de ce chapitre~:}
\begin{itemize}
    \item \textbf{CSIRT} (générique) et \textbf{CERT} (appellation
          protégée)~; l'ANTIC opère le CERT national.
    \item \textbf{Les rôles}~: responsable, analystes, métier, communication,
          juridique, main courante.
    \item \textbf{La main courante} trace décisions et actions, en direct.
    \item \textbf{On coopère} avec l'ANTIC et les tiers.
    \item \textbf{On désigne} la cellule avant la crise.
\end{itemize}
\end{boxresume}

\begin{boxinfo}
\textbf{Points de contrôle de l'élève.} Je sais~: distinguer CSIRT et CERT~;
citer les rôles~; tenir une main courante~; situer la coopération~;
dimensionner une cellule.
\end{boxinfo}

\bigskip

\begin{boxexo}[title={\ding{45}\quad Exercice~16.1. CSIRT ou CERT \niveaubase}]
Quelle est la différence entre CSIRT et CERT~?
\end{boxexo}

\begin{boxexo}[title={\ding{45}\quad Exercice~16.2. Rôles \niveaubase}]
Citez quatre rôles d'une cellule de réponse.
\end{boxexo}

\begin{boxexo}[title={\ding{45}\quad Exercice~16.3. Main courante \niveaubase}]
Que consigne-t-on dans une main courante~?
\end{boxexo}

\begin{boxexo}[title={\ding{45}\quad Exercice~16.4. Écrit \niveaumoyen}]
Pourquoi « ce qui n'est pas écrit n'a pas été décidé »~?
\end{boxexo}

\begin{boxexo}[title={\ding{45}\quad Exercice~16.5. ANTIC \niveaumoyen}]
Quand et pourquoi informer l'ANTIC~?
\end{boxexo}

\begin{boxexo}[title={\ding{45}\quad Exercice~16.6. MAYO \niveaumoyen}]
Qu'a coûté à LiZbiZ l'absence de CSIRT préétabli~?
\end{boxexo}

\begin{boxexo}[title={\ding{45}\quad Exercice~16.7. Communication \niveaumoyen}]
Pourquoi un rôle communication dédié dans la cellule~?
\end{boxexo}

\begin{boxexo}[title={\ding{45}\quad Exercice~16.8. Opposabilité \niveaumoyen}]
En quoi la main courante est-elle une preuve de diligence~?
\end{boxexo}

\begin{boxexo}[title={\ding{45}\quad Exercice~16.9. Charte de CSIRT \niveauexpert}]
Rédigez la charte du CSIRT minimal de LiZbiZ~: périmètre, rôles, suppléants,
coordonnées.
\end{boxexo}

\begin{boxexo}[title={\ding{45}\quad Exercice~16.10. Modèle de main courante \niveauexpert}]
Concevez le modèle de main courante de LiZbiZ.
\end{boxexo}

\bigskip

\begin{boxinfo}
\textbf{Pour aller plus loin.}
\begin{itemize}
    \item ISO/IEC 27035~\cite{iso27035}~; NIST SP~800-61r2~\cite{nist80061r2}.
\end{itemize}
\end{boxinfo}

\begin{boxinfo}
\textbf{Transition.} Pour anticiper, la cellule s'appuie sur le renseignement
sur la menace~: le Chapitre~\ref{chap:renseignement} le traite.
\end{boxinfo}
}{}
        \IfFileExists{partie5/P05_C17_threat_intelligence.tex}{% ============================================================================
% Fichier  : partie5/P05_C17_threat_intelligence.tex
% Titre    : Le Renseignement sur la Menace
% Auteur   : MINKA MI NGUIDJOÏ Thierry Emmanuel
% ============================================================================
% Sources : MITRE ATT&CK ; Kill Chain (Lockheed Martin, 2011) ; Diamond Model
% (2013) ; STIX/TAXII (OASIS). Opération MAYO.
% ============================================================================

\chapter{Le Renseignement sur la Menace}
\label{chap:renseignement}
\chaptermark{Le renseignement sur la menace}

\epigraph{%
    « Connaître l'adversaire, c'est transformer une alerte en histoire. »%
}{Principe du cours}

\niveauChapitre{Intermédiaire à expert}{Chapitre~\ref{chap:csirt}~;
Chapitre~\ref{chap:regles}.}

\begin{boxobjectifs}
\begin{itemize}
    \item Définir le renseignement sur la menace et ses usages.
    \item Utiliser la chaîne d'attaque (Kill Chain) et ATT\&CK.
    \item Comprendre le modèle du Diamant.
    \item Situer les formats d'échange (STIX/TAXII) et les indicateurs.
    \item Mapper un incident aux techniques de l'adversaire.
\end{itemize}
\end{boxobjectifs}

\bigskip

% ============================================================================
\section*{Question directrice : l'attaque de LiZbiZ ressemble-t-elle à d'autres, et qu'en déduire~?}
% ============================================================================

Une attaque n'est presque jamais inédite~: elle suit des étapes connues et
emploie des techniques répertoriées. Le renseignement sur la menace permet de
reconnaître le mode opératoire, d'anticiper la suite et d'écrire de meilleures
règles.

\medskip
\textbf{Plan du chapitre.}
\begin{enumerate}
    \item Qu'est-ce que le renseignement (section~\ref{sec:c17-definir}).
    \item Chaîne d'attaque et ATT\&CK (section~\ref{sec:c17-attack}).
    \item Le modèle du Diamant (section~\ref{sec:c17-diamant}).
    \item Échanger le renseignement (section~\ref{sec:c17-echanger}).
\end{enumerate}

\separateur

% ============================================================================
\section{Qu'est-ce que le Renseignement}
\label{sec:c17-definir}
% ============================================================================

\begin{boxdef}
\termecle{Renseignement sur la menace}~: information analysée sur les
adversaires, leurs motivations, leurs techniques et leurs indicateurs, qui
éclaire la détection, la réponse et la décision. Il transforme des données
brutes en connaissance actionnable.
\end{boxdef}

% ============================================================================
\section{Chaîne d'Attaque et ATT\&CK}
\label{sec:c17-attack}
% ============================================================================

La \termecle{chaîne d'attaque} (Lockheed Martin, 2011) décrit une intrusion
en étapes~: reconnaissance, armement, livraison, exploitation, installation,
pilotage, actions sur objectif. \termecle{MITRE ATT\&CK} affine en tactiques
et techniques observées. Les deux se complètent~: la chaîne raconte
l'histoire, ATT\&CK en nomme les gestes.

\begin{tableaucours}{L{4.2cm}L{4.4cm}L{3.6cm}}{%
    \tableauHeader{Étape (LOGONE)} &
    \tableauHeader{Tactique ATT\&CK} &
    \tableauHeader{Trace}
}
Accès initial (compte de service) & Accès aux identifiants & Connexion anormale \\
Installation (WinUpd32) & Persistance & Service installé \\
Pilotage (connexions régulières) & Commande et contrôle & Flux réguliers \\
Exfiltration (40 Go) & Exfiltration & Volume sortant \\
Chiffrement & Impact & Fichiers chiffrés \\
\end{tableaucours}
\vspace{-0.8em}
\begin{center}{\small\itshape Tableau~17.1. L'opération LOGONE mappée à ATT\&CK~\cite{mitreattack,killchain2011}}\end{center}

% ============================================================================
\section{Le Modèle du Diamant}
\label{sec:c17-diamant}
% ============================================================================

\begin{boxdef}
\termecle{Modèle du Diamant} (2013)~: cadre d'analyse d'une intrusion par
quatre sommets reliés~: l'\emph{adversaire}, ses \emph{capacités}, son
\emph{infrastructure}, et la \emph{victime}. Relier deux sommets connus permet
d'inférer les autres.
\end{boxdef}

Exemple~: connaissant l'infrastructure de pilotage (l'adresse externe) et la
victime (LiZbiZ), on peut rapprocher l'incident d'autres victimes partageant
la même infrastructure, et ainsi caractériser l'adversaire.

% ============================================================================
\section{Échanger le Renseignement}
\label{sec:c17-echanger}
% ============================================================================

Les \termecle{indicateurs de compromission} (adresses, condensats, noms de
domaine) s'échangent entre organisations dans des formats normalisés
(\termecle{STIX} pour décrire, \termecle{TAXII} pour transporter). L'ANTIC et
les CERT sectoriels jouent ce rôle de partage au niveau national.

\begin{boxalert}
\textbf{Prudence.} Un indicateur partagé peut vieillir ou être faux.

\smallskip
\textbf{Règle~:} on qualifie la source et la fraîcheur d'un indicateur avant
d'en faire une règle de blocage~: un indicateur périmé génère du bruit, un
indicateur erroné bloque du trafic légitime.
\end{boxalert}

\begin{boxscenario}{Opération MAYO. Le mode opératoire reconnu}
\textbf{Constatation.} Mappée à ATT\&CK et au Diamant, l'attaque de LiZbiZ
correspond à un mode opératoire de rançongiciel courant en Afrique centrale~:
accès par compte de service, persistance par service, exfiltration avant
chiffrement.

\medskip
\textbf{Analyse} (confiance moyenne). La reconnaissance du mode opératoire a
permis d'anticiper la recherche de persistances cachées et d'écrire des
règles ciblées. On ne nomme pas l'adversaire avec certitude~: NEK sur
l'attribution.
\end{boxscenario}

\begin{boxmetier}{Analyste SOC / Red Team}
\textbf{Du renseignement aux règles.} L'analyste transforme chaque
caractérisation d'adversaire en règles de détection et en points de chasse~:
le renseignement ne vaut que s'il améliore la défense.
\end{boxmetier}

\separateur

\begin{boxresume}
\textbf{Ce qu'il faut retenir de ce chapitre~:}
\begin{itemize}
    \item \textbf{Le renseignement} rend la menace actionnable.
    \item \textbf{La chaîne d'attaque} raconte, \textbf{ATT\&CK} nomme les
          gestes.
    \item \textbf{Le Diamant} relie adversaire, capacités, infrastructure,
          victime.
    \item \textbf{STIX/TAXII} échangent les indicateurs~; on qualifie leur
          fraîcheur.
    \item \textbf{L'attribution} reste souvent un NEK.
\end{itemize}
\end{boxresume}

\begin{boxinfo}
\textbf{Points de contrôle de l'élève.} Je sais~: définir le renseignement~;
mapper une attaque à la chaîne et à ATT\&CK~; appliquer le Diamant~; situer
STIX/TAXII~; transformer le renseignement en règles.
\end{boxinfo}

\bigskip

\begin{boxexo}[title={\ding{45}\quad Exercice~17.1. Chaîne \niveaubase}]
Citez trois étapes de la chaîne d'attaque.
\end{boxexo}

\begin{boxexo}[title={\ding{45}\quad Exercice~17.2. Diamant \niveaubase}]
Quels sont les quatre sommets du modèle du Diamant~?
\end{boxexo}

\begin{boxexo}[title={\ding{45}\quad Exercice~17.3. Indicateur \niveaubase}]
Donnez trois exemples d'indicateurs de compromission.
\end{boxexo}

\begin{boxexo}[title={\ding{45}\quad Exercice~17.4. Mapper \niveaumoyen}]
Mappez l'installation de \texttt{WinUpd32} à une tactique ATT\&CK.
\end{boxexo}

\begin{boxexo}[title={\ding{45}\quad Exercice~17.5. Inférer \niveaumoyen}]
Connaissant l'infrastructure et la victime, qu'infère-t-on avec le Diamant~?
\end{boxexo}

\begin{boxexo}[title={\ding{45}\quad Exercice~17.6. Indicateur périmé \niveaumoyen}]
Pourquoi qualifier la fraîcheur d'un indicateur avant d'en faire une règle~?
\end{boxexo}

\begin{boxexo}[title={\ding{45}\quad Exercice~17.7. Attribution \niveaumoyen}]
Pourquoi l'attribution reste-t-elle souvent un NEK~?
\end{boxexo}

\begin{boxexo}[title={\ding{45}\quad Exercice~17.8. STIX/TAXII \niveaumoyen}]
Quelle est la différence entre STIX et TAXII~?
\end{boxexo}

\begin{boxexo}[title={\ding{45}\quad Exercice~17.9. Mapping complet \niveauexpert}]
Dressez le mapping complet de l'opération LOGONE à ATT\&CK.
\end{boxexo}

\begin{boxexo}[title={\ding{45}\quad Exercice~17.10. Du renseignement aux règles \niveauexpert}]
À partir du mode opératoire reconnu, proposez trois règles et deux points de
chasse.
\end{boxexo}

\bigskip

\begin{boxinfo}
\textbf{Pour aller plus loin.}
\begin{itemize}
    \item MITRE ATT\&CK~\cite{mitreattack}~; chaîne d'attaque~\cite{killchain2011}~;
          modèle du Diamant~\cite{diamond2013}~; STIX/TAXII~\cite{stixtaxii}.
\end{itemize}
\end{boxinfo}

\begin{boxinfo}
\textbf{Transition.} Reste à agir concrètement sur l'incident~: le
Chapitre~\ref{chap:confinement} traite du confinement, de l'éradication et du
rétablissement.
\end{boxinfo}
}{}
        \IfFileExists{partie5/P05_C18_confinement_retablissement.tex}{% ============================================================================
% Fichier  : partie5/P05_C18_confinement_retablissement.tex
% Titre    : Confinement, Éradication, Rétablissement
% Auteur   : MINKA MI NGUIDJOÏ Thierry Emmanuel
% ============================================================================
% Sources : ISO/IEC 27035 ; NIST SP 800-61r2. Opération MAYO.
% ============================================================================

\chapter{Confinement, Éradication, Rétablissement}
\label{chap:confinement}
\chaptermark{Confiner, éradiquer, rétablir}

\epigraph{%
    « Contenir sans détruire la preuve, éradiquer sans oublier une porte,\\
    rétablir sans ramener l'attaquant. »%
}{Principe du cours}

\niveauChapitre{Intermédiaire à expert}{Chapitre~\ref{chap:cycle}.}

\begin{boxobjectifs}
\begin{itemize}
    \item Choisir une stratégie de confinement adaptée.
    \item Éradiquer complètement la présence de l'attaquant.
    \item Rétablir à partir de sauvegardes sûres.
    \item Vérifier qu'il n'y a pas de rechute.
    \item Décider, ou non, de payer une rançon.
\end{itemize}
\end{boxobjectifs}

\bigskip

% ============================================================================
\section*{Question directrice : comment remettre LiZbiZ en marche sans rouvrir la porte~?}
% ============================================================================

Après la qualification et la capture des preuves, il faut agir~: contenir,
éradiquer, rétablir. Chaque étape a ses pièges. Ce chapitre les parcourt, et
aborde la question difficile du paiement de la rançon.

\medskip
\textbf{Plan du chapitre.}
\begin{enumerate}
    \item Le confinement (section~\ref{sec:c18-confiner}).
    \item L'éradication (section~\ref{sec:c18-eradiquer}).
    \item Le rétablissement (section~\ref{sec:c18-retablir}).
    \item Faut-il payer~? (section~\ref{sec:c18-payer}).
\end{enumerate}

\separateur

% ============================================================================
\section{Le Confinement}
\label{sec:c18-confiner}
% ============================================================================

\begin{tableaucours}{L{3.6cm}L{4.6cm}L{4.0cm}}{%
    \tableauHeader{Stratégie} &
    \tableauHeader{Principe} &
    \tableauHeader{Quand}
}
Isolement réseau & Couper la machine du réseau sans l'éteindre & Préserve la preuve \\
Segmentation & Fermer un segment & Limiter la propagation \\
Blocage de comptes & Désactiver les comptes compromis & Couper l'accès \\
Arrêt (dernier recours) & Éteindre & Si la propagation l'exige \\
\end{tableaucours}
\vspace{-0.8em}
\begin{center}{\small\itshape Tableau~18.1. Stratégies de confinement}\end{center}

On privilégie l'isolement réseau, qui préserve la mémoire et la preuve
(Chapitre~\ref{chap:memoire}).

% ============================================================================
\section{L'Éradication}
\label{sec:c18-eradiquer}
% ============================================================================

\begin{boxprocedure}
\textbf{Éradiquer}
\begin{enumerate}
    \item Supprimer les accès de l'attaquant (comptes créés, clés déposées).
    \item Supprimer les persistances (services, tâches, clés de démarrage).
    \item Corriger la faille d'entrée (sinon l'attaquant revient).
    \item Changer les secrets exposés (mots de passe, clés).
    \item Vérifier l'éradication avant de rétablir.
\end{enumerate}
\end{boxprocedure}

\begin{boxalert}
\textbf{Erreur fréquente.} Supprimer le rançongiciel visible et croire la
machine propre.

\smallskip
\textbf{Règle~:} l'attaquant a souvent plusieurs portes (comptes, services,
accès distant). On éradique à partir de la chronologie forensique, qui liste
\emph{toutes} ses traces, pas seulement la charge finale.
\end{boxalert}

% ============================================================================
\section{Le Rétablissement}
\label{sec:c18-retablir}
% ============================================================================

On rétablit à partir de sauvegardes \termecle{vérifiées} et \termecle{hors
ligne} (une sauvegarde en ligne a pu être chiffrée elle aussi). On restaure
par étapes, en surveillant étroitement, prêt à reconfiner en cas de rechute.
La supervision renforcée (Partie~II) accompagne le retour à la normale.

% ============================================================================
\section{Faut-il Payer~?}
\label{sec:c18-payer}
% ============================================================================

\begin{boxscenario}{Opération MAYO. LiZbiZ ne paie pas}
\textbf{Constatation.} La rançon demandée s'élève à 25~millions de FCFA.
LiZbiZ dispose de sauvegardes hors ligne du serveur de suivi, vieilles de
cinq jours. Les données exfiltrées, elles, sont déjà perdues.

\medskip
\textbf{Analyse} (confiance élevée). Payer ne garantit ni la restitution ni
la non-divulgation, finance le crime, et peut être juridiquement risqué.
LiZbiZ choisit de ne pas payer~: elle restaure depuis les sauvegardes
(perte de cinq jours de données de suivi, reconstituées), notifie l'ANTIC et
les clients concernés par l'exfiltration. La décision est documentée.
\end{boxscenario}

\begin{boxCRO}
\textbf{La décision de payer, sous l'angle CRO.} Payer n'améliore pas la
\cro{-}{-}{-}~: les données exfiltrées restent exposées (C), rien ne garantit
l'intégrité des données « rendues » (F), et le paiement affaiblit la position
devant l'autorité (O). La préparation (sauvegardes) est la vraie réponse.
\end{boxCRO}

\begin{boxmetier}{Responsable de la réponse à incident et RSSI}
\textbf{La meilleure réponse au rançongiciel est antérieure à l'attaque.}
Des sauvegardes hors ligne testées rendent la question du paiement presque
théorique. Le RRI prépare cette réponse à froid.
\end{boxmetier}

\separateur

\begin{boxresume}
\textbf{Ce qu'il faut retenir de ce chapitre~:}
\begin{itemize}
    \item \textbf{Confiner} en isolant, pour préserver la preuve.
    \item \textbf{Éradiquer} toutes les portes, pas seulement la charge
          visible.
    \item \textbf{Corriger la faille d'entrée} avant de rétablir.
    \item \textbf{Rétablir} depuis des sauvegardes vérifiées hors ligne.
    \item \textbf{Payer} n'est pas une solution~; la préparation l'est.
\end{itemize}
\end{boxresume}

\begin{boxinfo}
\textbf{Points de contrôle de l'élève.} Je sais~: choisir une stratégie de
confinement~; éradiquer complètement~; rétablir sûrement~; vérifier l'absence
de rechute~; argumenter sur le paiement.
\end{boxinfo}

\bigskip

\begin{boxexo}[title={\ding{45}\quad Exercice~18.1. Confiner \niveaubase}]
Pourquoi préférer l'isolement réseau à l'extinction~?
\end{boxexo}

\begin{boxexo}[title={\ding{45}\quad Exercice~18.2. Portes \niveaubase}]
Citez trois « portes » qu'un attaquant peut laisser.
\end{boxexo}

\begin{boxexo}[title={\ding{45}\quad Exercice~18.3. Sauvegardes \niveaubase}]
Pourquoi rétablir depuis des sauvegardes hors ligne~?
\end{boxexo}

\begin{boxexo}[title={\ding{45}\quad Exercice~18.4. Éradication incomplète \niveaumoyen}]
Pourquoi supprimer seulement le rançongiciel visible ne suffit-il pas~?
\end{boxexo}

\begin{boxexo}[title={\ding{45}\quad Exercice~18.5. Faille d'entrée \niveaumoyen}]
Que se passe-t-il si l'on rétablit sans corriger la faille d'entrée~?
\end{boxexo}

\begin{boxexo}[title={\ding{45}\quad Exercice~18.6. Payer \niveaumoyen}]
Donnez trois raisons de ne pas payer une rançon.
\end{boxexo}

\begin{boxexo}[title={\ding{45}\quad Exercice~18.7. CRO \niveaumoyen}]
Analysez la décision de payer sous l'angle CRO.
\end{boxexo}

\begin{boxexo}[title={\ding{45}\quad Exercice~18.8. MAYO \niveaumoyen}]
Reconstituez la décision de LiZbiZ~: pourquoi ne pas payer, et comment
rétablir~?
\end{boxexo}

\begin{boxexo}[title={\ding{45}\quad Exercice~18.9. Plan de rétablissement \niveauexpert}]
Rédigez le plan de rétablissement de LiZbiZ après l'éradication.
\end{boxexo}

\begin{boxexo}[title={\ding{45}\quad Exercice~18.10. Politique rançongiciel \niveauexpert}]
Rédigez la politique de LiZbiZ sur le paiement des rançons, fondée sur la
préparation.
\end{boxexo}

\bigskip

\begin{boxinfo}
\textbf{Pour aller plus loin.}
\begin{itemize}
    \item ISO/IEC 27035~\cite{iso27035}~; NIST SP~800-61r2~\cite{nist80061r2}.
\end{itemize}
\end{boxinfo}

\begin{boxinfo}
\textbf{Transition.} L'incident traité, il reste à en tirer les leçons~: la
Partie~VI s'ouvre sur le retour d'expérience
(Chapitre~\ref{chap:retex}).
\end{boxinfo}
}{}

    \part{Apprendre et Mesurer}
    \label{part:mesurer}

        \IfFileExists{partie6/P06_C19_retex_amelioration.tex}{% ============================================================================
% Fichier  : partie6/P06_C19_retex_amelioration.tex
% Titre    : Retour d'Expérience et Amélioration Continue
% Auteur   : MINKA MI NGUIDJOÏ Thierry Emmanuel
% ============================================================================
% Sources : ISO/IEC 27035 (phase d'enseignements) ; culture juste. Op. MAYO.
% ============================================================================

\chapter{Retour d'Expérience et Amélioration Continue}
\label{chap:retex}
\chaptermark{Retour d'expérience}

\epigraph{%
    « Un incident dont on ne tire pas de leçon\\
    est un incident que l'on paiera deux fois. »%
}{Principe du cours}

\niveauChapitre{Fondamental}{Partie~V.}

\begin{boxobjectifs}
\begin{itemize}
    \item Conduire un retour d'expérience utile et sans procès.
    \item Distinguer causes immédiates et causes profondes.
    \item Transformer les leçons en actions suivies.
    \item Comprendre la culture juste.
    \item Boucler sur la supervision et la détection.
\end{itemize}
\end{boxobjectifs}

\bigskip

% ============================================================================
\section*{Question directrice : qu'a appris LiZbiZ, et comment éviter que cela se reproduise~?}
% ============================================================================

Une fois le service rétabli, la tentation est de tourner la page. C'est
l'erreur~: la phase d'enseignements de l'ISO 27035 est celle qui fait
progresser. Ce chapitre conduit le RETEX et le transforme en actions.

\medskip
\textbf{Plan du chapitre.}
\begin{enumerate}
    \item Conduire un RETEX (section~\ref{sec:c19-conduire}).
    \item Causes immédiates et profondes (section~\ref{sec:c19-causes}).
    \item Des leçons aux actions (section~\ref{sec:c19-actions}).
    \item La culture juste (section~\ref{sec:c19-culture}).
\end{enumerate}

\separateur

% ============================================================================
\section{Conduire un RETEX}
\label{sec:c19-conduire}
% ============================================================================

\begin{boxprocedure}
\textbf{Conduire un retour d'expérience}
\begin{enumerate}
    \item Réunir les acteurs dans les 72~heures, à partir de la main
          courante.
    \item Reconstituer la chronologie des faits (pas des responsabilités).
    \item Identifier ce qui a bien et mal fonctionné.
    \item Remonter aux causes profondes, au-delà du symptôme.
    \item Décider des actions, avec responsable et échéance.
\end{enumerate}
\end{boxprocedure}

% ============================================================================
\section{Causes Immédiates et Profondes}
\label{sec:c19-causes}
% ============================================================================

\begin{tableaucours}{L{4.4cm}L{7.8cm}}{%
    \tableauHeader{Niveau} &
    \tableauHeader{Exemple (GAROUA)}
}
Cause immédiate & Un compte de service à mot de passe faible compromis \\
Cause intermédiaire & Pas de supervision, journaux non lus \\
Cause profonde & Absence de fonction de détection et de gouvernance de la
sécurité \\
\end{tableaucours}
\vspace{-0.8em}
\begin{center}{\small\itshape Tableau~19.1. Remonter aux causes profondes}\end{center}

\begin{boxalert}
\textbf{Erreur fréquente.} S'arrêter à la cause immédiate (« le mot de passe
était faible ») et sanctionner.

\smallskip
\textbf{Règle~:} corriger la seule cause immédiate laisse la cause profonde
intacte~; l'incident reviendra sous une autre forme. On traite les trois
niveaux.
\end{boxalert}

% ============================================================================
\section{Des Leçons aux Actions}
\label{sec:c19-actions}
% ============================================================================

\begin{boxscenario}{Opération MAYO. Le plan d'action}
\textbf{Constatation.} Le RETEX de LiZbiZ débouche sur~: mettre en place une
supervision (SOC hybride), journaliser \texttt{FS-01}, durcir les comptes de
service, sauvegarder hors ligne, écrire des playbooks, former le personnel.
Chaque action a un responsable et une échéance.

\medskip
\textbf{Analyse} (confiance élevée). Les actions traitent les trois niveaux
de causes. Le RETEX n'est utile que si ses actions sont suivies~: LiZbiZ les
inscrit dans un plan revu chaque trimestre.
\end{boxscenario}

% ============================================================================
\section{La Culture Juste}
\label{sec:c19-culture}
% ============================================================================

\begin{boxdef}
\termecle{Culture juste}~: culture où l'on signale les erreurs et les
quasi-incidents sans crainte de sanction injuste. Elle distingue l'erreur
honnête (que l'on analyse) de la faute délibérée (que l'on sanctionne). Sans
elle, les incidents se cachent et ne s'apprennent pas.
\end{boxdef}

\begin{boxmetier}{RSSI}
\textbf{Protéger ceux qui signalent.} Le RSSI fait du RETEX un espace sans
procès~: le technicien qui a redémarré \texttt{FS-01} de bonne foi n'est pas
sanctionné, il devient l'exemple qui fonde la consigne de première
intervention. On apprend des erreurs, on ne les enterre pas.
\end{boxmetier}

\separateur

\begin{boxresume}
\textbf{Ce qu'il faut retenir de ce chapitre~:}
\begin{itemize}
    \item \textbf{Le RETEX} se conduit à partir de la main courante, sans
          procès.
    \item \textbf{Trois niveaux de causes}~: immédiate, intermédiaire,
          profonde.
    \item \textbf{On traite les trois}, sinon l'incident revient.
    \item \textbf{Des actions suivies}, avec responsable et échéance.
    \item \textbf{La culture juste} fait remonter ce qu'il faut apprendre.
\end{itemize}
\end{boxresume}

\begin{boxinfo}
\textbf{Points de contrôle de l'élève.} Je sais~: conduire un RETEX~;
distinguer les niveaux de causes~; transformer les leçons en actions~;
expliquer la culture juste~; boucler sur la supervision.
\end{boxinfo}

\bigskip

\begin{boxexo}[title={\ding{45}\quad Exercice~19.1. Délai \niveaubase}]
Dans quel délai conduire le RETEX~? À partir de quel document~?
\end{boxexo}

\begin{boxexo}[title={\ding{45}\quad Exercice~19.2. Causes \niveaubase}]
Distinguez cause immédiate et cause profonde dans GAROUA.
\end{boxexo}

\begin{boxexo}[title={\ding{45}\quad Exercice~19.3. Culture juste \niveaubase}]
Qu'est-ce que la culture juste~?
\end{boxexo}

\begin{boxexo}[title={\ding{45}\quad Exercice~19.4. Symptôme \niveaumoyen}]
Pourquoi corriger la seule cause immédiate ne suffit-il pas~?
\end{boxexo}

\begin{boxexo}[title={\ding{45}\quad Exercice~19.5. Actions \niveaumoyen}]
Reliez chaque action du plan de MAYO à un niveau de cause.
\end{boxexo}

\begin{boxexo}[title={\ding{45}\quad Exercice~19.6. Sans procès \niveaumoyen}]
Pourquoi le RETEX ne doit-il pas être un procès~?
\end{boxexo}

\begin{boxexo}[title={\ding{45}\quad Exercice~19.7. Le technicien \niveaumoyen}]
Comment traiter le cas du technicien qui a redémarré \texttt{FS-01}~?
\end{boxexo}

\begin{boxexo}[title={\ding{45}\quad Exercice~19.8. Suivi \niveaumoyen}]
Comment s'assurer que les actions du RETEX sont réellement suivies~?
\end{boxexo}

\begin{boxexo}[title={\ding{45}\quad Exercice~19.9. Compte rendu de RETEX \niveauexpert}]
Rédigez le compte rendu de RETEX de l'opération GAROUA.
\end{boxexo}

\begin{boxexo}[title={\ding{45}\quad Exercice~19.10. Plan d'amélioration \niveauexpert}]
Construisez le plan d'amélioration trimestriel de LiZbiZ issu du RETEX.
\end{boxexo}

\bigskip

\begin{boxinfo}
\textbf{Pour aller plus loin.}
\begin{itemize}
    \item ISO/IEC 27035~\cite{iso27035}.
\end{itemize}
\end{boxinfo}

\begin{boxinfo}
\textbf{Transition.} Reste à mesurer tout cela~: le
Chapitre~\ref{chap:mesurer} clôt le cours sur la mesure et la recherche.
\end{boxinfo}
}{}
        \IfFileExists{partie6/P06_C20_mesurer_recherche.tex}{% ============================================================================
% Fichier  : partie6/P06_C20_mesurer_recherche.tex
% Titre    : Mesurer la Supervision et la Réponse
% Auteur   : MINKA MI NGUIDJOÏ Thierry Emmanuel
% ============================================================================
% Sources : ISO/IEC 27004 ; cadre OUI/NON/NEK de l'auditologie (Minka). MBAM.
% ============================================================================

\chapter{Mesurer la Supervision et la Réponse}
\label{chap:mesurer}
\chaptermark{Mesurer et conclure}

\epigraph{%
    « Mesurer la sécurité, c'est aussi savoir dire ce que l'on ignore. »%
}{Principe du cours}

\niveauChapitre{Avancé}{L'ensemble du cours~; notions d'auditologie.}

\begin{boxobjectifs}
\begin{itemize}
    \item Choisir des indicateurs de supervision et de réponse honnêtes.
    \item Appliquer le verdict à trois valeurs OUI, NON, NEK.
    \item Éviter les indicateurs trompeurs.
    \item Inscrire la supervision dans un cycle d'amélioration.
    \item Ouvrir sur la recherche en détection et réponse.
\end{itemize}
\end{boxobjectifs}

\bigskip

% ============================================================================
\section*{Question directrice : LiZbiZ est-elle « prête » pour le prochain incident~?}
% ============================================================================

Six mois après, la question revient~: « Sommes-nous prêts~? » La réponse
honnête n'est ni oui ni non en bloc, mais un ensemble de verdicts précis, et,
pour ce qui n'a pu être éprouvé, un franc NEK. Ce chapitre clôt le cours sur
la mesure.

\medskip
\textbf{Plan du chapitre.}
\begin{enumerate}
    \item Des indicateurs honnêtes (section~\ref{sec:c20-indicateurs}).
    \item Le verdict à trois valeurs (section~\ref{sec:c20-verdict}).
    \item Les indicateurs trompeurs (section~\ref{sec:c20-trompeurs}).
    \item Le cycle et la recherche (section~\ref{sec:c20-cycle}).
\end{enumerate}

\separateur

% ============================================================================
\section{Des Indicateurs Honnêtes}
\label{sec:c20-indicateurs}
% ============================================================================

\begin{tableaucours}{L{4.8cm}L{7.4cm}}{%
    \tableauHeader{Indicateur} &
    \tableauHeader{Ce qu'il dit}
}
Délai de détection, de réaction & Rapidité à voir et à contenir \\
Couverture de journalisation & Part du parc réellement surveillé \\
Couverture de détection (ATT\&CK) & Part des techniques que l'on verrait \\
Taux de réussite des exercices & Préparation de la cellule \\
Taux de correction des actions RETEX & Capacité à apprendre \\
\end{tableaucours}
\vspace{-0.8em}
\begin{center}{\small\itshape Tableau~20.1. Indicateurs de supervision et de réponse~\cite{iso27004}}\end{center}

% ============================================================================
\section{Le Verdict à Trois Valeurs}
\label{sec:c20-verdict}
% ============================================================================

\begin{boxdef}
\termecle{Verdict à trois valeurs}~: à chaque question de préparation, la
réponse est \sigle{OUI} (démontré, avec preuve), \sigle{NON} (réfuté, avec
preuve) ou \sigle{NEK} (indéterminé~: les preuves manquent pour trancher). Le
NEK n'est pas un aveu de faiblesse~: c'est une exigence d'honnêteté qui
désigne où porter l'effort.
\end{boxdef}

\begin{tableaucours}{L{2.0cm}L{4.8cm}L{5.4cm}}{%
    \tableauHeader{Verdict} &
    \tableauHeader{Question} &
    \tableauHeader{Exemple (LiZbiZ)}
}
OUI & Détecte-t-on un compte de service anormal~? & Règle écrite et éprouvée \\
NON & Détecte-t-on l'exfiltration depuis tout serveur~? & FS-01 journalisé, pas les autres \\
NEK & Résisterait-on à un attaquant plus patient~? & Jamais éprouvé par un exercice \\
\end{tableaucours}
\vspace{-0.8em}
\begin{center}{\small\itshape Tableau~20.2. Le verdict OUI / NON / NEK appliqué~\cite{minka2026auditologie}}\end{center}

\begin{boxalert}
\textbf{Violation typique.} Conclure OUI (« nous sommes prêts ») faute d'avoir
éprouvé.

\smallskip
\textbf{Règle~:} ce qui n'a pas été testé par un exercice reste NEK, pas OUI.
L'absence d'incident n'est pas une preuve de préparation.
\end{boxalert}

% ============================================================================
\section{Les Indicateurs Trompeurs}
\label{sec:c20-trompeurs}
% ============================================================================

\begin{tableaucours}{L{4.4cm}L{7.8cm}}{%
    \tableauHeader{Indicateur trompeur} &
    \tableauHeader{Pourquoi il trompe}
}
Nombre d'alertes & Dépend des règles, pas du risque \\
« Aucun incident ce mois-ci » & Peut signifier « rien détecté » \\
Nombre d'outils déployés & L'outil n'est pas la détection \\
Temps passé au SOC & Mesure l'effort, pas le résultat \\
\end{tableaucours}
\vspace{-0.8em}
\begin{center}{\small\itshape Tableau~20.3. Indicateurs à manier avec prudence}\end{center}

% ============================================================================
\section{Le Cycle et la Recherche}
\label{sec:c20-cycle}
% ============================================================================

La supervision s'inscrit dans un cycle~: détecter, répondre, apprendre,
améliorer la détection, éprouver par des exercices, repartir des NEK. La
recherche prolonge ce cycle~: mesure de la couverture de détection,
automatisation de la réponse, renseignement adapté au contexte
d'Afrique centrale, et application du cadre de l'auditologie à la réduction
de l'incertitude en supervision.

\begin{boxscenario}{Opération MBAM. L'exercice qui lève des NEK}
\textbf{Constatation.} Pour éprouver sa préparation, LiZbiZ participe à un
exercice d'équipe (type équipe violette) avec NYONG SÉCURITÉ~: 30~techniques
jouées, 18~détectées. Les 12~non détectées étaient des NEK~: on ne savait pas
si on les verrait. Désormais, on sait que non, et on écrit les règles.

\medskip
\textbf{Analyse} (confiance élevée). L'exercice transforme des NEK en OUI ou
en NON étayés. C'est le seul moyen honnête de répondre « sommes-nous
prêts~? ».
\end{boxscenario}

\begin{boxmetier}{RSSI, analyste SOC, analyste forensique}
\textbf{Dire le vrai, NEK compris.} Le professionnel ne cède pas au OUI
rassurant~: il distingue ce qu'il a démontré, réfuté, et ce qu'il n'a pu
trancher. C'est à ce triple langage qu'on reconnaît un bilan honnête.
\end{boxmetier}

\separateur

\begin{boxresume}
\textbf{Ce qu'il faut retenir de ce chapitre~:}
\begin{itemize}
    \item \textbf{Des indicateurs} de délai, couverture, exercices,
          correction.
    \item \textbf{OUI, NON, NEK}~: le verdict honnête.
    \item \textbf{L'absence d'incident} n'est pas une preuve de préparation.
    \item \textbf{Les indicateurs d'activité} trompent.
    \item \textbf{Les exercices} transforment les NEK en verdicts étayés.
\end{itemize}
\end{boxresume}

\begin{boxinfo}
\textbf{Points de contrôle de l'élève.} Je sais~: choisir des indicateurs
honnêtes~; appliquer OUI/NON/NEK~; reconnaître un indicateur trompeur~;
décrire le cycle~; situer les pistes de recherche.
\end{boxinfo}

\bigskip

\begin{boxexo}[title={\ding{45}\quad Exercice~20.1. Indicateurs \niveaubase}]
Citez trois indicateurs honnêtes de supervision.
\end{boxexo}

\begin{boxexo}[title={\ding{45}\quad Exercice~20.2. Trois valeurs \niveaubase}]
Donnez le sens de OUI, NON et NEK.
\end{boxexo}

\begin{boxexo}[title={\ding{45}\quad Exercice~20.3. Trompeur \niveaubase}]
Pourquoi « aucun incident ce mois-ci » peut-il tromper~?
\end{boxexo}

\begin{boxexo}[title={\ding{45}\quad Exercice~20.4. Classer \niveaumoyen}]
Attribuez OUI, NON ou NEK~: détection d'un compte de service anormal~;
résistance à un attaquant patient~; détection de l'exfiltration partout.
\end{boxexo}

\begin{boxexo}[title={\ding{45}\quad Exercice~20.5. Pas de test \niveaumoyen}]
Pourquoi ce qui n'a pas été éprouvé reste-t-il NEK~?
\end{boxexo}

\begin{boxexo}[title={\ding{45}\quad Exercice~20.6. MBAM \niveaumoyen}]
Comment l'exercice MBAM transforme-t-il des NEK en verdicts étayés~?
\end{boxexo}

\begin{boxexo}[title={\ding{45}\quad Exercice~20.7. Note au DG \niveaumoyen}]
Répondez à « sommes-nous prêts~? » en OUI / NON / NEK pour LiZbiZ.
\end{boxexo}

\begin{boxexo}[title={\ding{45}\quad Exercice~20.8. Outils \niveaumoyen}]
Pourquoi le nombre d'outils déployés n'est-il pas un bon indicateur~?
\end{boxexo}

\begin{boxexo}[title={\ding{45}\quad Exercice~20.9. Tableau de bord \niveauexpert}]
Construisez le tableau de bord annuel de supervision et de réponse de LiZbiZ,
NEK compris.
\end{boxexo}

\begin{boxexo}[title={\ding{45}\quad Exercice~20.10. Question de recherche \niveauexpert}]
Formulez une question de recherche sur la mesure de la couverture de
détection en contexte d'Afrique centrale.
\end{boxexo}

\bigskip

\begin{boxinfo}
\textbf{Pour aller plus loin.}
\begin{itemize}
    \item ISO/IEC 27004~\cite{iso27004}~; cadre de
          l'auditologie~\cite{minka2026auditologie}.
\end{itemize}
\end{boxinfo}

\begin{boxinfo}
\textbf{Fin du cours.} De l'écran de rançon au verdict mesuré, SIS589 aura
montré qu'une attaque se voit venir, s'investigue proprement et se maîtrise~:
superviser, investiguer, répondre, apprendre, mesurer, sans jamais confondre
ce que l'on constate et ce que l'on en déduit.
\end{boxinfo}
}{}

% --------------------------------------------------------------------------
% ANNEXES
% --------------------------------------------------------------------------
\appendix

    %% Annexe A : Glossaire
\chapter{Glossaire}
\label{ann:glossaire}

Ce glossaire reprend les notions définies dans les chapitres. Le renvoi
indique le chapitre où la notion est présentée.

\begin{description}
    \item[Admissibilité] Aptitude d'une preuve à être reçue par un tiers, selon la manière dont elle a été recueillie et documentée (Ch.~\ref{chap:preuve}).
    \item[Analyse forensique] Recherche méthodique, documentée et opposable des traces d'un événement (Ch.~\ref{chap:fondements}).
    \item[ANTIC] Agence Nationale des Technologies de l'Information et de la Communication, autorité camerounaise de la sécurité électronique (Ch.~\ref{chap:droit}).
    \item[Canal de pilotage] Connexions régulières de faible volume par lesquelles un attaquant commande un code malveillant (Ch.~\ref{chap:reseau}).
    \item[CERT] Appellation protégée d'un CSIRT~; l'ANTIC opère le CERT national (Ch.~\ref{chap:csirt}).
    \item[Chaîne d'attaque] Description d'une intrusion en étapes, de la reconnaissance à l'action sur objectif (Ch.~\ref{chap:renseignement}).
    \item[Chaîne de custody] Continuité documentée de la détention d'une preuve (Ch.~\ref{chap:fondements}).
    \item[Chasse] Recherche proactive de ce qui n'a pas déclenché d'alerte (Ch.~\ref{chap:soc}).
    \item[Condensat] Empreinte numérique qui prouve l'intégrité d'une copie (Ch.~\ref{chap:fondements}).
    \item[Corréler] Rapprocher automatiquement des événements de sources différentes pour faire émerger une alerte (Ch.~\ref{chap:siem}).
    \item[CSIRT] Équipe de réponse aux incidents de sécurité d'un périmètre (Ch.~\ref{chap:csirt}).
    \item[Culture juste] Culture où l'on signale les erreurs sans crainte de sanction injuste (Ch.~\ref{chap:retex}).
    \item[Délai de présence] Durée entre la première compromission et sa détection (Ch.~\ref{chap:pourquoi}).
    \item[Exfiltration] Transfert sortant anormal de données vers une destination contrôlée par l'attaquant (Ch.~\ref{chap:reseau}).
    \item[Faux négatif] Comportement visé qu'une règle laisse passer (Ch.~\ref{chap:regles}).
    \item[Faux positif] Alerte sans incident réel (Ch.~\ref{chap:regles}).
    \item[Indicateur de compromission] Élément observable (adresse, condensat, domaine) trahissant une attaque (Ch.~\ref{chap:renseignement}).
    \item[Ingénierie de détection] Écriture et maintenance des règles de détection (Ch.~\ref{chap:regles}).
    \item[ISO/IEC 27035] Norme de gestion des incidents de sécurité, cycle en cinq phases (Ch.~\ref{chap:normes}).
    \item[ISO/IEC 27037] Lignes directrices de collecte et de préservation de la preuve numérique (Ch.~\ref{chap:normes}).
    \item[MITRE ATT\&CK] Base de connaissances des tactiques et techniques des attaquants (Ch.~\ref{chap:renseignement}).
    \item[Modèle du Diamant] Cadre d'analyse d'une intrusion~: adversaire, capacités, infrastructure, victime (Ch.~\ref{chap:renseignement}).
    \item[NIST SP 800-86] Guide d'intégration de la forensique à la réponse~: collecte, examen, analyse, rapport (Ch.~\ref{chap:normes}).
    \item[Ordre de volatilité] Ordre d'acquisition des preuves, du plus fugace au plus stable (Ch.~\ref{chap:fondements}).
    \item[RFC 3227] Lignes directrices de collecte et d'archivage des preuves (Ch.~\ref{chap:normes}).
    \item[SIEM] Plateforme de collecte, corrélation et alerte sur les journaux (Ch.~\ref{chap:siem}).
    \item[SOC] Centre d'opérations de sécurité~: détecter, analyser, traiter en continu (Ch.~\ref{chap:soc}).
    \item[STIX / TAXII] Formats d'échange du renseignement sur la menace (Ch.~\ref{chap:renseignement}).
    \item[Supervision de sécurité] Collecte et analyse continues des traces pour détecter les attaques (Ch.~\ref{chap:pourquoi}).
    \item[Trilemme CRO] Confidentialité, Fiabilité, Opposabilité~: les trois dimensions de toute trace et de tout acte d'enquête (Ch.~\ref{chap:prologue}).
    \item[Verdict à trois valeurs] OUI (démontré), NON (réfuté), NEK (indéterminé) (Ch.~\ref{chap:mesurer}).
    \item[Vrai positif] Alerte sur un comportement réellement visé (Ch.~\ref{chap:regles}).
\end{description}

    %% Annexe B : Référentiel des outils
\chapter{Référentiel des Outils}
\label{ann:outils}

Les outils cités sont des exemples. Le cours insiste sur la méthode~: l'outil
sert la lecture, il ne la remplace pas. Tout usage se fait dans un cadre
autorisé ou en laboratoire fermé.

\begin{tableaucours}{L{4.2cm}L{5.4cm}L{2.8cm}}{%
    \tableauHeader{Catégorie d'outil} &
    \tableauHeader{Usage} &
    \tableauHeader{Chapitre, lab}
}
Collecte de journaux & Centraliser, normaliser, horodater & Ch.~\ref{chap:journaux}, SAFI-LAB-02 \\
SIEM & Corréler, alerter, cas d'usage & Ch.~\ref{chap:siem}, SAFI-LAB-03 \\
Sondes IDS/IPS & Détecter réseau et hôte & Ch.~\ref{chap:idsips} \\
Acquisition forensique & Copie bit à bit, condensat & Ch.~\ref{chap:fondements}, SAFI-LAB-05 \\
Analyse système & Artefacts, chronologie & Ch.~\ref{chap:systeme}, SAFI-LAB-06 \\
Acquisition mémoire & Capture de la mémoire vive & Ch.~\ref{chap:memoire} \\
Analyse de flux & Journaux de flux exportés & Ch.~\ref{chap:reseau}, SAFI-LAB-07 \\
Référentiel ATT\&CK & Planifier, mapper, mesurer & Ch.~\ref{chap:renseignement}, SAFI-LAB-09 \\
Tableur & Chronologie, tableau de bord, couverture & Ch.~\ref{chap:gouvernance}, SAFI-LAB-10 \\
\end{tableaucours}

\begin{boxalert}
\textbf{Rappel.} Aucun de ces outils ne se dirige vers un système réel sans
autorisation. L'entraînement se fait en laboratoire fermé.
\end{boxalert}

    %% Annexe C : Modèles et procédures
\chapter{Modèles et Procédures}
\label{ann:templates}

\section*{Consigne de première intervention}
\begin{enumerate}
    \item Ne pas éteindre, ne pas redémarrer, ne pas « nettoyer ».
    \item Isoler la machine du réseau (sans l'éteindre).
    \item Prévenir le responsable de la réponse à incident.
    \item Noter l'heure et ce qui a été constaté.
    \item Attendre la capture des preuves avant toute action.
\end{enumerate}

\section*{Fiche de constat}
\begin{tableaucours}{L{3.4cm}L{8.8cm}}{%
    \tableauHeader{Rubrique} & \tableauHeader{Contenu}
}
Identifiant & T-xx, titre court \\
Constatation & Le fait vérifiable \\
Preuve & Artefact, condensat, source \\
Analyse & Interprétation, niveau de confiance \\
Gravité & Cotation justifiée \\
Recommandation & Correction concrète \\
\end{tableaucours}

\section*{Fiche de chaîne de custody}
\begin{tableaucours}{L{3.4cm}L{8.8cm}}{%
    \tableauHeader{Rubrique} & \tableauHeader{Contenu}
}
Pièce & Numéro, description, condensat \\
Acquisition & Qui, quand, comment, outil \\
Transferts & De qui à qui, quand, pourquoi \\
Conservation & Lieu, protection \\
\end{tableaucours}

\section*{Plan du rapport d'incident}
\begin{enumerate}
    \item Synthèse pour la direction.
    \item Chronologie des faits.
    \item Analyse (constatation / interprétation).
    \item Impact et obligations (ANTIC, données).
    \item Actions~: confinement, éradication, rétablissement.
    \item Retour d'expérience et recommandations.
\end{enumerate}

    %% Annexe D : Textes juridiques et normatifs
\chapter{Textes Juridiques et Normatifs}
\label{ann:textes}

Seuls les éléments vérifiés dans le texte officiel sont cités dans les
chapitres. Ce tableau indique où les trouver.

\section*{Droit camerounais et africain}
\begin{tableaucours}{L{5.0cm}L{4.6cm}L{2.6cm}}{%
    \tableauHeader{Texte} & \tableauHeader{Objet} & \tableauHeader{Chapitre}
}
Loi n° 2010/012 du 21 déc. 2010 & Cybersécurité et cybercriminalité & \ref{chap:droit} \\
Loi n° 2010/013 du 21 déc. 2010 & Communications électroniques & \ref{chap:droit} \\
Loi n° 2024/017 du 23 déc. 2024 & Protection des données personnelles & \ref{chap:droit} \\
Convention de Malabo (UA, 2014) & Cybersécurité et données, Afrique & \ref{chap:droit} \\
Convention de Budapest (2001) & Cybercriminalité, preuve, coopération & \ref{chap:droit} \\
\end{tableaucours}

\section*{Normes et guides internationaux}
\begin{tableaucours}{L{4.4cm}L{5.2cm}L{2.6cm}}{%
    \tableauHeader{Référence} & \tableauHeader{Objet} & \tableauHeader{Chapitre}
}
ISO/IEC 27035 & Gestion des incidents & \ref{chap:normes} \\
ISO/IEC 27037 & Preuve numérique & \ref{chap:normes} \\
ISO/IEC 27004 & Mesure de la sécurité & \ref{chap:gouvernance} \\
NIST SP 800-86 & Forensique et réponse & \ref{chap:normes} \\
NIST SP 800-61 Rev. 2 & Traitement des incidents & \ref{chap:cycle} \\
RFC 3227 & Collecte et archivage des preuves & \ref{chap:normes} \\
MITRE ATT\&CK & Tactiques et techniques & \ref{chap:renseignement} \\
CIS Controls v8 & Mesures de sécurité prioritaires & \ref{chap:soc} \\
\end{tableaucours}

    %% Annexe E : Bibliographie commentée
\chapter{Bibliographie Commentée}
\label{ann:bibliographie}

Pour chaque partie, quelques lectures. Les références complètes figurent dans
la bibliographie du livre.

\section*{Parties I et II · Superviser, collecter, détecter}
\begin{itemize}
    \item NIST SP 800-61r2~\cite{nist80061r2}~: le traitement des incidents.
    \item CIS Controls v8~\cite{cis2021controls}~: la place de la supervision.
    \item MITRE ATT\&CK~\cite{mitreattack}~: le vocabulaire des techniques.
\end{itemize}

\section*{Partie III · Cadre normatif et juridique}
\begin{itemize}
    \item ISO/IEC 27035~\cite{iso27035} et 27037~\cite{iso27037}~: incident et preuve.
    \item NIST SP 800-86~\cite{nist80086}, RFC 3227~\cite{rfc3227}~: forensique et volatilité.
    \item Lois 2010/012~\cite{loi2010012}, 2024/017~\cite{loi2024017}~; Malabo~\cite{malabo2014}~; Budapest~\cite{budapest2001}.
\end{itemize}

\section*{Partie IV · L'analyse forensique}
\begin{itemize}
    \item NIST SP 800-86~\cite{nist80086}~: la démarche forensique.
    \item ISO/IEC 27037~\cite{iso27037}~: la préservation.
\end{itemize}

\section*{Partie V · La réponse à incident}
\begin{itemize}
    \item ISO/IEC 27035~\cite{iso27035}~: le cycle de réponse.
    \item Chaîne d'attaque~\cite{killchain2011}, Diamant~\cite{diamond2013}, STIX/TAXII~\cite{stixtaxii}.
\end{itemize}

\section*{Partie VI · Apprendre et mesurer}
\begin{itemize}
    \item ISO/IEC 27004~\cite{iso27004}~: la mesure.
    \item T.~E.~Minka Mi Nguidjoï~\cite{minka2026auditologie}~: le verdict à trois valeurs.
\end{itemize}

    %% Annexe F : Contacts et ressources
\chapter{Contacts et Ressources}
\label{ann:contacts}

Les adresses sont celles des sites institutionnels à la date de rédaction~; à
vérifier.

\begin{tableaucours}{L{4.4cm}L{4.8cm}L{3.2cm}}{%
    \tableauHeader{Organisation} & \tableauHeader{Rôle} & \tableauHeader{Site}
}
ANTIC & Sécurité électronique, CERT national (Cameroun) & \texttt{www.antic.cm} \\
MINPOSTEL & Tutelle des communications électroniques & \texttt{www.minpostel.gov.cm} \\
Union africaine & Convention de Malabo & \texttt{au.int} \\
Conseil de l'Europe & Convention de Budapest & \texttt{coe.int} \\
MITRE & ATT\&CK & \texttt{attack.mitre.org} \\
NIST & Publications SP 800 & \texttt{www.nist.gov} \\
ISO & Normes ISO et ISO/IEC & \texttt{www.iso.org} \\
\end{tableaucours}

\begin{boxinfo}
\textbf{La plateforme du cours.} Les fiches d'auto-évaluation et les dix labs
de NDAP SUKULU pour ce cours sont accessibles par le code \texttt{SIS589}.
\end{boxinfo}


% --------------------------------------------------------------------------
% ÉLÉMENTS FINAUX
% --------------------------------------------------------------------------
\backmatter

    \nocite{*}
    \printbibliography[heading=bibintoc, title={Références bibliographiques}]

    \printindex

% ==============================================================================
% 4ÈME DE COUVERTURE
% ==============================================================================
\clearpage
\thispagestyle{empty}
\pagecolor{CouleurPrimaire}

\begin{tikzpicture}[remember picture, overlay]
    \fill[CouleurAccent]
        (current page.south west) rectangle
        ([yshift=1.8cm] current page.south east);
    \fill[CouleurAccent]
        ([yshift=-1.4cm] current page.north west) rectangle
        ([yshift=-1.7cm] current page.north east);
\end{tikzpicture}

\vspace*{2.5cm}

\begin{center}
    {\color{CouleurAccent}\rule{0.6\linewidth}{1pt}}\\[1.5ex]
    {\color{white}\Large\bfseries À propos de ce cours}\\[1.5ex]
    {\color{CouleurAccent}\rule{0.6\linewidth}{1pt}}
\end{center}

\vspace{1cm}

\begin{center}
\begin{minipage}{0.75\textwidth}
\color{white!85!gray}\setstretch{1.4}
Quand un système d'information s'arrête, la question décisive n'est pas
technique : elle porte sur qui avait décidé, et qui devait rendre compte.
Ce cours apprend à y répondre avant la panne.

\medskip

Il progresse en six parties, précédées d'un prologue : les fondements de la
gouvernance des systèmes d'information, le référentiel COBIT~2019, le cadre
normatif et juridique, ITIL~4 et les processus de gestion des services, les
défis de l'agilité, du nuage et de l'intelligence artificielle, et enfin la
recherche.

\medskip

Chaque notion est ancrée à la fois dans les standards internationaux (ISO,
ISACA, IIA, INTOSAI) et dans le droit camerounais et communautaire (CEMAC,
OHADA). Le cours s'adresse aux étudiants du Master Professionnel en
Management des Systèmes d'Information et aux professionnels qui exercent comme
DSI, responsables GRC, auditeurs des SI, responsables des services
informatiques ou RSSI.
\end{minipage}
\end{center}

\vfill

\begin{center}
    {\color{CouleurAccent}\rule{0.7\linewidth}{1.2pt}}\\[1.2ex]
    {\color{white}\large\bfseries MINKA MI NGUIDJOÏ Thierry Emmanuel}\\[0.4ex]
    {\color{white!80!gray}\small
        CISA~·~CISM~·~CGEIT~·~CRISC~·~CDPSE~·~ISO 27001~·~ISO 22301~·~ISO 9001}\\[0.3ex]
    {\color{white!70!gray}\footnotesize\itshape
        Chercheur au laboratoire d'Ingénierie Mathématiques et des Systèmes d'Information}\\[0.3ex]
    {\color{white!70!gray}\footnotesize\itshape
        Chercheur~|~Cryptologue~|~Auditeur Senior~|~Expert Judiciaire~|~GRC \& Sécurité de l'Information}\\[1.2ex]
    {\color{CouleurAccent}\rule{0.7\linewidth}{1.2pt}}
\end{center}

\vspace{1cm}

\afterpage{\nopagecolor}
\end{document}
